% Options for packages loaded elsewhere
% Options for packages loaded elsewhere
\PassOptionsToPackage{unicode}{hyperref}
\PassOptionsToPackage{hyphens}{url}
\PassOptionsToPackage{dvipsnames,svgnames,x11names}{xcolor}
%
\documentclass[
  french,
  letterpaper,
  open=any]{scrbook}
\usepackage{xcolor}
\usepackage[top=25mm,bottom=25mm,left=20mm,right=20mm]{geometry}
\usepackage{amsmath,amssymb}
\setcounter{secnumdepth}{-\maxdimen} % remove section numbering
\usepackage{iftex}
\ifPDFTeX
  \usepackage[T1]{fontenc}
  \usepackage[utf8]{inputenc}
  \usepackage{textcomp} % provide euro and other symbols
\else % if luatex or xetex
  \usepackage{unicode-math} % this also loads fontspec
  \defaultfontfeatures{Scale=MatchLowercase}
  \defaultfontfeatures[\rmfamily]{Ligatures=TeX,Scale=1}
\fi
\usepackage{lmodern}
\ifPDFTeX\else
  % xetex/luatex font selection
\fi
% Use upquote if available, for straight quotes in verbatim environments
\IfFileExists{upquote.sty}{\usepackage{upquote}}{}
\IfFileExists{microtype.sty}{% use microtype if available
  \usepackage[]{microtype}
  \UseMicrotypeSet[protrusion]{basicmath} % disable protrusion for tt fonts
}{}
\makeatletter
\@ifundefined{KOMAClassName}{% if non-KOMA class
  \IfFileExists{parskip.sty}{%
    \usepackage{parskip}
  }{% else
    \setlength{\parindent}{0pt}
    \setlength{\parskip}{6pt plus 2pt minus 1pt}}
}{% if KOMA class
  \KOMAoptions{parskip=half}}
\makeatother
% Make \paragraph and \subparagraph free-standing
\makeatletter
\ifx\paragraph\undefined\else
  \let\oldparagraph\paragraph
  \renewcommand{\paragraph}{
    \@ifstar
      \xxxParagraphStar
      \xxxParagraphNoStar
  }
  \newcommand{\xxxParagraphStar}[1]{\oldparagraph*{#1}\mbox{}}
  \newcommand{\xxxParagraphNoStar}[1]{\oldparagraph{#1}\mbox{}}
\fi
\ifx\subparagraph\undefined\else
  \let\oldsubparagraph\subparagraph
  \renewcommand{\subparagraph}{
    \@ifstar
      \xxxSubParagraphStar
      \xxxSubParagraphNoStar
  }
  \newcommand{\xxxSubParagraphStar}[1]{\oldsubparagraph*{#1}\mbox{}}
  \newcommand{\xxxSubParagraphNoStar}[1]{\oldsubparagraph{#1}\mbox{}}
\fi
\makeatother

\usepackage{color}
\usepackage{fancyvrb}
\newcommand{\VerbBar}{|}
\newcommand{\VERB}{\Verb[commandchars=\\\{\}]}
\DefineVerbatimEnvironment{Highlighting}{Verbatim}{commandchars=\\\{\}}
% Add ',fontsize=\small' for more characters per line
\usepackage{framed}
\definecolor{shadecolor}{RGB}{241,243,245}
\newenvironment{Shaded}{\begin{snugshade}}{\end{snugshade}}
\newcommand{\AlertTok}[1]{\textcolor[rgb]{0.68,0.00,0.00}{#1}}
\newcommand{\AnnotationTok}[1]{\textcolor[rgb]{0.37,0.37,0.37}{#1}}
\newcommand{\AttributeTok}[1]{\textcolor[rgb]{0.40,0.45,0.13}{#1}}
\newcommand{\BaseNTok}[1]{\textcolor[rgb]{0.68,0.00,0.00}{#1}}
\newcommand{\BuiltInTok}[1]{\textcolor[rgb]{0.00,0.23,0.31}{#1}}
\newcommand{\CharTok}[1]{\textcolor[rgb]{0.13,0.47,0.30}{#1}}
\newcommand{\CommentTok}[1]{\textcolor[rgb]{0.37,0.37,0.37}{#1}}
\newcommand{\CommentVarTok}[1]{\textcolor[rgb]{0.37,0.37,0.37}{\textit{#1}}}
\newcommand{\ConstantTok}[1]{\textcolor[rgb]{0.56,0.35,0.01}{#1}}
\newcommand{\ControlFlowTok}[1]{\textcolor[rgb]{0.00,0.23,0.31}{\textbf{#1}}}
\newcommand{\DataTypeTok}[1]{\textcolor[rgb]{0.68,0.00,0.00}{#1}}
\newcommand{\DecValTok}[1]{\textcolor[rgb]{0.68,0.00,0.00}{#1}}
\newcommand{\DocumentationTok}[1]{\textcolor[rgb]{0.37,0.37,0.37}{\textit{#1}}}
\newcommand{\ErrorTok}[1]{\textcolor[rgb]{0.68,0.00,0.00}{#1}}
\newcommand{\ExtensionTok}[1]{\textcolor[rgb]{0.00,0.23,0.31}{#1}}
\newcommand{\FloatTok}[1]{\textcolor[rgb]{0.68,0.00,0.00}{#1}}
\newcommand{\FunctionTok}[1]{\textcolor[rgb]{0.28,0.35,0.67}{#1}}
\newcommand{\ImportTok}[1]{\textcolor[rgb]{0.00,0.46,0.62}{#1}}
\newcommand{\InformationTok}[1]{\textcolor[rgb]{0.37,0.37,0.37}{#1}}
\newcommand{\KeywordTok}[1]{\textcolor[rgb]{0.00,0.23,0.31}{\textbf{#1}}}
\newcommand{\NormalTok}[1]{\textcolor[rgb]{0.00,0.23,0.31}{#1}}
\newcommand{\OperatorTok}[1]{\textcolor[rgb]{0.37,0.37,0.37}{#1}}
\newcommand{\OtherTok}[1]{\textcolor[rgb]{0.00,0.23,0.31}{#1}}
\newcommand{\PreprocessorTok}[1]{\textcolor[rgb]{0.68,0.00,0.00}{#1}}
\newcommand{\RegionMarkerTok}[1]{\textcolor[rgb]{0.00,0.23,0.31}{#1}}
\newcommand{\SpecialCharTok}[1]{\textcolor[rgb]{0.37,0.37,0.37}{#1}}
\newcommand{\SpecialStringTok}[1]{\textcolor[rgb]{0.13,0.47,0.30}{#1}}
\newcommand{\StringTok}[1]{\textcolor[rgb]{0.13,0.47,0.30}{#1}}
\newcommand{\VariableTok}[1]{\textcolor[rgb]{0.07,0.07,0.07}{#1}}
\newcommand{\VerbatimStringTok}[1]{\textcolor[rgb]{0.13,0.47,0.30}{#1}}
\newcommand{\WarningTok}[1]{\textcolor[rgb]{0.37,0.37,0.37}{\textit{#1}}}

\usepackage{longtable,booktabs,array}
\newcounter{none} % for unnumbered tables
\usepackage{calc} % for calculating minipage widths
% Correct order of tables after \paragraph or \subparagraph
\usepackage{etoolbox}
\makeatletter
\patchcmd\longtable{\par}{\if@noskipsec\mbox{}\fi\par}{}{}
\makeatother
% Allow footnotes in longtable head/foot
\IfFileExists{footnotehyper.sty}{\usepackage{footnotehyper}}{\usepackage{footnote}}
\makesavenoteenv{longtable}
\usepackage{graphicx}
\makeatletter
\newsavebox\pandoc@box
\newcommand*\pandocbounded[1]{% scales image to fit in text height/width
  \sbox\pandoc@box{#1}%
  \Gscale@div\@tempa{\textheight}{\dimexpr\ht\pandoc@box+\dp\pandoc@box\relax}%
  \Gscale@div\@tempb{\linewidth}{\wd\pandoc@box}%
  \ifdim\@tempb\p@<\@tempa\p@\let\@tempa\@tempb\fi% select the smaller of both
  \ifdim\@tempa\p@<\p@\scalebox{\@tempa}{\usebox\pandoc@box}%
  \else\usebox{\pandoc@box}%
  \fi%
}
% Set default figure placement to htbp
\def\fps@figure{htbp}
\makeatother


% definitions for citeproc citations
\NewDocumentCommand\citeproctext{}{}
\NewDocumentCommand\citeproc{mm}{%
  \begingroup\def\citeproctext{#2}\cite{#1}\endgroup}
\makeatletter
 % allow citations to break across lines
 \let\@cite@ofmt\@firstofone
 % avoid brackets around text for \cite:
 \def\@biblabel#1{}
 \def\@cite#1#2{{#1\if@tempswa , #2\fi}}
\makeatother
\newlength{\cslhangindent}
\setlength{\cslhangindent}{1.5em}
\newlength{\csllabelwidth}
\setlength{\csllabelwidth}{3em}
\newenvironment{CSLReferences}[2] % #1 hanging-indent, #2 entry-spacing
 {\begin{list}{}{%
  \setlength{\itemindent}{0pt}
  \setlength{\leftmargin}{0pt}
  \setlength{\parsep}{0pt}
  % turn on hanging indent if param 1 is 1
  \ifodd #1
   \setlength{\leftmargin}{\cslhangindent}
   \setlength{\itemindent}{-1\cslhangindent}
  \fi
  % set entry spacing
  \setlength{\itemsep}{#2\baselineskip}}}
 {\end{list}}
\usepackage{calc}
\newcommand{\CSLBlock}[1]{\hfill\break\parbox[t]{\linewidth}{\strut\ignorespaces#1\strut}}
\newcommand{\CSLLeftMargin}[1]{\parbox[t]{\csllabelwidth}{\strut#1\strut}}
\newcommand{\CSLRightInline}[1]{\parbox[t]{\linewidth - \csllabelwidth}{\strut#1\strut}}
\newcommand{\CSLIndent}[1]{\hspace{\cslhangindent}#1}

\ifLuaTeX
\usepackage[bidi=basic,shorthands=off]{babel}
\else
\usepackage[bidi=default,shorthands=off]{babel}
\fi
\ifLuaTeX
  \usepackage{selnolig} % disable illegal ligatures
\fi


\setlength{\emergencystretch}{3em} % prevent overfull lines

\providecommand{\tightlist}{%
  \setlength{\itemsep}{0pt}\setlength{\parskip}{0pt}}



 


\usepackage{graphicx}
\usepackage{ragged2e}
\makeatletter
\renewcommand{\maketitle}{}
\AtBeginDocument{\frontmatter\pagenumbering{gobble}\pagestyle{empty}}
\let\ps@plain\ps@empty
\makeatother
\makeatletter
\@ifpackageloaded{tcolorbox}{}{\usepackage[skins,breakable]{tcolorbox}}
\@ifpackageloaded{fontawesome5}{}{\usepackage{fontawesome5}}
\definecolor{quarto-callout-color}{HTML}{909090}
\definecolor{quarto-callout-note-color}{HTML}{0758E5}
\definecolor{quarto-callout-important-color}{HTML}{CC1914}
\definecolor{quarto-callout-warning-color}{HTML}{EB9113}
\definecolor{quarto-callout-tip-color}{HTML}{00A047}
\definecolor{quarto-callout-caution-color}{HTML}{FC5300}
\definecolor{quarto-callout-color-frame}{HTML}{acacac}
\definecolor{quarto-callout-note-color-frame}{HTML}{4582ec}
\definecolor{quarto-callout-important-color-frame}{HTML}{d9534f}
\definecolor{quarto-callout-warning-color-frame}{HTML}{f0ad4e}
\definecolor{quarto-callout-tip-color-frame}{HTML}{02b875}
\definecolor{quarto-callout-caution-color-frame}{HTML}{fd7e14}
\makeatother
\makeatletter
\@ifpackageloaded{bookmark}{}{\usepackage{bookmark}}
\makeatother
\makeatletter
\@ifpackageloaded{caption}{}{\usepackage{caption}}
\AtBeginDocument{%
\ifdefined\contentsname
  \renewcommand*\contentsname{Table des matières}
\else
  \newcommand\contentsname{Table des matières}
\fi
\ifdefined\listfigurename
  \renewcommand*\listfigurename{Liste des Figures}
\else
  \newcommand\listfigurename{Liste des Figures}
\fi
\ifdefined\listtablename
  \renewcommand*\listtablename{Liste des Tables}
\else
  \newcommand\listtablename{Liste des Tables}
\fi
\ifdefined\figurename
  \renewcommand*\figurename{Figure}
\else
  \newcommand\figurename{Figure}
\fi
\ifdefined\tablename
  \renewcommand*\tablename{Table}
\else
  \newcommand\tablename{Table}
\fi
}
\@ifpackageloaded{float}{}{\usepackage{float}}
\floatstyle{ruled}
\@ifundefined{c@chapter}{\newfloat{codelisting}{h}{lop}}{\newfloat{codelisting}{h}{lop}[chapter]}
\floatname{codelisting}{Listing}
\newcommand*\listoflistings{\listof{codelisting}{Liste des Listings}}
\usepackage{amsthm}
\theoremstyle{definition}
\newtheorem{exercise}{Exercice}[section]
\theoremstyle{definition}
\newtheorem{definition}{Définition}[section]
\theoremstyle{remark}
\AtBeginDocument{\renewcommand*{\proofname}{Preuve}}
\newtheorem*{remark}{Remarque}
\newtheorem*{solution}{Solution}
\newtheorem{refremark}{Remarque}[section]
\newtheorem{refsolution}{Solution}[section]
\makeatother
\makeatletter
\makeatother
\makeatletter
\@ifpackageloaded{caption}{}{\usepackage{caption}}
\@ifpackageloaded{subcaption}{}{\usepackage{subcaption}}
\makeatother
\usepackage{bookmark}
\IfFileExists{xurl.sty}{\usepackage{xurl}}{} % add URL line breaks if available
\urlstyle{same}
% fallback for those not using the hyperref driver hyperxmp:
\makeatletter
\@ifundefined{xmpquote}{\newcommand{\xmpquote}[1]{#1}}{}
\makeatother
\hypersetup{
  pdftitle={Analyse et visualisation de données},
  pdfauthor={IAOUSSE M'barek},
  pdflang={fr},
  colorlinks=true,
  linkcolor={Maroon},
  filecolor={Maroon},
  citecolor={Blue},
  urlcolor={Blue},
  pdfcreator={LaTeX via pandoc}}


\title{Analyse et visualisation de données}
\author{IAOUSSE M'barek}
\date{2026}
\begin{document}
\frontmatter
\maketitle


\mainmatter
\bookmarksetup{startatroot}

\chapter{}\label{section}

\makeatletter
\let\ps@plain\ps@empty
\makeatother
\pagenumbering{gobble}
\renewcommand{\thepage}{}
\pagestyle{empty}
\thispagestyle{empty}
\centering
\vspace*{2cm}
\includegraphics[width=6cm]{images/logo_estc.png}
\vfill
{\Huge\bfseries Analyse et visualisation de données\par}
\vspace{1cm}
{\large De la donnée brute au tableau de bord\par}
\vspace{0.5cm}
{\Large IAOUSSE M'barek\par}
\vspace{0.5cm}
{\large 2026\par}
\newpage
\thispagestyle{empty}
\centering
\vspace*{\fill}
\includegraphics[width=0.9\textwidth]{images/qrcode.png}
\vspace*{\fill}
\newpage

\small

\textbf{Ce manuel} a été élaboré pour servir de support de cours
principal pour l'élément \emph{Analyse et visualisation de données}
dispensé dans le cadre de la Licence Professionnelle DBT-DAIA
(\emph{Développement d'Applications et Intelligence Artificielle}) en
Temps Aménagé à l'École Supérieure de Technologie de Casablanca (ESTC).

\textbf{L'auteur}: M'barek IAOUSSE, Enseignant-Chercheur à l'EST
Casablanca, Université Hassan II de Casablanca.

\begin{tcolorbox}[enhanced jigsaw, arc=.35mm, bottomrule=.15mm, breakable, colback=white, colframe=quarto-callout-color-frame, left=2mm, leftrule=.75mm, opacityback=0, rightrule=.15mm, toprule=.15mm]

\begin{figure}[H]

\begin{minipage}[c]{0.23\linewidth}
\pandocbounded{\includegraphics[keepaspectratio]{images/CC/by.png}}\end{minipage}%
%
\begin{minipage}[c]{0.02\linewidth}
~\end{minipage}%
%
\begin{minipage}[c]{0.76\linewidth}
\justifying\small Sauf indications contraires, le contenu de ce manuel
électronique est disponible en vertu des termes de la
\href{https://creativecommons.org/licenses/by/4.0/deed.fr}{Licence
Creative Commons Attribution 4.0 International}.\end{minipage}%

\end{figure}%

\textbf{Autorisations accordées :}\\
\textbf{Partager} -- copier, distribuer et communiquer le matériel par
tous moyens et sous tous formats.\\
\textbf{Adapter} -- remixer, transformer et créer à partir du matériel
pour toute utilisation, y compris commerciale.

\textbf{Sous réserve de citer l'auteur.}

\end{tcolorbox}

\filbreak

© M'barek IAOUSSE (2026)\\
Pour citer ce manuel : M'barek IAOUSSE. (2026). \textbf{Analyse et
visualisation de données}. École Supérieure de Technologie de
Casablanca. Sous licence CC BY 4.0.

\bookmarksetup{startatroot}

\chapter{Préface}\label{pruxe9face}

\pagenumbering{roman}
\setcounter{page}{1}
\makeatletter
\def\ps@preface{%
  \let\@mkboth\@gobbletwo
  \def\@oddhead{}%
  \def\@evenhead{}%
  \def\@oddfoot{\hfil\thepage\hfil}%
  \let\@evenfoot\@oddfoot
}
\makeatother
\pagestyle{preface}
\thispagestyle{preface}
\phantomsection
%{\usekomafont{chapter}\strut Préface\par}

Le présent manuel est le support de cours officiel de l'élément
\textbf{Analyse et Visualisation des Données}, dispensé dans le cadre
d'une formation en Développement d'Applications et Intelligence
Artificielle.

La communauté étudiante visée par ce module provient de parcours
techniques variés (DUT, BTS, DTS ou diplômes équivalents), ayant déjà
acquis un bagage de base en programmation Python. La familiarité
préalable avec la Data Science, l'analyse de données ou le Machine
Learning n'est pas un prérequis : ce manuel se concentre volontairement
sur ce qui précède et accompagne la modélisation --- comprendre,
préparer, explorer et représenter une donnée --- sans aborder les
algorithmes d'apprentissage automatique, qui font l'objet d'un module
distinct.

Ce manuel vise à outiller les étudiants d'une \textbf{démarche
méthodique et reproductible} face à un jeu de données : savoir le
questionner avant de le modéliser, savoir le nettoyer avant de
l'analyser, et savoir le représenter avant de le communiquer. Il
introduit progressivement les statistiques descriptives nécessaires,
l'écosystème Python (\texttt{pandas}, \texttt{numpy},
\texttt{matplotlib}, \texttt{seaborn}), les principes d'un entrepôt de
données, puis la construction d'une solution de Business Intelligence
avec Power BI et Taipy. Un cas pratique final relie ces notions dans un
même projet.

La progression est volontairement graduelle : les chapitres 00 à 02
présentent l'analyse, la préparation et l'exploration ; le chapitre 03
organise les données pour l'analytique ; le chapitre 04 transforme ces
données en solution BI ; le chapitre 05 réunit l'ensemble dans une étude
de cas guidée.

\subsection*{Prérequis}

Le contenu de ce manuel est organisé pour des personnes ayant déjà une
base et une expérience avec :

\begin{itemize}
\tightlist
\item
  la programmation en langage Python (variables, boucles, fonctions,
  structures de données de base) ;
\item
  la manipulation élémentaire de fichiers et de tableaux.
\end{itemize}

Aucune connaissance préalable en analyse de données, en statistique ou
en data visualization n'est nécessaire : les indicateurs utiles sont
introduits et illustrés progressivement.

\subsection*{Livres et ressources complémentaires}

Ce manuel s'inspire des ouvrages de référence en analyse de données et
en visualisation, notamment \emph{Python for Data Analysis}
(\citeproc{ref-mckinney2022}{McKinney 2022}), consacré à la manipulation
de données avec \texttt{pandas}, ainsi que \emph{Storytelling with Data}
(\citeproc{ref-knaflic2015}{Knaflic 2015}), référence sur la
communication visuelle de résultats analytiques.

L'ensemble des exemples et exercices pratiques est développé en
utilisant \href{https://python.org}{Python} comme langage principal,
avec \texttt{pandas} et \texttt{numpy} pour la manipulation des données,
et \texttt{matplotlib}/\texttt{seaborn} pour la visualisation statique.

\begin{tcolorbox}[enhanced jigsaw, arc=.35mm, bottomrule=.15mm, bottomtitle=1mm, breakable, colback=white, colbacktitle=quarto-callout-note-color!10!white, colframe=quarto-callout-note-color-frame, coltitle=black, left=2mm, leftrule=.75mm, opacityback=0, opacitybacktitle=0.6, rightrule=.15mm, title=\textcolor{quarto-callout-note-color}{\faInfo}\hspace{0.5em}{Note}, titlerule=0mm, toprule=.15mm, toptitle=1mm]

Les bibliothèques scientifiques Python évoluant rapidement, veillez à
utiliser un environnement virtuel basé sur Python 3.10 ou supérieur pour
garantir la compatibilité des extraits de code présentés dans ce manuel.

\end{tcolorbox}

\clearpage
\tableofcontents
\clearpage

\clearpage
\setcounter{chapter}{0}
\pagenumbering{arabic}
\pagestyle{headings}
\makeatletter
\let\ps@plain\ps@headings
\makeatother

\bookmarksetup{startatroot}

\chapter{Introduction à l'analyse de
données}\label{introduction-uxe0-lanalyse-de-donnuxe9es}

\begin{tcolorbox}[enhanced jigsaw, arc=.35mm, bottomrule=.15mm, bottomtitle=1mm, breakable, colback=white, colbacktitle=quarto-callout-tip-color!10!white, colframe=quarto-callout-tip-color-frame, coltitle=black, left=2mm, leftrule=.75mm, opacityback=0, opacitybacktitle=0.6, rightrule=.15mm, title=\textcolor{quarto-callout-tip-color}{\faLightbulb}\hspace{0.5em}{Objectifs du chapitre}, titlerule=0mm, toprule=.15mm, toptitle=1mm]

À la fin de ce chapitre introductif, l'étudiant sera capable de :

\begin{itemize}
\tightlist
\item
  \textbf{Définir} l'analyse de données et la \textbf{situer} par
  rapport à la Data Science, la Business Intelligence (BI) et le Machine
  Learning ;
\item
  \textbf{Distinguer} les quatre types d'analyse : descriptive,
  diagnostique, prédictive et prescriptive ;
\item
  \textbf{Identifier} les grandes étapes du cycle de vie d'un projet
  d'analyse de données ;
\item
  \textbf{Identifier} les outils de base utilisés dans ce cours (Python,
  JupyterLab, \texttt{pandas}, \texttt{numpy}, \texttt{uv}) et les
  tendances récentes de l'écosystème.
\end{itemize}

\end{tcolorbox}

\section{1. Qu'est-ce que l'analyse de données
?}\label{quest-ce-que-lanalyse-de-donnuxe9es}

\subsection{1.1. Définition}\label{duxe9finition}

L'\textbf{analyse de données} (\emph{Data Analysis}) désigne l'ensemble
des démarches permettant d'inspecter, de nettoyer, de transformer et de
modéliser statistiquement des données, dans le but d'en extraire des
informations utiles, de soutenir une décision ou de valider une
hypothèse.

\begin{tcolorbox}[enhanced jigsaw, arc=.35mm, bottomrule=.15mm, bottomtitle=1mm, breakable, colback=white, colbacktitle=quarto-callout-tip-color!10!white, colframe=quarto-callout-tip-color-frame, coltitle=black, left=2mm, leftrule=.75mm, opacityback=0, opacitybacktitle=0.6, rightrule=.15mm, title=\textcolor{quarto-callout-tip-color}{\faLightbulb}\hspace{0.5em}{Analyse de données}, titlerule=0mm, toprule=.15mm, toptitle=1mm]

\begin{definition}[]\protect\hypertarget{def-analyse-donnees}{}\label{def-analyse-donnees}

L'analyse de données consiste à transformer des données brutes, souvent
volumineuses et désorganisées, en \textbf{informations exploitables}, à
travers une succession d'étapes de collecte, de nettoyage, d'exploration
et de représentation.

\end{definition}

\end{tcolorbox}

Contrairement à ce que l'on pense parfois, l'essentiel du travail d'un
analyste ne consiste pas à produire des graphiques, mais à
\textbf{comprendre les données} : d'où elles viennent, ce qu'elles
représentent réellement, quelles sont leurs limites et leurs biais. La
visualisation n'intervient qu'en bout de chaîne, pour communiquer un
résultat déjà établi.

\subsection{1.2. Data Analysis, Data Science, Business Intelligence,
Machine Learning : quelles frontières
?}\label{data-analysis-data-science-business-intelligence-machine-learning-quelles-frontiuxe8res}

Ces quatre termes sont souvent confondus. Ils partagent une matière
première commune --- la donnée --- mais diffèrent par leurs objectifs et
leurs méthodes.

\begin{longtable}[]{@{}
  >{\raggedright\arraybackslash}p{(\linewidth - 4\tabcolsep) * \real{0.3333}}
  >{\raggedright\arraybackslash}p{(\linewidth - 4\tabcolsep) * \real{0.3333}}
  >{\raggedright\arraybackslash}p{(\linewidth - 4\tabcolsep) * \real{0.3333}}@{}}
\caption{Quatre disciplines complémentaires autour de la
donnée}\label{tbl-disciplines-donnees}\tabularnewline
\toprule\noalign{}
\begin{minipage}[b]{\linewidth}\raggedright
Domaine
\end{minipage} & \begin{minipage}[b]{\linewidth}\raggedright
Question typique
\end{minipage} & \begin{minipage}[b]{\linewidth}\raggedright
Résultat produit
\end{minipage} \\
\midrule\noalign{}
\endfirsthead
\toprule\noalign{}
\begin{minipage}[b]{\linewidth}\raggedright
Domaine
\end{minipage} & \begin{minipage}[b]{\linewidth}\raggedright
Question typique
\end{minipage} & \begin{minipage}[b]{\linewidth}\raggedright
Résultat produit
\end{minipage} \\
\midrule\noalign{}
\endhead
\bottomrule\noalign{}
\endlastfoot
\textbf{Business Intelligence (BI)} & « Que s'est-il passé, et où en
sommes-nous ? » & Tableaux de bord, indicateurs (KPI) suivis dans le
temps \\
\textbf{Analyse de données (Data Analysis)} & « Pourquoi cela s'est-il
produit ? Que peut-on en conclure ? » & Rapports, visualisations,
recommandations argumentées \\
\textbf{Data Science} & « Que va-t-il probablement se passer ? » &
Modèles statistiques et modèles de Machine Learning \\
\textbf{Machine Learning} & « Comment automatiser une prédiction ou une
décision à grande échelle ? » & Un modèle entraîné, déployé en
production \\
\end{longtable}

\begin{tcolorbox}[enhanced jigsaw, arc=.35mm, bottomrule=.15mm, bottomtitle=1mm, breakable, colback=white, colbacktitle=quarto-callout-warning-color!10!white, colframe=quarto-callout-warning-color-frame, coltitle=black, left=2mm, leftrule=.75mm, opacityback=0, opacitybacktitle=0.6, rightrule=.15mm, title=\textcolor{quarto-callout-warning-color}{\faExclamationTriangle}\hspace{0.5em}{Périmètre de ce cours}, titlerule=0mm, toprule=.15mm, toptitle=1mm]

Ce module se concentre sur les deux premières colonnes de la
Table~\ref{tbl-disciplines-donnees} : la \textbf{préparation},
l'\textbf{exploration} et la \textbf{visualisation} des données. Les
algorithmes de Machine Learning (régression, classification,
clustering\ldots) ne sont volontairement \textbf{pas traités ici} : ils
font l'objet d'un module dédié. Une bonne partie du travail présenté
dans ce cours reste néanmoins un prérequis indispensable à tout projet
de Machine Learning.

\end{tcolorbox}

En pratique, ces frontières sont poreuses : un même professionnel peut,
selon les entreprises, être appelé \emph{Data Analyst} et exercer des
missions de \emph{Data Scientist}, ou inversement. Ce qui compte est de
comprendre \textbf{la nature de la question posée} plutôt que le titre
du poste.

\section{2. Les quatre types d'analyse de
données}\label{les-quatre-types-danalyse-de-donnuxe9es}

On distingue classiquement quatre niveaux d'analyse, de la simple
observation du passé jusqu'à la recommandation d'action, chacun
s'appuyant sur le précédent.

\begin{figure}

\centering{

\pandocbounded{\includegraphics[keepaspectratio]{index_files/mediabag/chapitres/../images/types-analyse-donnees.pdf}}

}

\caption{\label{fig-types-analyse}Les quatre types d'analyse de données}

\end{figure}%

\begin{longtable}[]{@{}
  >{\raggedright\arraybackslash}p{(\linewidth - 4\tabcolsep) * \real{0.3333}}
  >{\raggedright\arraybackslash}p{(\linewidth - 4\tabcolsep) * \real{0.3333}}
  >{\raggedright\arraybackslash}p{(\linewidth - 4\tabcolsep) * \real{0.3333}}@{}}
\caption{Les quatre types d'analyse de
données}\label{tbl-types-analyse}\tabularnewline
\toprule\noalign{}
\begin{minipage}[b]{\linewidth}\raggedright
Type d'analyse
\end{minipage} & \begin{minipage}[b]{\linewidth}\raggedright
Question posée
\end{minipage} & \begin{minipage}[b]{\linewidth}\raggedright
Exemple
\end{minipage} \\
\midrule\noalign{}
\endfirsthead
\toprule\noalign{}
\begin{minipage}[b]{\linewidth}\raggedright
Type d'analyse
\end{minipage} & \begin{minipage}[b]{\linewidth}\raggedright
Question posée
\end{minipage} & \begin{minipage}[b]{\linewidth}\raggedright
Exemple
\end{minipage} \\
\midrule\noalign{}
\endhead
\bottomrule\noalign{}
\endlastfoot
\textbf{Descriptive} & Que s'est-il passé ? & Le chiffre d'affaires du
mois dernier par région \\
\textbf{Diagnostique} & Pourquoi cela s'est-il passé ? & Pourquoi les
ventes ont-elles chuté en mars dans le nord ? \\
\textbf{Prédictive} & Que va-t-il probablement se passer ? & Quel sera
le chiffre d'affaires du mois prochain ? \\
\textbf{Prescriptive} & Que devrait-on faire ? & Quelle action marketing
maximiserait les ventes du trimestre ? \\
\end{longtable}

\begin{tcolorbox}[enhanced jigsaw, arc=.35mm, bottomrule=.15mm, bottomtitle=1mm, breakable, colback=white, colbacktitle=quarto-callout-tip-color!10!white, colframe=quarto-callout-tip-color-frame, coltitle=black, left=2mm, leftrule=.75mm, opacityback=0, opacitybacktitle=0.6, rightrule=.15mm, title=\textcolor{quarto-callout-tip-color}{\faLightbulb}\hspace{0.5em}{Exemple filé : une chaîne de supermarchés marocaine}, titlerule=0mm, toprule=.15mm, toptitle=1mm]

\begin{itemize}
\tightlist
\item
  \textbf{Descriptive} : « Les ventes de la région de Casablanca ont
  baissé de 12 \% en mars par rapport à février. »
\item
  \textbf{Diagnostique} : en croisant les ventes avec le calendrier et
  la météo, on découvre que la baisse coïncide avec des travaux routiers
  ayant réduit l'affluence dans deux magasins.
\item
  \textbf{Prédictive} \emph{(hors périmètre de ce cours, traité en
  Machine Learning)} : un modèle estime les ventes attendues pour le
  mois suivant compte tenu de la reprise du trafic.
\item
  \textbf{Prescriptive} : sur la base du diagnostic, l'équipe décide de
  renforcer temporairement les promotions dans les magasins concernés.
\end{itemize}

\end{tcolorbox}

Ce cours se concentre principalement sur les analyses
\textbf{descriptive} et \textbf{diagnostique} : elles reposent
entièrement sur les techniques de préparation, d'exploration et de
visualisation présentées dans les chapitres suivants, et ne nécessitent
pas d'algorithme d'apprentissage automatique.

\section{3. Les premiers indicateurs à
connaître}\label{les-premiers-indicateurs-uxe0-connauxeetre}

Avant de construire un graphique ou un tableau de bord, il faut savoir
décrire simplement les données. Une \textbf{observation} est une ligne
du tableau (par exemple une commande) et une \textbf{variable} est une
colonne (par exemple le montant de la commande). Une variable peut être
:

\begin{itemize}
\tightlist
\item
  \textbf{qualitative} : une catégorie ou une étiquette, comme une ville
  ou un moyen de paiement ;
\item
  \textbf{quantitative} : un nombre sur lequel on peut calculer, comme
  un prix, une quantité ou une durée ;
\item
  \textbf{temporelle} : une date ou une période, que l'on peut ordonner
  dans le temps.
\end{itemize}

\subsection{3.1. Compter et comparer}\label{compter-et-comparer}

Le \textbf{nombre d'observations} indique la taille du jeu de données.
Pour une catégorie, on utilise un effectif (combien de commandes à
Casablanca ?) et une fréquence (quelle part des commandes vient de
Casablanca ?).

Si 30 commandes sur 120 viennent de Casablanca :

\[
\text{fréquence} = \frac{30}{120} = 0{,}25 = 25\%.
\]

Pour comparer deux périodes, on peut calculer une variation :

\[
\text{variation en \%} =
\frac{\text{valeur finale} - \text{valeur initiale}}
{\text{valeur initiale}}\times 100.
\]

Une vente qui passe de 200 à 240 MAD a donc augmenté de
\(\frac{240-200}{200}\times 100=20\%\).

\subsection{3.2. Résumer une variable
numérique}\label{ruxe9sumer-une-variable-numuxe9rique}

Les indicateurs de tendance centrale et de dispersion seront utilisés
tout au long du manuel :

{\def\LTcaptype{none} % do not increment counter
\begin{longtable}[]{@{}
  >{\raggedright\arraybackslash}p{(\linewidth - 4\tabcolsep) * \real{0.3333}}
  >{\raggedright\arraybackslash}p{(\linewidth - 4\tabcolsep) * \real{0.3333}}
  >{\raggedright\arraybackslash}p{(\linewidth - 4\tabcolsep) * \real{0.3333}}@{}}
\toprule\noalign{}
\begin{minipage}[b]{\linewidth}\raggedright
Indicateur
\end{minipage} & \begin{minipage}[b]{\linewidth}\raggedright
Question à laquelle il répond
\end{minipage} & \begin{minipage}[b]{\linewidth}\raggedright
Exemple d'interprétation
\end{minipage} \\
\midrule\noalign{}
\endhead
\bottomrule\noalign{}
\endlastfoot
Moyenne & Quel serait le partage équitable du total ? & 10 commandes
totalisent 2 000 MAD : moyenne = 200 MAD. \\
Médiane & Quelle valeur partage les observations en deux groupes égaux ?
& La moitié des commandes vaut au plus la médiane. \\
Minimum / maximum & Quelles sont les valeurs extrêmes observées ? & La
commande la moins chère et la plus chère. \\
Quartiles \(Q_1\) et \(Q_3\) & Où se situent les 25 \% et 75 \% des
observations ? & 50 \% des valeurs sont entre \(Q_1\) et \(Q_3\). \\
Écart interquartile (IQR) & Quelle est la largeur de la moitié centrale
? & \(IQR=Q_3-Q_1\). \\
Écart-type & Les valeurs sont-elles plutôt proches de la moyenne ? &
Plus il est grand, plus les valeurs sont dispersées autour de la
moyenne. \\
\end{longtable}
}

La moyenne est calculée en additionnant les valeurs puis en divisant par
leur nombre. La médiane nécessite de trier les valeurs : avec un nombre
impair de valeurs, c'est la valeur du milieu ; avec un nombre pair,
c'est la moyenne des deux valeurs centrales. Une valeur très élevée peut
augmenter fortement la moyenne, mais elle influence beaucoup moins la
médiane. C'est pourquoi les deux indicateurs doivent souvent être
présentés ensemble.

\begin{Shaded}
\begin{Highlighting}[]
\NormalTok{resume }\OperatorTok{=}\NormalTok{ df[}\StringTok{"montant"}\NormalTok{].describe()}
\BuiltInTok{print}\NormalTok{(resume[[}\StringTok{"count"}\NormalTok{, }\StringTok{"mean"}\NormalTok{, }\StringTok{"50\%"}\NormalTok{, }\StringTok{"min"}\NormalTok{, }\StringTok{"max"}\NormalTok{]])}
\end{Highlighting}
\end{Shaded}

La commande \texttt{describe()} fournit notamment l'effectif
(\texttt{count}), la moyenne (\texttt{mean}), les quartiles
(\texttt{25\%}, \texttt{50\%}, \texttt{75\%}), le minimum et le maximum.
Elle ne remplace pas l'interprétation : un indicateur doit toujours être
accompagné de son unité, de sa période et de la population étudiée.

\begin{tcolorbox}[enhanced jigsaw, arc=.35mm, bottomrule=.15mm, bottomtitle=1mm, breakable, colback=white, colbacktitle=quarto-callout-warning-color!10!white, colframe=quarto-callout-warning-color-frame, coltitle=black, left=2mm, leftrule=.75mm, opacityback=0, opacitybacktitle=0.6, rightrule=.15mm, title=\textcolor{quarto-callout-warning-color}{\faExclamationTriangle}\hspace{0.5em}{Un indicateur n'est pas une conclusion}, titlerule=0mm, toprule=.15mm, toptitle=1mm]

Une moyenne de 250 MAD ne dit pas si les commandes sont toutes proches
de 250 MAD ou si quelques très grosses commandes compensent beaucoup de
petites commandes. Pour comprendre cela, on complète la moyenne par la
médiane, les quartiles ou un graphique.

\end{tcolorbox}

\section{4. Le cycle de vie d'un projet d'analyse de
données}\label{le-cycle-de-vie-dun-projet-danalyse-de-donnuxe9es}

Un projet d'analyse de données suit une succession d'étapes, rarement
linéaire, que l'on peut résumer ainsi :

\begin{figure}

\centering{

\pandocbounded{\includegraphics[keepaspectratio]{index_files/mediabag/chapitres/../images/cycle-vie-analyse-donnees.pdf}}

}

\caption{\label{fig-cycle-vie-analyse}Le cycle de vie d'un projet
d'analyse de données}

\end{figure}%

\begin{enumerate}
\def\labelenumi{\arabic{enumi}.}
\tightlist
\item
  \textbf{Cadrage} : formuler clairement la question métier à laquelle
  l'analyse doit répondre.
\item
  \textbf{Collecte} : rassembler les données nécessaires (bases de
  données, fichiers, API, web scraping\ldots).
\item
  \textbf{Nettoyage} : traiter les valeurs manquantes, les valeurs
  aberrantes, les doublons et les incohérences --- l'étape la plus
  chronophage en pratique.
\item
  \textbf{Transformation} : encoder, normaliser, recalculer certaines
  variables pour les rendre exploitables.
\item
  \textbf{Exploration (EDA)} : décrire statistiquement les données, en
  dégager des tendances, des corrélations, des anomalies.
\item
  \textbf{Visualisation et communication} : représenter graphiquement
  les résultats et les traduire en recommandations compréhensibles par
  un public non technique.
\end{enumerate}

\begin{tcolorbox}[enhanced jigsaw, arc=.35mm, bottomrule=.15mm, bottomtitle=1mm, breakable, colback=white, colbacktitle=quarto-callout-note-color!10!white, colframe=quarto-callout-note-color-frame, coltitle=black, left=2mm, leftrule=.75mm, opacityback=0, opacitybacktitle=0.6, rightrule=.15mm, title=\textcolor{quarto-callout-note-color}{\faInfo}\hspace{0.5em}{Repère}, titlerule=0mm, toprule=.15mm, toptitle=1mm]

Ce découpage est proche du cadre \textbf{OSEMN} (\emph{Obtain, Scrub,
Explore, Model, iNterpret}), popularisé par la communauté Data Science.
Nous l'adaptons ici en retirant l'étape de modélisation prédictive («
Model »), hors périmètre de ce cours, et en insistant davantage sur la
communication des résultats.

\end{tcolorbox}

Comme pour tout cycle de projet, ces étapes sont \textbf{itératives} :
une exploration révélant un problème de qualité de données ramène
généralement à l'étape de nettoyage, et une visualisation peu
convaincante amène souvent à revisiter l'exploration.

\section{5. Écosystème d'outils Python pour l'analyse de
données}\label{uxe9cosystuxe8me-doutils-python-pour-lanalyse-de-donnuxe9es}

\subsection{5.1. Bibliothèques cœur}\label{bibliothuxe8ques-cux153ur}

Comme pour le Machine Learning, Python s'impose comme langage principal
pour l'analyse de données, grâce à la richesse de son écosystème
scientifique et à la lisibilité de sa syntaxe.

\begin{longtable}[]{@{}
  >{\raggedright\arraybackslash}p{(\linewidth - 2\tabcolsep) * \real{0.5000}}
  >{\raggedright\arraybackslash}p{(\linewidth - 2\tabcolsep) * \real{0.5000}}@{}}
\caption{Bibliothèques Python cœur de ce
cours}\label{tbl-bibliotheques-eda}\tabularnewline
\toprule\noalign{}
\begin{minipage}[b]{\linewidth}\raggedright
Bibliothèque
\end{minipage} & \begin{minipage}[b]{\linewidth}\raggedright
Rôle
\end{minipage} \\
\midrule\noalign{}
\endfirsthead
\toprule\noalign{}
\begin{minipage}[b]{\linewidth}\raggedright
Bibliothèque
\end{minipage} & \begin{minipage}[b]{\linewidth}\raggedright
Rôle
\end{minipage} \\
\midrule\noalign{}
\endhead
\bottomrule\noalign{}
\endlastfoot
\texttt{pandas} & Manipulation de tableaux de données
(\texttt{DataFrame}) : lecture, nettoyage, agrégation \\
\texttt{numpy} & Calcul numérique et vectoriel sous-jacent à
\texttt{pandas} \\
\texttt{matplotlib} & Bibliothèque de visualisation de référence,
hautement personnalisable \\
\texttt{seaborn} & Visualisation statistique construite sur
\texttt{matplotlib}, orientée exploration rapide \\
\end{longtable}

\subsection{5.2. Environnements de
travail}\label{environnements-de-travail}

Comme dans le reste de la formation, nous privilégions les
\textbf{notebooks Jupyter} (JupyterLab, Google Colab) pour l'exploration
interactive et la visualisation immédiate cellule par cellule, et
\textbf{Visual Studio Code} pour consolider un script d'analyse
reproductible destiné à être réexécuté ou partagé.

\subsection{\texorpdfstring{5.3. Gestion de projet avec
\texttt{uv}}{5.3. Gestion de projet avec uv}}\label{gestion-de-projet-avec-uv}

Comme pour les autres modules de la formation, les projets de ce cours
sont organisés avec \href{https://docs.astral.sh/uv/}{\texttt{uv}}, qui
centralise la gestion de l'environnement virtuel, des dépendances
(\texttt{pyproject.toml}, \texttt{uv.lock}) et de l'exécution des
scripts et notebooks. Se référer au module précédent pour le détail de
son fonctionnement.

\subsection{5.4. Tendances récentes de
l'écosystème}\label{tendances-ruxe9centes-de-luxe9cosystuxe8me}

\begin{tcolorbox}[enhanced jigsaw, arc=.35mm, bottomrule=.15mm, bottomtitle=1mm, breakable, colback=white, colbacktitle=quarto-callout-important-color!10!white, colframe=quarto-callout-important-color-frame, coltitle=black, left=2mm, leftrule=.75mm, opacityback=0, opacitybacktitle=0.6, rightrule=.15mm, title=\textcolor{quarto-callout-important-color}{\faExclamation}\hspace{0.5em}{Zoom sur les tendances actuelles}, titlerule=0mm, toprule=.15mm, toptitle=1mm]

L'écosystème de l'analyse de données évolue rapidement. Sans remplacer
les outils historiques présentés ci-dessus, plusieurs bibliothèques
récentes gagnent en popularité et méritent d'être connues :

\begin{itemize}
\tightlist
\item
  \textbf{\href{https://pola.rs}{\texttt{Polars}}} : une alternative à
  \texttt{pandas}, écrite en Rust, offrant des performances nettement
  supérieures sur de gros volumes de données grâce à l'exécution
  parallèle et à l'évaluation paresseuse (\emph{lazy evaluation}).
\item
  \textbf{\href{https://duckdb.org}{\texttt{DuckDB}}} : un moteur de
  base de données analytique embarqué, permettant d'exécuter des
  requêtes SQL directement sur des fichiers CSV ou Parquet, sans
  serveur, avec des performances proches d'un entrepôt de données.
\item
  \textbf{Outils de profilage automatique} (\texttt{ydata-profiling},
  \texttt{Sweetviz}) : génèrent en une commande un rapport exploratoire
  complet d'un jeu de données (statistiques, distributions,
  corrélations, valeurs manquantes) --- voir le chapitre sur
  l'exploration des données.
\item
  \textbf{Assistants d'analyse par IA générative} : des outils
  permettant d'interroger un \texttt{DataFrame} en langage naturel
  commencent à s'intégrer aux environnements de data analysis, sans se
  substituer à une compréhension rigoureuse des données par l'analyste.
\end{itemize}

Ces outils ne sont pas obligatoires dans ce cours, mais seront
mentionnés ponctuellement à titre de complément, car ils représentent
une évolution significative des pratiques professionnelles actuelles.

\end{tcolorbox}

\section{6. Les métiers de la
donnée}\label{les-muxe9tiers-de-la-donnuxe9e}

\begin{longtable}[]{@{}
  >{\raggedright\arraybackslash}p{(\linewidth - 4\tabcolsep) * \real{0.3333}}
  >{\raggedright\arraybackslash}p{(\linewidth - 4\tabcolsep) * \real{0.3333}}
  >{\raggedright\arraybackslash}p{(\linewidth - 4\tabcolsep) * \real{0.3333}}@{}}
\caption{Les principaux métiers autour de la
donnée}\label{tbl-metiers-donnee}\tabularnewline
\toprule\noalign{}
\begin{minipage}[b]{\linewidth}\raggedright
Métier
\end{minipage} & \begin{minipage}[b]{\linewidth}\raggedright
Mission principale
\end{minipage} & \begin{minipage}[b]{\linewidth}\raggedright
Compétences clés
\end{minipage} \\
\midrule\noalign{}
\endfirsthead
\toprule\noalign{}
\begin{minipage}[b]{\linewidth}\raggedright
Métier
\end{minipage} & \begin{minipage}[b]{\linewidth}\raggedright
Mission principale
\end{minipage} & \begin{minipage}[b]{\linewidth}\raggedright
Compétences clés
\end{minipage} \\
\midrule\noalign{}
\endhead
\bottomrule\noalign{}
\endlastfoot
\textbf{Data Analyst} & Explorer, visualiser et interpréter les données
pour éclairer une décision & SQL, \texttt{pandas}, visualisation,
statistiques descriptives \\
\textbf{BI Analyst / BI Developer} & Construire et maintenir des
tableaux de bord et indicateurs récurrents & SQL, outils de BI (Power
BI, Tableau), modélisation de données \\
\textbf{Data Scientist} & Concevoir des modèles statistiques et
prédictifs & Statistique inférentielle, Machine Learning,
programmation \\
\textbf{Data Engineer} & Construire et fiabiliser les pipelines
d'acquisition et de stockage des données & Bases de données,
orchestration, ingénierie logicielle \\
\end{longtable}

Ce cours prépare avant tout au cœur du métier de \textbf{Data Analyst},
tout en posant des bases indispensables aux trois autres profils.

\begin{tcolorbox}[enhanced jigsaw, arc=.35mm, bottomrule=.15mm, bottomtitle=1mm, breakable, colback=white, colbacktitle=quarto-callout-note-color!10!white, colframe=quarto-callout-note-color-frame, coltitle=black, left=2mm, leftrule=.75mm, opacityback=0, opacitybacktitle=0.6, rightrule=.15mm, title=\textcolor{quarto-callout-note-color}{\faInfo}\hspace{0.5em}{Auto-évaluation}, titlerule=0mm, toprule=.15mm, toptitle=1mm]

Avant de passer au chapitre suivant, assurez-vous de pouvoir répondre «
oui » aux questions suivantes :

\begin{itemize}
\tightlist
\item[$\square$]
  Je peux expliquer la différence entre analyse de données, Data
  Science, BI et Machine Learning.
\item[$\square$]
  Je peux citer et illustrer les quatre types d'analyse de données.
\item[$\square$]
  Je peux énumérer, dans l'ordre, les étapes du cycle de vie d'un projet
  d'analyse de données.
\item[$\square$]
  Je connais au moins deux outils récents qui complètent \texttt{pandas}
  et \texttt{matplotlib}.
\end{itemize}

\end{tcolorbox}

\section{Exercices}\label{exercices}

\begin{exercise}[Vrai ou
Faux]\protect\hypertarget{exr-intro-ad-1}{}\label{exr-intro-ad-1}

~

\begin{enumerate}
\def\labelenumi{\alph{enumi}.}
\tightlist
\item
  La Business Intelligence répond principalement à la question «
  pourquoi ? ».
\item
  Une analyse diagnostique s'appuie généralement sur une analyse
  descriptive préalable.
\item
  Le Machine Learning est un prérequis indispensable pour produire une
  analyse descriptive.
\end{enumerate}

\end{exercise}

\begin{tcolorbox}[enhanced jigsaw, arc=.35mm, bottomrule=.15mm, bottomtitle=1mm, breakable, colback=white, colbacktitle=quarto-callout-tip-color!10!white, colframe=quarto-callout-tip-color-frame, coltitle=black, left=2mm, leftrule=.75mm, opacityback=0, opacitybacktitle=0.6, rightrule=.15mm, title=\textcolor{quarto-callout-tip-color}{\faLightbulb}\hspace{0.5em}{Solution de l'exercice Exercice~\ref{exr-intro-ad-1}}, titlerule=0mm, toprule=.15mm, toptitle=1mm]

\begin{enumerate}
\def\labelenumi{\alph{enumi}.}
\tightlist
\item
  \textbf{Faux.} La BI répond surtout à « que s'est-il passé, et où en
  sommes-nous ? » (analyse descriptive) ; c'est l'analyse diagnostique
  qui répond au « pourquoi ».
\item
  \textbf{Vrai.} On ne peut expliquer une variation (diagnostic) sans
  l'avoir d'abord constatée et quantifiée (description).
\item
  \textbf{Faux.} L'analyse descriptive repose sur des statistiques et
  des visualisations ; aucun modèle prédictif n'est nécessaire.
\end{enumerate}

\end{tcolorbox}

\begin{exercise}[Classer les
questions]\protect\hypertarget{exr-intro-ad-2}{}\label{exr-intro-ad-2}

Pour chacune des questions suivantes, indiquez le type d'analyse
correspondant (descriptive, diagnostique, prédictive ou prescriptive) :

\begin{enumerate}
\def\labelenumi{\alph{enumi}.}
\tightlist
\item
  Combien de visiteurs le site web a-t-il reçus la semaine dernière ?
\item
  Pourquoi le taux d'abandon de panier a-t-il augmenté en juillet ?
\item
  Quelle campagne publicitaire faut-il privilégier pour la rentrée ?
\end{enumerate}

\end{exercise}

\begin{tcolorbox}[enhanced jigsaw, arc=.35mm, bottomrule=.15mm, bottomtitle=1mm, breakable, colback=white, colbacktitle=quarto-callout-tip-color!10!white, colframe=quarto-callout-tip-color-frame, coltitle=black, left=2mm, leftrule=.75mm, opacityback=0, opacitybacktitle=0.6, rightrule=.15mm, title=\textcolor{quarto-callout-tip-color}{\faLightbulb}\hspace{0.5em}{Solution de l'exercice Exercice~\ref{exr-intro-ad-2}}, titlerule=0mm, toprule=.15mm, toptitle=1mm]

\begin{enumerate}
\def\labelenumi{\alph{enumi}.}
\tightlist
\item
  \textbf{Descriptive} --- simple constat chiffré.
\item
  \textbf{Diagnostique} --- recherche d'une cause à une variation
  observée.
\item
  \textbf{Prescriptive} --- recommandation d'action.
\end{enumerate}

\end{tcolorbox}

\begin{exercise}[Étude de cas
courte]\protect\hypertarget{exr-intro-ad-3}{}\label{exr-intro-ad-3}

Une association marocaine de protection de l'environnement souhaite
comprendre l'évolution du volume de déchets collectés dans trois villes
sur les deux dernières années, puis identifier les facteurs expliquant
les écarts entre villes.

\begin{enumerate}
\def\labelenumi{\alph{enumi}.}
\tightlist
\item
  Quelles étapes du cycle de vie présenté en
  Figure~\ref{fig-cycle-vie-analyse} cette association devra-t-elle
  suivre ?
\item
  À quel(s) type(s) d'analyse (Table~\ref{tbl-types-analyse}) ce projet
  correspond-il ?
\end{enumerate}

\end{exercise}

\begin{tcolorbox}[enhanced jigsaw, arc=.35mm, bottomrule=.15mm, bottomtitle=1mm, breakable, colback=white, colbacktitle=quarto-callout-tip-color!10!white, colframe=quarto-callout-tip-color-frame, coltitle=black, left=2mm, leftrule=.75mm, opacityback=0, opacitybacktitle=0.6, rightrule=.15mm, title=\textcolor{quarto-callout-tip-color}{\faLightbulb}\hspace{0.5em}{Solution de l'exercice Exercice~\ref{exr-intro-ad-3}}, titlerule=0mm, toprule=.15mm, toptitle=1mm]

\begin{enumerate}
\def\labelenumi{\alph{enumi}.}
\tightlist
\item
  Cadrage de la question (comparer trois villes) → collecte des données
  de collecte de déchets → nettoyage (unités, doublons, données
  manquantes par ville/mois) → exploration statistique → visualisation
  comparative dans le temps et par ville.
\item
  \textbf{Descriptive} (évolution des volumes) puis
  \textbf{diagnostique} (recherche des facteurs expliquant les écarts
  entre villes).
\end{enumerate}

\end{tcolorbox}

\begin{exercise}[Moyenne, médiane et
fréquence]\protect\hypertarget{exr-intro-ad-4}{}\label{exr-intro-ad-4}

Les montants de cinq commandes sont : 80, 100, 100, 120 et 600 MAD.

\begin{enumerate}
\def\labelenumi{\alph{enumi}.}
\tightlist
\item
  Calculez la moyenne et la médiane.
\item
  Quel indicateur décrit le mieux une commande « typique » ? Pourquoi ?
\item
  Quel est l'effectif et la fréquence des commandes d'au moins 100 MAD ?
\end{enumerate}

\end{exercise}

\begin{tcolorbox}[enhanced jigsaw, arc=.35mm, bottomrule=.15mm, bottomtitle=1mm, breakable, colback=white, colbacktitle=quarto-callout-tip-color!10!white, colframe=quarto-callout-tip-color-frame, coltitle=black, left=2mm, leftrule=.75mm, opacityback=0, opacitybacktitle=0.6, rightrule=.15mm, title=\textcolor{quarto-callout-tip-color}{\faLightbulb}\hspace{0.5em}{Solution de l'exercice Exercice~\ref{exr-intro-ad-4}}, titlerule=0mm, toprule=.15mm, toptitle=1mm]

\begin{enumerate}
\def\labelenumi{\alph{enumi}.}
\tightlist
\item
  La moyenne vaut \((80+100+100+120+600)/5=200\) MAD. La médiane est la
  troisième valeur triée : \textbf{100 MAD}.
\item
  La médiane est plus représentative : la commande de 600 MAD tire
  fortement la moyenne vers le haut.
\item
  Quatre commandes sur cinq sont d'au moins 100 MAD : effectif
  \textbf{4}, fréquence \textbf{80 \%}.
\end{enumerate}

\end{tcolorbox}

\begin{exercise}[Lire une
variation]\protect\hypertarget{exr-intro-ad-5}{}\label{exr-intro-ad-5}

Le nombre de visiteurs d'un site passe de 2 000 en avril à 2 400 en mai,
puis à 2 280 en juin.

\begin{enumerate}
\def\labelenumi{\alph{enumi}.}
\tightlist
\item
  Calculez la variation en pourcentage entre avril et mai.
\item
  Calculez la variation en pourcentage entre mai et juin.
\item
  La baisse de mai à juin annule-t-elle entièrement la hausse précédente
  ?
\end{enumerate}

\end{exercise}

\begin{tcolorbox}[enhanced jigsaw, arc=.35mm, bottomrule=.15mm, bottomtitle=1mm, breakable, colback=white, colbacktitle=quarto-callout-tip-color!10!white, colframe=quarto-callout-tip-color-frame, coltitle=black, left=2mm, leftrule=.75mm, opacityback=0, opacitybacktitle=0.6, rightrule=.15mm, title=\textcolor{quarto-callout-tip-color}{\faLightbulb}\hspace{0.5em}{Solution de l'exercice Exercice~\ref{exr-intro-ad-5}}, titlerule=0mm, toprule=.15mm, toptitle=1mm]

\begin{enumerate}
\def\labelenumi{\alph{enumi}.}
\tightlist
\item
  \((2400-2000)/2000 \times 100 = 20\%\).
\item
  \((2280-2400)/2400 \times 100 = -5\%\).
\item
  Non. Le niveau de juin reste supérieur à avril :
  \((2280-2000)/2000 \times 100 = 14\%\) environ.
\end{enumerate}

\end{tcolorbox}

\begin{exercise}[Identifier les indicateurs
utiles]\protect\hypertarget{exr-intro-ad-6}{}\label{exr-intro-ad-6}

Pour chacune des questions, choisissez l'indicateur le plus directement
utile : effectif, fréquence, moyenne, médiane, minimum/maximum ou IQR.

\begin{enumerate}
\def\labelenumi{\alph{enumi}.}
\tightlist
\item
  Quelle part des clients utilise le paiement par carte ?
\item
  Quel est le montant central des salaires ?
\item
  Quelle est la largeur de la moitié centrale des durées de livraison ?
\item
  Combien de commandes ont été enregistrées ?
\end{enumerate}

\end{exercise}

\begin{tcolorbox}[enhanced jigsaw, arc=.35mm, bottomrule=.15mm, bottomtitle=1mm, breakable, colback=white, colbacktitle=quarto-callout-tip-color!10!white, colframe=quarto-callout-tip-color-frame, coltitle=black, left=2mm, leftrule=.75mm, opacityback=0, opacitybacktitle=0.6, rightrule=.15mm, title=\textcolor{quarto-callout-tip-color}{\faLightbulb}\hspace{0.5em}{Solution de l'exercice Exercice~\ref{exr-intro-ad-6}}, titlerule=0mm, toprule=.15mm, toptitle=1mm]

\begin{enumerate}
\def\labelenumi{\alph{enumi}.}
\tightlist
\item
  \textbf{Fréquence} ; b. \textbf{médiane} ; c.~\textbf{IQR} ;
  d.~\textbf{effectif}.
\end{enumerate}

\end{tcolorbox}

\begin{exercise}[Reconnaître le type de
variable]\protect\hypertarget{exr-intro-ad-7}{}\label{exr-intro-ad-7}

Classez les variables suivantes en qualitative, quantitative ou
temporelle : \texttt{quartier}, \texttt{nombre\_de\_produits},
\texttt{date\_commande}, \texttt{satisfaction} codée « faible / moyenne
/ élevée ».

\end{exercise}

\begin{tcolorbox}[enhanced jigsaw, arc=.35mm, bottomrule=.15mm, bottomtitle=1mm, breakable, colback=white, colbacktitle=quarto-callout-tip-color!10!white, colframe=quarto-callout-tip-color-frame, coltitle=black, left=2mm, leftrule=.75mm, opacityback=0, opacitybacktitle=0.6, rightrule=.15mm, title=\textcolor{quarto-callout-tip-color}{\faLightbulb}\hspace{0.5em}{Solution de l'exercice Exercice~\ref{exr-intro-ad-7}}, titlerule=0mm, toprule=.15mm, toptitle=1mm]

\texttt{quartier} est \textbf{qualitative nominale} ;
\texttt{nombre\_de\_produits} est \textbf{quantitative discrète} ;
\texttt{date\_commande} est \textbf{temporelle} ; \texttt{satisfaction}
est \textbf{qualitative ordinale}, car ses catégories possèdent un
ordre.

\end{tcolorbox}

\begin{exercise}[Choisir une
analyse]\protect\hypertarget{exr-intro-ad-8}{}\label{exr-intro-ad-8}

Une responsable constate que les ventes sont plus faibles en août. Elle
veut d'abord mesurer la baisse, puis comprendre si elle concerne
seulement une région.

\begin{enumerate}
\def\labelenumi{\alph{enumi}.}
\tightlist
\item
  Quel type d'analyse répond à la première question ?
\item
  Quelle comparaison simple peut aider pour la seconde ?
\item
  Pourquoi faut-il éviter de conclure à une cause avant d'avoir comparé
  les données ?
\end{enumerate}

\end{exercise}

\begin{tcolorbox}[enhanced jigsaw, arc=.35mm, bottomrule=.15mm, bottomtitle=1mm, breakable, colback=white, colbacktitle=quarto-callout-tip-color!10!white, colframe=quarto-callout-tip-color-frame, coltitle=black, left=2mm, leftrule=.75mm, opacityback=0, opacitybacktitle=0.6, rightrule=.15mm, title=\textcolor{quarto-callout-tip-color}{\faLightbulb}\hspace{0.5em}{Solution de l'exercice Exercice~\ref{exr-intro-ad-8}}, titlerule=0mm, toprule=.15mm, toptitle=1mm]

\begin{enumerate}
\def\labelenumi{\alph{enumi}.}
\tightlist
\item
  Une analyse \textbf{descriptive} : comparer les ventes d'août à celles
  d'une période de référence.
\item
  Comparer les ventes par région et par mois, avec des effectifs, des
  totaux ou des moyennes.
\item
  Une baisse peut coïncider avec plusieurs facteurs ; une comparaison
  permet d'abord de vérifier si le phénomène est général ou localisé.
  Elle ne prouve pas à elle seule une causalité.
\end{enumerate}

\end{tcolorbox}

\bookmarksetup{startatroot}

\chapter{Préparation des données}\label{pruxe9paration-des-donnuxe9es}

\begin{tcolorbox}[enhanced jigsaw, arc=.35mm, bottomrule=.15mm, bottomtitle=1mm, breakable, colback=white, colbacktitle=quarto-callout-tip-color!10!white, colframe=quarto-callout-tip-color-frame, coltitle=black, left=2mm, leftrule=.75mm, opacityback=0, opacitybacktitle=0.6, rightrule=.15mm, title=\textcolor{quarto-callout-tip-color}{\faLightbulb}\hspace{0.5em}{Objectifs du chapitre}, titlerule=0mm, toprule=.15mm, toptitle=1mm]

À la fin de ce chapitre, l'étudiant sera capable de :

\begin{itemize}
\tightlist
\item
  \textbf{Identifier} les principales sources de collecte de données
  (bases de données, fichiers plats, API, web scraping) ;
\item
  \textbf{Diagnostiquer} les problèmes de qualité d'un jeu de données
  (valeurs manquantes, aberrantes, dupliquées) ;
\item
  \textbf{Appliquer} des stratégies adaptées de traitement des valeurs
  manquantes et aberrantes avec \texttt{pandas} ;
\item
  \textbf{Transformer} des variables (encodage, normalisation,
  discrétisation) en vue de leur exploration ou de leur visualisation.
\end{itemize}

\end{tcolorbox}

\begin{quote}
Si les données sont de mauvaise qualité, corrompues ou biaisées, aucune
analyse --- aussi soignée soit-elle --- ne pourra produire de
conclusions fiables (\emph{Garbage In, Garbage Out}).

--- Principe fondamental de l'analyse de données
\end{quote}

Avant de pouvoir explorer ou visualiser un jeu de données, il faut
d'abord le \textbf{rassembler} puis le \textbf{rendre exploitable}. Ce
chapitre couvre les étapes 2, 3 et 4 du cycle de vie présenté au
chapitre précédent (Figure~\ref{fig-cycle-vie-analyse}) : où trouver les
données, comment détecter et corriger leurs défauts, puis comment les
transformer pour l'exploration et la visualisation.

\begin{tcolorbox}[enhanced jigsaw, arc=.35mm, bottomrule=.15mm, bottomtitle=1mm, breakable, colback=white, colbacktitle=quarto-callout-note-color!10!white, colframe=quarto-callout-note-color-frame, coltitle=black, left=2mm, leftrule=.75mm, opacityback=0, opacitybacktitle=0.6, rightrule=.15mm, title=\textcolor{quarto-callout-note-color}{\faInfo}\hspace{0.5em}{Réflexe de débutant : connaître son tableau}, titlerule=0mm, toprule=.15mm, toptitle=1mm]

Avant toute modification, conservez une copie du fichier original et
rédigez un petit \textbf{dictionnaire de données}. Pour chaque colonne,
notez son nom, sa signification, son unité, son type attendu et les
valeurs possibles. Par exemple, \texttt{prix} peut être un nombre
positif en MAD, tandis que \texttt{type\_bien} peut prendre les valeurs
« appartement », « maison » ou « terrain ». Ce dictionnaire évite de
corriger une valeur qui était en réalité valide.

\end{tcolorbox}

\section{1. Collecte de données}\label{collecte-de-donnuxe9es}

\subsection{1.1. Les bases de données}\label{les-bases-de-donnuxe9es}

La source la plus courante en entreprise reste la base de données
interne :

\begin{itemize}
\tightlist
\item
  \textbf{relationnelle (SQL)} : données organisées en tables reliées
  par des clés (PostgreSQL, MySQL, SQL Server), interrogées avec le
  langage \textbf{SQL} ;
\item
  \textbf{non-relationnelle (NoSQL)} : données stockées sous une forme
  plus flexible --- documents JSON (MongoDB), paires clé-valeur (Redis)
  --- adaptées à des données peu structurées ou à très grande échelle.
\end{itemize}

\texttt{pandas} permet de charger directement le résultat d'une requête
SQL dans un \texttt{DataFrame} :

\begin{Shaded}
\begin{Highlighting}[]
\ImportTok{import}\NormalTok{ pandas }\ImportTok{as}\NormalTok{ pd}
\ImportTok{import}\NormalTok{ sqlite3}

\NormalTok{connexion }\OperatorTok{=}\NormalTok{ sqlite3.}\ExtensionTok{connect}\NormalTok{(}\StringTok{"ventes.db"}\NormalTok{)}
\NormalTok{df }\OperatorTok{=}\NormalTok{ pd.read\_sql(}\StringTok{"SELECT * FROM transactions WHERE annee = 2025"}\NormalTok{, connexion)}
\end{Highlighting}
\end{Shaded}

\subsection{1.2. Les fichiers plats et formats de
données}\label{les-fichiers-plats-et-formats-de-donnuxe9es}

En dehors des bases de données, les données sont fréquemment échangées
sous forme de fichiers :

\begin{longtable}[]{@{}
  >{\raggedright\arraybackslash}p{(\linewidth - 4\tabcolsep) * \real{0.3333}}
  >{\raggedright\arraybackslash}p{(\linewidth - 4\tabcolsep) * \real{0.3333}}
  >{\raggedright\arraybackslash}p{(\linewidth - 4\tabcolsep) * \real{0.3333}}@{}}
\caption{Principaux formats de données rencontrés en analyse de
données}\label{tbl-formats-donnees}\tabularnewline
\toprule\noalign{}
\begin{minipage}[b]{\linewidth}\raggedright
Format
\end{minipage} & \begin{minipage}[b]{\linewidth}\raggedright
Caractéristiques
\end{minipage} & \begin{minipage}[b]{\linewidth}\raggedright
Lecture avec \texttt{pandas}
\end{minipage} \\
\midrule\noalign{}
\endfirsthead
\toprule\noalign{}
\begin{minipage}[b]{\linewidth}\raggedright
Format
\end{minipage} & \begin{minipage}[b]{\linewidth}\raggedright
Caractéristiques
\end{minipage} & \begin{minipage}[b]{\linewidth}\raggedright
Lecture avec \texttt{pandas}
\end{minipage} \\
\midrule\noalign{}
\endhead
\bottomrule\noalign{}
\endlastfoot
\textbf{CSV} & Texte tabulaire, très répandu, lisible par un humain &
\texttt{pd.read\_csv()} \\
\textbf{JSON} & Structure hiérarchique (objets imbriqués), courant pour
les API & \texttt{pd.read\_json()} \\
\textbf{Excel (\texttt{.xlsx})} & Tableur, souvent utilisé en entreprise
& \texttt{pd.read\_excel()} \\
\textbf{Parquet} & Format binaire columnaire, compressé, optimisé pour
l'analytique & \texttt{pd.read\_parquet()} \\
\end{longtable}

\begin{tcolorbox}[enhanced jigsaw, arc=.35mm, bottomrule=.15mm, bottomtitle=1mm, breakable, colback=white, colbacktitle=quarto-callout-important-color!10!white, colframe=quarto-callout-important-color-frame, coltitle=black, left=2mm, leftrule=.75mm, opacityback=0, opacitybacktitle=0.6, rightrule=.15mm, title=\textcolor{quarto-callout-important-color}{\faExclamation}\hspace{0.5em}{Zoom sur les tendances : le format Parquet}, titlerule=0mm, toprule=.15mm, toptitle=1mm]

Contrairement au CSV, le format \textbf{Parquet} stocke les données par
colonne plutôt que par ligne. Cela permet de ne lire que les colonnes
réellement nécessaires à une analyse, de bénéficier d'une compression
efficace et de conserver les types de données (un entier reste un
entier, pas une chaîne de caractères à réinterpréter). Il est
aujourd'hui privilégié pour les jeux de données volumineux, et s'utilise
nativement avec \texttt{pandas}, \texttt{Polars} ou \texttt{DuckDB}.

\end{tcolorbox}

\subsection{1.3. Les API}\label{les-api}

Une \textbf{API} (\emph{Application Programming Interface}) permet
d'obtenir des données structurées de façon normalisée, sans avoir à
analyser du HTML :

\begin{itemize}
\tightlist
\item
  les \textbf{API REST}, les plus répandues, où chaque ressource est
  accessible via une URL et une méthode HTTP (\texttt{GET},
  \texttt{POST}\ldots), avec une réponse généralement au format JSON ;
\item
  les \textbf{API GraphQL}, plus récentes, qui permettent au client de
  préciser exactement les champs souhaités en une seule requête.
\end{itemize}

\begin{Shaded}
\begin{Highlighting}[]
\ImportTok{import}\NormalTok{ requests}
\ImportTok{import}\NormalTok{ pandas }\ImportTok{as}\NormalTok{ pd}

\NormalTok{reponse }\OperatorTok{=}\NormalTok{ requests.get(}\StringTok{"https://api.exemple.org/ventes?annee=2025"}\NormalTok{)}
\NormalTok{df }\OperatorTok{=}\NormalTok{ pd.json\_normalize(reponse.json()[}\StringTok{"resultats"}\NormalTok{])}
\end{Highlighting}
\end{Shaded}

\subsection{\texorpdfstring{1.4. Le \emph{web
scraping}}{1.4. Le web scraping}}\label{le-web-scraping}

Lorsque les données ne sont disponibles que sur des pages web, on
recourt au \textbf{web scraping} : l'extraction automatisée
d'informations à partir du code HTML. Les bibliothèques Python de
référence sont \texttt{BeautifulSoup} (analyse d'une page unique) et
\texttt{Scrapy} (collecte à grande échelle, avec gestion de la
pagination et des files d'attente).

\begin{tcolorbox}[enhanced jigsaw, arc=.35mm, bottomrule=.15mm, bottomtitle=1mm, breakable, colback=white, colbacktitle=quarto-callout-warning-color!10!white, colframe=quarto-callout-warning-color-frame, coltitle=black, left=2mm, leftrule=.75mm, opacityback=0, opacitybacktitle=0.6, rightrule=.15mm, title=\textcolor{quarto-callout-warning-color}{\faExclamationTriangle}\hspace{0.5em}{Éthique et légalité du web scraping}, titlerule=0mm, toprule=.15mm, toptitle=1mm]

Avant de collecter des données sur un site, il convient de vérifier son
fichier \texttt{robots.txt}, ses conditions générales d'utilisation,
ainsi que le cadre réglementaire applicable à la protection des données
personnelles (RGPD, ou la loi marocaine 09-08), en particulier si des
données à caractère personnel sont concernées.

\end{tcolorbox}

\subsection{1.5. Open data et capteurs /
IoT}\label{open-data-et-capteurs-iot}

Deux autres sources méritent d'être mentionnées :

\begin{itemize}
\tightlist
\item
  les \textbf{portails d'open data} (données publiques mises à
  disposition par des institutions, ex. données démographiques,
  économiques ou climatiques) ;
\item
  les \textbf{capteurs et objets connectés (IoT)}, qui génèrent des flux
  de données en continu (température, géolocalisation, fréquentation),
  souvent collectés via des files de messages (ex. MQTT, Kafka) avant
  d'être stockés pour l'analyse.
\end{itemize}

\section{2. Nettoyage des données}\label{nettoyage-des-donnuxe9es}

Une fois collectées, les données contiennent presque toujours des
imperfections. Le nettoyage est l'étape la plus chronophage d'un projet
d'analyse --- souvent estimée à 60-80 \% du temps réel passé sur un
projet.

Un nettoyage fiable vérifie au moins cinq dimensions :

\begin{enumerate}
\def\labelenumi{\arabic{enumi}.}
\tightlist
\item
  \textbf{complétude} : les informations importantes sont-elles
  présentes ?
\item
  \textbf{validité} : les valeurs respectent-elles les règles attendues
  (un prix positif, une date valide) ?
\item
  \textbf{unicité} : une même observation n'est-elle pas enregistrée
  deux fois ?
\item
  \textbf{cohérence} : les colonnes et les unités sont-elles compatibles
  entre elles ?
\item
  \textbf{actualité} : les dates couvrent-elles la période utile à la
  question ?
\end{enumerate}

Il faut garder une trace de chaque correction : colonne concernée, règle
appliquée, nombre de lignes modifiées et justification. Cette
documentation permet de revenir en arrière et d'expliquer les résultats.

\subsection{2.1. Les valeurs manquantes}\label{les-valeurs-manquantes}

Une valeur manquante peut avoir des origines très différentes, ce qui
conditionne la stratégie à adopter :

\begin{longtable}[]{@{}
  >{\raggedright\arraybackslash}p{(\linewidth - 4\tabcolsep) * \real{0.3333}}
  >{\raggedright\arraybackslash}p{(\linewidth - 4\tabcolsep) * \real{0.3333}}
  >{\raggedright\arraybackslash}p{(\linewidth - 4\tabcolsep) * \real{0.3333}}@{}}
\caption{Typologie des valeurs
manquantes}\label{tbl-typologie-manquantes}\tabularnewline
\toprule\noalign{}
\begin{minipage}[b]{\linewidth}\raggedright
Type
\end{minipage} & \begin{minipage}[b]{\linewidth}\raggedright
Signification
\end{minipage} & \begin{minipage}[b]{\linewidth}\raggedright
Exemple
\end{minipage} \\
\midrule\noalign{}
\endfirsthead
\toprule\noalign{}
\begin{minipage}[b]{\linewidth}\raggedright
Type
\end{minipage} & \begin{minipage}[b]{\linewidth}\raggedright
Signification
\end{minipage} & \begin{minipage}[b]{\linewidth}\raggedright
Exemple
\end{minipage} \\
\midrule\noalign{}
\endhead
\bottomrule\noalign{}
\endlastfoot
\textbf{MCAR} (\emph{Missing Completely At Random}) & L'absence ne
dépend d'aucune variable, ni observée ni cachée & Un capteur tombe en
panne aléatoirement \\
\textbf{MAR} (\emph{Missing At Random}) & L'absence dépend d'une autre
variable observée & Les hommes répondent moins souvent à une question
sur le revenu \\
\textbf{MNAR} (\emph{Missing Not At Random}) & L'absence dépend de la
valeur manquante elle-même & Les personnes aux revenus très élevés
refusent de répondre \\
\end{longtable}

Ces trois étiquettes servent surtout à réfléchir à la cause d'une
absence. Il n'est pas demandé de réaliser un test statistique pour les
distinguer dans ce cours : l'étudiant doit surtout éviter d'imputer
automatiquement sans comprendre ce que signifie la valeur manquante.

Avec \texttt{pandas}, on détecte et on traite les valeurs manquantes
ainsi :

\begin{Shaded}
\begin{Highlighting}[]
\CommentTok{\# Détection}
\NormalTok{df.isna().}\BuiltInTok{sum}\NormalTok{()                      }\CommentTok{\# nombre de valeurs manquantes par colonne}
\NormalTok{df.isna().mean().}\BuiltInTok{round}\NormalTok{(}\DecValTok{2}\NormalTok{) }\OperatorTok{*} \DecValTok{100}      \CommentTok{\# pourcentage de valeurs manquantes}

\CommentTok{\# Suppression}
\NormalTok{df.dropna(subset}\OperatorTok{=}\NormalTok{[}\StringTok{"prix"}\NormalTok{])           }\CommentTok{\# supprime les lignes sans prix}

\CommentTok{\# Imputation}
\NormalTok{df[}\StringTok{"surface"}\NormalTok{] }\OperatorTok{=}\NormalTok{ df[}\StringTok{"surface"}\NormalTok{].fillna(df[}\StringTok{"surface"}\NormalTok{].median())}
\NormalTok{df[}\StringTok{"quartier"}\NormalTok{] }\OperatorTok{=}\NormalTok{ df[}\StringTok{"quartier"}\NormalTok{].fillna(}\StringTok{"Non renseigné"}\NormalTok{)}
\end{Highlighting}
\end{Shaded}

\begin{tcolorbox}[enhanced jigsaw, arc=.35mm, bottomrule=.15mm, bottomtitle=1mm, breakable, colback=white, colbacktitle=quarto-callout-tip-color!10!white, colframe=quarto-callout-tip-color-frame, coltitle=black, left=2mm, leftrule=.75mm, opacityback=0, opacitybacktitle=0.6, rightrule=.15mm, title=\textcolor{quarto-callout-tip-color}{\faLightbulb}\hspace{0.5em}{Comment choisir entre suppression et imputation ?}, titlerule=0mm, toprule=.15mm, toptitle=1mm]

\begin{itemize}
\tightlist
\item
  Si la proportion de valeurs manquantes est \textbf{faible} (quelques
  pourcents) et que le mécanisme semble MCAR, la \textbf{suppression}
  des lignes concernées est souvent acceptable.
\item
  Si la proportion est \textbf{élevée}, ou si l'absence porte sur une
  variable centrale de l'analyse, on préfère une \textbf{imputation}
  (moyenne, médiane, mode, ou une valeur explicite comme « Non renseigné
  »), en gardant à l'esprit qu'elle introduit un biais qu'il faut
  documenter.
\end{itemize}

\end{tcolorbox}

\subsection{\texorpdfstring{2.2. Les valeurs aberrantes
(\emph{outliers})}{2.2. Les valeurs aberrantes (outliers)}}\label{les-valeurs-aberrantes-outliers}

Une valeur aberrante est une observation qui s'écarte fortement du reste
de la distribution. Deux méthodes usuelles permettent de les détecter :

\begin{itemize}
\tightlist
\item
  la méthode du \textbf{score Z}, qui signale une observation aberrante
  lorsque \(|z| = \left|\dfrac{x - \bar{x}}{\sigma}\right| > 3\) ;
\item
  la méthode de l'\textbf{écart interquartile (IQR)}, plus robuste aux
  distributions asymétriques : on considère comme aberrante toute valeur
  en dehors de l'intervalle
  \([Q_1 - 1{,}5 \times IQR,\ Q_3 + 1{,}5 \times IQR]\), avec
  \(IQR = Q_3 - Q_1\).
\end{itemize}

\begin{Shaded}
\begin{Highlighting}[]
\NormalTok{q1, q3 }\OperatorTok{=}\NormalTok{ df[}\StringTok{"prix"}\NormalTok{].quantile([}\FloatTok{0.25}\NormalTok{, }\FloatTok{0.75}\NormalTok{])}
\NormalTok{iqr }\OperatorTok{=}\NormalTok{ q3 }\OperatorTok{{-}}\NormalTok{ q1}
\NormalTok{borne\_basse, borne\_haute }\OperatorTok{=}\NormalTok{ q1 }\OperatorTok{{-}} \FloatTok{1.5} \OperatorTok{*}\NormalTok{ iqr, q3 }\OperatorTok{+} \FloatTok{1.5} \OperatorTok{*}\NormalTok{ iqr}

\NormalTok{aberrants }\OperatorTok{=}\NormalTok{ df[(df[}\StringTok{"prix"}\NormalTok{] }\OperatorTok{\textless{}}\NormalTok{ borne\_basse) }\OperatorTok{|}\NormalTok{ (df[}\StringTok{"prix"}\NormalTok{] }\OperatorTok{\textgreater{}}\NormalTok{ borne\_haute)]}
\end{Highlighting}
\end{Shaded}

La \textbf{boîte à moustaches} (\emph{boxplot}), présentée au chapitre
suivant sur la visualisation, est l'outil graphique privilégié pour
repérer visuellement ces valeurs.

\begin{tcolorbox}[enhanced jigsaw, arc=.35mm, bottomrule=.15mm, bottomtitle=1mm, breakable, colback=white, colbacktitle=quarto-callout-warning-color!10!white, colframe=quarto-callout-warning-color-frame, coltitle=black, left=2mm, leftrule=.75mm, opacityback=0, opacitybacktitle=0.6, rightrule=.15mm, title=\textcolor{quarto-callout-warning-color}{\faExclamationTriangle}\hspace{0.5em}{Une valeur aberrante n'est pas toujours une erreur}, titlerule=0mm, toprule=.15mm, toptitle=1mm]

Avant de supprimer une valeur aberrante, il faut se demander si elle
correspond à une \textbf{erreur de saisie} (à corriger ou supprimer) ou
à un \textbf{phénomène réel mais rare} (par exemple une vente
exceptionnelle), qu'il peut être pertinent de conserver, voire d'étudier
spécifiquement.

\end{tcolorbox}

\subsection{2.3. Doublons et
incohérences}\label{doublons-et-incohuxe9rences}

\begin{Shaded}
\begin{Highlighting}[]
\NormalTok{df.duplicated().}\BuiltInTok{sum}\NormalTok{()                     }\CommentTok{\# nombre de lignes strictement dupliquées}
\NormalTok{df }\OperatorTok{=}\NormalTok{ df.drop\_duplicates()}

\NormalTok{df[}\StringTok{"ville"}\NormalTok{] }\OperatorTok{=}\NormalTok{ df[}\StringTok{"ville"}\NormalTok{].}\BuiltInTok{str}\NormalTok{.strip().}\BuiltInTok{str}\NormalTok{.title()   }\CommentTok{\# espaces superflus, casse}
\NormalTok{df[}\StringTok{"date"}\NormalTok{] }\OperatorTok{=}\NormalTok{ pd.to\_datetime(df[}\StringTok{"date"}\NormalTok{], errors}\OperatorTok{=}\StringTok{"coerce"}\NormalTok{)  }\CommentTok{\# type de données cohérent}
\end{Highlighting}
\end{Shaded}

Les incohérences les plus fréquentes concernent la \textbf{casse} et les
\textbf{espaces superflus} dans les chaînes de caractères (« casablanca
» vs « Casablanca »), les \textbf{types de données} mal interprétés (une
date lue comme texte), et les \textbf{unités hétérogènes} (des prix
exprimés tantôt en MAD, tantôt en euros).

\subsection{2.4. Vers une démarche de qualité des
données}\label{vers-une-duxe9marche-de-qualituxe9-des-donnuxe9es}

\begin{tcolorbox}[enhanced jigsaw, arc=.35mm, bottomrule=.15mm, bottomtitle=1mm, breakable, colback=white, colbacktitle=quarto-callout-important-color!10!white, colframe=quarto-callout-important-color-frame, coltitle=black, left=2mm, leftrule=.75mm, opacityback=0, opacitybacktitle=0.6, rightrule=.15mm, title=\textcolor{quarto-callout-important-color}{\faExclamation}\hspace{0.5em}{Zoom sur les tendances : la validation systématique de la qualité}, titlerule=0mm, toprule=.15mm, toptitle=1mm]

Sur des projets d'analyse récurrents (tableaux de bord alimentés
quotidiennement, par exemple), il devient risqué de vérifier la qualité
des données « à l'œil » à chaque exécution. Des bibliothèques comme
\textbf{\href{https://greatexpectations.io}{\texttt{Great\ Expectations}}}
permettent de définir des règles de validation explicites (« la colonne
\texttt{prix} doit toujours être positive », « la colonne \texttt{date}
ne doit jamais être manquante ») et de les exécuter automatiquement à
chaque nouvelle collecte de données, afin de détecter une dérive de
qualité avant qu'elle n'affecte l'analyse.

\end{tcolorbox}

\section{3. Transformation des
données}\label{transformation-des-donnuxe9es}

Une fois nettoyées, les données doivent souvent être reformatées avant
l'exploration ou la visualisation.

\subsection{3.1. Encodage des variables
catégorielles}\label{encodage-des-variables-catuxe9gorielles}

De nombreux outils statistiques et graphiques manipulent plus facilement
des valeurs numériques que du texte. On distingue deux stratégies
principales :

\begin{itemize}
\tightlist
\item
  l'\textbf{encodage ordinal}, adapté aux catégories qui possèdent un
  ordre naturel (ex. « Faible » \textless{} « Moyen » \textless{} «
  Élevé ») ;
\item
  l'\textbf{encodage one-hot} (variables indicatrices), adapté aux
  catégories sans ordre (ex. les quartiers d'une ville), qui crée une
  colonne binaire par catégorie.
\end{itemize}

\begin{Shaded}
\begin{Highlighting}[]
\CommentTok{\# Encodage ordinal}
\NormalTok{ordre }\OperatorTok{=}\NormalTok{ \{}\StringTok{"Faible"}\NormalTok{: }\DecValTok{0}\NormalTok{, }\StringTok{"Moyen"}\NormalTok{: }\DecValTok{1}\NormalTok{, }\StringTok{"Élevé"}\NormalTok{: }\DecValTok{2}\NormalTok{\}}
\NormalTok{df[}\StringTok{"niveau\_encode"}\NormalTok{] }\OperatorTok{=}\NormalTok{ df[}\StringTok{"niveau"}\NormalTok{].}\BuiltInTok{map}\NormalTok{(ordre)}

\CommentTok{\# Encodage one{-}hot}
\NormalTok{df\_encode }\OperatorTok{=}\NormalTok{ pd.get\_dummies(df, columns}\OperatorTok{=}\NormalTok{[}\StringTok{"quartier"}\NormalTok{])}
\end{Highlighting}
\end{Shaded}

\begin{tcolorbox}[enhanced jigsaw, arc=.35mm, bottomrule=.15mm, bottomtitle=1mm, breakable, colback=white, colbacktitle=quarto-callout-note-color!10!white, colframe=quarto-callout-note-color-frame, coltitle=black, left=2mm, leftrule=.75mm, opacityback=0, opacitybacktitle=0.6, rightrule=.15mm, title=\textcolor{quarto-callout-note-color}{\faInfo}\hspace{0.5em}{Note}, titlerule=0mm, toprule=.15mm, toptitle=1mm]

L'encodage détaillé des variables (avec \texttt{scikit-learn}, gestion
des catégories rares, etc.) sera approfondi dans le module de Machine
Learning. Dans ce cours, il est surtout utile pour faciliter certains
graphiques et calculs de corrélation présentés au chapitre suivant.

\end{tcolorbox}

\subsection{3.2. Normalisation et
standardisation}\label{normalisation-et-standardisation}

Lorsque plusieurs variables numériques ont des échelles très différentes
(ex. une surface en \(m^2\) et un prix en millions de MAD), il est
parfois utile de les ramener à une échelle comparable, notamment pour
certaines visualisations combinées :

\begin{itemize}
\tightlist
\item
  la \textbf{normalisation min-max} ramène chaque valeur dans
  l'intervalle \([0, 1]\) :
  \(x' = \dfrac{x - x_{min}}{x_{max} - x_{min}}\) ;
\item
  la \textbf{standardisation} (score Z) centre la variable autour de 0
  avec un écart-type de 1 : \(x' = \dfrac{x - \bar{x}}{\sigma}\).
\end{itemize}

\begin{Shaded}
\begin{Highlighting}[]
\NormalTok{df[}\StringTok{"prix\_norm"}\NormalTok{] }\OperatorTok{=}\NormalTok{ (df[}\StringTok{"prix"}\NormalTok{] }\OperatorTok{{-}}\NormalTok{ df[}\StringTok{"prix"}\NormalTok{].}\BuiltInTok{min}\NormalTok{()) }\OperatorTok{/}\NormalTok{ (df[}\StringTok{"prix"}\NormalTok{].}\BuiltInTok{max}\NormalTok{() }\OperatorTok{{-}}\NormalTok{ df[}\StringTok{"prix"}\NormalTok{].}\BuiltInTok{min}\NormalTok{())}
\NormalTok{df[}\StringTok{"prix\_standard"}\NormalTok{] }\OperatorTok{=}\NormalTok{ (df[}\StringTok{"prix"}\NormalTok{] }\OperatorTok{{-}}\NormalTok{ df[}\StringTok{"prix"}\NormalTok{].mean()) }\OperatorTok{/}\NormalTok{ df[}\StringTok{"prix"}\NormalTok{].std()}
\end{Highlighting}
\end{Shaded}

\subsection{\texorpdfstring{3.3. Discrétisation
(\emph{binning})}{3.3. Discrétisation (binning)}}\label{discruxe9tisation-binning}

La discrétisation consiste à transformer une variable numérique continue
en catégories, souvent utile pour simplifier une visualisation (par
exemple des tranches d'âge plutôt qu'un âge exact) :

\begin{Shaded}
\begin{Highlighting}[]
\NormalTok{df[}\StringTok{"tranche\_age"}\NormalTok{] }\OperatorTok{=}\NormalTok{ pd.cut(}
\NormalTok{    df[}\StringTok{"age"}\NormalTok{],}
\NormalTok{    bins}\OperatorTok{=}\NormalTok{[}\DecValTok{0}\NormalTok{, }\DecValTok{18}\NormalTok{, }\DecValTok{35}\NormalTok{, }\DecValTok{60}\NormalTok{, }\DecValTok{120}\NormalTok{],}
\NormalTok{    labels}\OperatorTok{=}\NormalTok{[}\StringTok{"Mineur"}\NormalTok{, }\StringTok{"Jeune adulte"}\NormalTok{, }\StringTok{"Adulte"}\NormalTok{, }\StringTok{"Senior"}\NormalTok{]}
\NormalTok{)}
\end{Highlighting}
\end{Shaded}

\subsection{3.4. Quelques transformations usuelles
supplémentaires}\label{quelques-transformations-usuelles-suppluxe9mentaires}

\begin{itemize}
\tightlist
\item
  \textbf{Extraction de composantes temporelles} : à partir d'une
  colonne \texttt{date}, on extrait fréquemment l'année, le mois, le
  jour de semaine (\texttt{df{[}"date"{]}.dt.month}), utiles pour
  l'exploration de séries temporelles.
\item
  \textbf{Agrégats dérivés} : création de nouvelles colonnes résumant
  une information (ex. le prix au mètre carré =
  \texttt{prix\ /\ surface}), qui facilitent grandement l'exploration et
  la comparaison.
\end{itemize}

\begin{tcolorbox}[enhanced jigsaw, arc=.35mm, bottomrule=.15mm, bottomtitle=1mm, breakable, colback=white, colbacktitle=quarto-callout-warning-color!10!white, colframe=quarto-callout-warning-color-frame, coltitle=black, left=2mm, leftrule=.75mm, opacityback=0, opacitybacktitle=0.6, rightrule=.15mm, title=\textcolor{quarto-callout-warning-color}{\faExclamationTriangle}\hspace{0.5em}{Ne pas transformer sans objectif}, titlerule=0mm, toprule=.15mm, toptitle=1mm]

Une transformation n'est pas automatiquement une amélioration. Par
exemple, remplacer toutes les valeurs manquantes par zéro change la
signification de la variable : « absence d'information » devient «
valeur nulle ». Avant chaque transformation, demandez-vous quelle
question elle facilite et si son effet sera documenté.

\end{tcolorbox}

\begin{tcolorbox}[enhanced jigsaw, arc=.35mm, bottomrule=.15mm, bottomtitle=1mm, breakable, colback=white, colbacktitle=quarto-callout-note-color!10!white, colframe=quarto-callout-note-color-frame, coltitle=black, left=2mm, leftrule=.75mm, opacityback=0, opacitybacktitle=0.6, rightrule=.15mm, title=\textcolor{quarto-callout-note-color}{\faInfo}\hspace{0.5em}{Auto-évaluation}, titlerule=0mm, toprule=.15mm, toptitle=1mm]

Avant de passer au chapitre suivant, assurez-vous de pouvoir répondre «
oui » aux questions suivantes :

\begin{itemize}
\tightlist
\item[$\square$]
  Je sais citer trois sources possibles de collecte de données et un cas
  d'usage pour chacune.
\item[$\square$]
  Je sais distinguer les valeurs manquantes MCAR, MAR et MNAR et adapter
  ma stratégie en conséquence.
\item[$\square$]
  Je sais détecter des valeurs aberrantes avec la méthode de l'IQR.
\item[$\square$]
  Je sais choisir entre encodage ordinal et one-hot selon la nature
  d'une variable catégorielle.
\end{itemize}

\end{tcolorbox}

\section{Exercices}\label{exercices-1}

\begin{exercise}[Choisir une stratégie de traitement des valeurs
manquantes]\protect\hypertarget{exr-prep-1}{}\label{exr-prep-1}

Un jeu de données de candidatures à un emploi comporte une colonne
\texttt{salaire\_souhaite} manquante pour 40 \% des lignes. Une analyse
rapide montre que les candidats les plus expérimentés (souvent plus
exigeants sur le salaire) sont surreprésentés parmi les valeurs
manquantes.

\begin{enumerate}
\def\labelenumi{\alph{enumi}.}
\tightlist
\item
  À quel type de valeur manquante (Table~\ref{tbl-typologie-manquantes})
  correspond cette situation ?
\item
  Recommanderiez-vous une suppression ou une imputation ? Justifiez.
\end{enumerate}

\end{exercise}

\begin{tcolorbox}[enhanced jigsaw, arc=.35mm, bottomrule=.15mm, bottomtitle=1mm, breakable, colback=white, colbacktitle=quarto-callout-tip-color!10!white, colframe=quarto-callout-tip-color-frame, coltitle=black, left=2mm, leftrule=.75mm, opacityback=0, opacitybacktitle=0.6, rightrule=.15mm, title=\textcolor{quarto-callout-tip-color}{\faLightbulb}\hspace{0.5em}{Solution de l'exercice Exercice~\ref{exr-prep-1}}, titlerule=0mm, toprule=.15mm, toptitle=1mm]

\begin{enumerate}
\def\labelenumi{\alph{enumi}.}
\tightlist
\item
  Il s'agit d'un cas \textbf{MNAR} : l'absence de réponse dépend de la
  valeur elle-même (les candidats avec un salaire souhaité élevé
  répondent moins).
\item
  Ni la suppression pure ni une imputation par la moyenne ne sont
  satisfaisantes ici, car elles biaiseraient l'analyse à la baisse. Il
  est préférable de \textbf{documenter explicitement} cette non-réponse
  (par exemple une catégorie « Non communiqué ») plutôt que de l'imputer
  silencieusement, et d'en tenir compte dans l'interprétation des
  résultats.
\end{enumerate}

\end{tcolorbox}

\begin{exercise}[Détection
d'outliers]\protect\hypertarget{exr-prep-2}{}\label{exr-prep-2}

Pour une variable \texttt{prix} avec \(Q_1 = 200\,000\) et
\(Q_3 = 450\,000\) (en MAD), calculez les bornes de détection des
valeurs aberrantes selon la méthode de l'IQR.

\end{exercise}

\begin{tcolorbox}[enhanced jigsaw, arc=.35mm, bottomrule=.15mm, bottomtitle=1mm, breakable, colback=white, colbacktitle=quarto-callout-tip-color!10!white, colframe=quarto-callout-tip-color-frame, coltitle=black, left=2mm, leftrule=.75mm, opacityback=0, opacitybacktitle=0.6, rightrule=.15mm, title=\textcolor{quarto-callout-tip-color}{\faLightbulb}\hspace{0.5em}{Solution de l'exercice Exercice~\ref{exr-prep-2}}, titlerule=0mm, toprule=.15mm, toptitle=1mm]

\(IQR = 450\,000 - 200\,000 = 250\,000\).

\begin{itemize}
\tightlist
\item
  Borne basse : \(200\,000 - 1{,}5 \times 250\,000 = -175\,000\) (donc,
  en pratique, aucune valeur basse ne sera jamais considérée comme
  aberrante ici, un prix ne pouvant être négatif).
\item
  Borne haute : \(450\,000 + 1{,}5 \times 250\,000 = 825\,000\).
\end{itemize}

Tout bien affiché à plus de 825 000 MAD serait donc signalé comme
potentiellement aberrant.

\end{tcolorbox}

\begin{exercise}[Encodage]\protect\hypertarget{exr-prep-3}{}\label{exr-prep-3}

Pour chacune des variables suivantes, indiquez si un encodage
\textbf{ordinal} ou \textbf{one-hot} est le plus adapté, et pourquoi :

\begin{enumerate}
\def\labelenumi{\alph{enumi}.}
\tightlist
\item
  Le niveau d'étude (« Bac », « Bac+2 », « Bac+3 », « Bac+5 ») ;
\item
  La couleur préférée d'un client (« Rouge », « Bleu », « Vert »).
\end{enumerate}

\end{exercise}

\begin{tcolorbox}[enhanced jigsaw, arc=.35mm, bottomrule=.15mm, bottomtitle=1mm, breakable, colback=white, colbacktitle=quarto-callout-tip-color!10!white, colframe=quarto-callout-tip-color-frame, coltitle=black, left=2mm, leftrule=.75mm, opacityback=0, opacitybacktitle=0.6, rightrule=.15mm, title=\textcolor{quarto-callout-tip-color}{\faLightbulb}\hspace{0.5em}{Solution de l'exercice Exercice~\ref{exr-prep-3}}, titlerule=0mm, toprule=.15mm, toptitle=1mm]

\begin{enumerate}
\def\labelenumi{\alph{enumi}.}
\tightlist
\item
  \textbf{Encodage ordinal} : les niveaux d'étude possèdent un ordre
  naturel croissant.
\item
  \textbf{Encodage one-hot} : les couleurs n'ont pas d'ordre logique
  entre elles ; leur attribuer un ordre arbitraire introduirait une
  fausse relation numérique.
\end{enumerate}

\end{tcolorbox}

\begin{exercise}[Diagnostic de
qualité]\protect\hypertarget{exr-prep-4}{}\label{exr-prep-4}

Un fichier de ventes contient 1 000 lignes. La colonne \texttt{date}
comporte 20 valeurs manquantes, \texttt{prix} contient 5 valeurs
négatives et 12 lignes sont des doublons exacts.

\begin{enumerate}
\def\labelenumi{\alph{enumi}.}
\tightlist
\item
  Quel pourcentage de dates est manquant ?
\item
  Quel problème doit être vérifié en priorité avant de calculer le
  chiffre d'affaires ?
\item
  Les doublons doivent-ils être supprimés automatiquement ? Justifiez.
\end{enumerate}

\end{exercise}

\begin{tcolorbox}[enhanced jigsaw, arc=.35mm, bottomrule=.15mm, bottomtitle=1mm, breakable, colback=white, colbacktitle=quarto-callout-tip-color!10!white, colframe=quarto-callout-tip-color-frame, coltitle=black, left=2mm, leftrule=.75mm, opacityback=0, opacitybacktitle=0.6, rightrule=.15mm, title=\textcolor{quarto-callout-tip-color}{\faLightbulb}\hspace{0.5em}{Solution de l'exercice Exercice~\ref{exr-prep-4}}, titlerule=0mm, toprule=.15mm, toptitle=1mm]

\begin{enumerate}
\def\labelenumi{\alph{enumi}.}
\tightlist
\item
  \(20/1000 \times 100 = 2\%\).
\item
  Les \textbf{prix négatifs}, car ils rendent le total incohérent et
  peuvent signaler une règle métier particulière (remboursement) ou une
  erreur.
\item
  Non. Il faut vérifier si les lignes sont réellement identiques ou si
  elles représentent deux transactions légitimes saisies avec les mêmes
  valeurs. Une suppression aveugle pourrait supprimer de vraies ventes.
\end{enumerate}

\end{tcolorbox}

\begin{exercise}[Choisir une stratégie
d'imputation]\protect\hypertarget{exr-prep-5}{}\label{exr-prep-5}

Associez chaque situation à une stratégie raisonnable : supprimer
quelques lignes, utiliser la médiane, créer la catégorie « Non renseigné
» ou ne pas imputer sans investigation.

\begin{enumerate}
\def\labelenumi{\alph{enumi}.}
\tightlist
\item
  1 \% des lignes n'ont pas de surface, et la surface est numérique.
\item
  35 \% des lignes n'ont pas de quartier.
\item
  Une date manquante empêche de savoir à quelle période appartient une
  vente.
\item
  Une variable \texttt{remise} a 2 \% de valeurs manquantes dans une
  distribution très asymétrique.
\end{enumerate}

\end{exercise}

\begin{tcolorbox}[enhanced jigsaw, arc=.35mm, bottomrule=.15mm, bottomtitle=1mm, breakable, colback=white, colbacktitle=quarto-callout-tip-color!10!white, colframe=quarto-callout-tip-color-frame, coltitle=black, left=2mm, leftrule=.75mm, opacityback=0, opacitybacktitle=0.6, rightrule=.15mm, title=\textcolor{quarto-callout-tip-color}{\faLightbulb}\hspace{0.5em}{Solution de l'exercice Exercice~\ref{exr-prep-5}}, titlerule=0mm, toprule=.15mm, toptitle=1mm]

\begin{enumerate}
\def\labelenumi{\alph{enumi}.}
\tightlist
\item
  \textbf{Médiane} (ou suppression des rares lignes si elles ne sont pas
  centrales), après vérification.
\item
  \textbf{« Non renseigné »}, afin de ne pas perdre un tiers du jeu de
  données.
\item
  \textbf{Ne pas imputer sans investigation} : il faut retrouver la date
  dans la source ou exclure la vente d'une analyse temporelle.
\item
  \textbf{Médiane} plutôt que moyenne, car elle est plus robuste à
  l'asymétrie ; la décision doit être documentée.
\end{enumerate}

\end{tcolorbox}

\begin{exercise}[Calculer une normalisation
min-max]\protect\hypertarget{exr-prep-6}{}\label{exr-prep-6}

Une variable \texttt{score} varie de 20 à 80. Calculez la valeur
normalisée d'un score égal à 50 avec :

\[x'=\frac{x-x_{\min}}{x_{\max}-x_{\min}}.\]

Que signifie le résultat ?

\end{exercise}

\begin{tcolorbox}[enhanced jigsaw, arc=.35mm, bottomrule=.15mm, bottomtitle=1mm, breakable, colback=white, colbacktitle=quarto-callout-tip-color!10!white, colframe=quarto-callout-tip-color-frame, coltitle=black, left=2mm, leftrule=.75mm, opacityback=0, opacitybacktitle=0.6, rightrule=.15mm, title=\textcolor{quarto-callout-tip-color}{\faLightbulb}\hspace{0.5em}{Solution de l'exercice Exercice~\ref{exr-prep-6}}, titlerule=0mm, toprule=.15mm, toptitle=1mm]

\(x'=(50-20)/(80-20)=30/60=\textbf{0,5}\). Le score 50 se situe au
milieu de l'intervalle observé. La normalisation ne signifie pas que 50
est la moyenne : elle indique sa position entre le minimum et le
maximum.

\end{tcolorbox}

\begin{exercise}[Unités et
cohérence]\protect\hypertarget{exr-prep-7}{}\label{exr-prep-7}

Une colonne \texttt{distance} contient les valeurs 5, 12 et 8 pour trois
trajets, mais le dictionnaire indique que certaines lignes sont en
kilomètres et d'autres en mètres.

\begin{enumerate}
\def\labelenumi{\alph{enumi}.}
\tightlist
\item
  Pourquoi est-il dangereux de calculer une moyenne directement ?
\item
  Quelle action doit être réalisée avant l'analyse ?
\item
  Est-ce un problème de valeur manquante, de validité ou de cohérence ?
\end{enumerate}

\end{exercise}

\begin{tcolorbox}[enhanced jigsaw, arc=.35mm, bottomrule=.15mm, bottomtitle=1mm, breakable, colback=white, colbacktitle=quarto-callout-tip-color!10!white, colframe=quarto-callout-tip-color-frame, coltitle=black, left=2mm, leftrule=.75mm, opacityback=0, opacitybacktitle=0.6, rightrule=.15mm, title=\textcolor{quarto-callout-tip-color}{\faLightbulb}\hspace{0.5em}{Solution de l'exercice Exercice~\ref{exr-prep-7}}, titlerule=0mm, toprule=.15mm, toptitle=1mm]

\begin{enumerate}
\def\labelenumi{\alph{enumi}.}
\tightlist
\item
  Les nombres n'ont pas la même unité : la moyenne mélangerait des
  grandeurs incomparables.
\item
  Identifier l'unité de chaque ligne puis convertir toutes les valeurs
  dans une unité commune, par exemple le kilomètre.
\item
  C'est principalement un problème de \textbf{cohérence} ; certaines
  valeurs peuvent aussi être invalides si leur unité est inconnue.
\end{enumerate}

\end{tcolorbox}

\begin{exercise}[Discrétisation et
interprétation]\protect\hypertarget{exr-prep-8}{}\label{exr-prep-8}

On regroupe les âges selon les intervalles \texttt{{[}0,18{[}},
\texttt{{[}18,35{[}}, \texttt{{[}35,60{[}} et \texttt{{[}60,120{]}}.

\begin{enumerate}
\def\labelenumi{\alph{enumi}.}
\tightlist
\item
  Dans quelle classe place-t-on les âges 18, 34, 35 et 60 ?
\item
  Quel avantage apporte cette transformation pour un tableau de bord ?
\item
  Quel risque faut-il signaler ?
\end{enumerate}

\end{exercise}

\begin{tcolorbox}[enhanced jigsaw, arc=.35mm, bottomrule=.15mm, bottomtitle=1mm, breakable, colback=white, colbacktitle=quarto-callout-tip-color!10!white, colframe=quarto-callout-tip-color-frame, coltitle=black, left=2mm, leftrule=.75mm, opacityback=0, opacitybacktitle=0.6, rightrule=.15mm, title=\textcolor{quarto-callout-tip-color}{\faLightbulb}\hspace{0.5em}{Solution de l'exercice Exercice~\ref{exr-prep-8}}, titlerule=0mm, toprule=.15mm, toptitle=1mm]

\begin{enumerate}
\def\labelenumi{\alph{enumi}.}
\tightlist
\item
  18 et 34 vont dans \texttt{{[}18,35{[}}, 35 dans \texttt{{[}35,60{[}},
  et 60 dans \texttt{{[}60,120{]}}.
\item
  Elle simplifie la lecture et permet de comparer des groupes d'âge.
\item
  Elle fait perdre de la précision : deux personnes de 18 et 34 ans sont
  dans le même groupe malgré un écart important. Les bornes doivent donc
  être documentées.
\end{enumerate}

\end{tcolorbox}

\bookmarksetup{startatroot}

\chapter{Exploration des données}\label{exploration-des-donnuxe9es}

\begin{tcolorbox}[enhanced jigsaw, arc=.35mm, bottomrule=.15mm, bottomtitle=1mm, breakable, colback=white, colbacktitle=quarto-callout-tip-color!10!white, colframe=quarto-callout-tip-color-frame, coltitle=black, left=2mm, leftrule=.75mm, opacityback=0, opacitybacktitle=0.6, rightrule=.15mm, title=\textcolor{quarto-callout-tip-color}{\faLightbulb}\hspace{0.5em}{Objectifs du chapitre}, titlerule=0mm, toprule=.15mm, toptitle=1mm]

À la fin de ce chapitre, l'étudiant sera capable de :

\begin{itemize}
\tightlist
\item
  \textbf{Calculer et interpréter} les principales statistiques
  descriptives (tendance centrale, dispersion, forme) ;
\item
  \textbf{Distinguer} l'analyse univariée, bivariée et multivariée ;
\item
  \textbf{Mesurer} une corrélation ou une association entre deux
  variables et en connaître les limites ;
\item
  \textbf{Construire} des tableaux d'agrégation (\texttt{groupby},
  tableaux croisés dynamiques) pour résumer un jeu de données ;
\item
  \textbf{Utiliser} un outil de profilage automatique pour accélérer une
  première exploration.
\end{itemize}

\end{tcolorbox}

L'\textbf{analyse exploratoire de données} (\emph{Exploratory Data
Analysis}, EDA), popularisée par le statisticien John Tukey, consiste à
examiner un jeu de données déjà nettoyé afin d'en dégager les tendances
principales, les relations entre variables et les éventuelles anomalies
restantes --- avant toute conclusion définitive ou toute visualisation
destinée à un public externe.

\section{1. Statistiques descriptives}\label{statistiques-descriptives}

\subsection{1.1. Tendance centrale}\label{tendance-centrale}

\begin{longtable}[]{@{}
  >{\raggedright\arraybackslash}p{(\linewidth - 4\tabcolsep) * \real{0.3333}}
  >{\raggedright\arraybackslash}p{(\linewidth - 4\tabcolsep) * \real{0.3333}}
  >{\raggedright\arraybackslash}p{(\linewidth - 4\tabcolsep) * \real{0.3333}}@{}}
\caption{Indicateurs de tendance
centrale}\label{tbl-tendance-centrale}\tabularnewline
\toprule\noalign{}
\begin{minipage}[b]{\linewidth}\raggedright
Indicateur
\end{minipage} & \begin{minipage}[b]{\linewidth}\raggedright
Définition
\end{minipage} & \begin{minipage}[b]{\linewidth}\raggedright
Sensibilité aux valeurs extrêmes
\end{minipage} \\
\midrule\noalign{}
\endfirsthead
\toprule\noalign{}
\begin{minipage}[b]{\linewidth}\raggedright
Indicateur
\end{minipage} & \begin{minipage}[b]{\linewidth}\raggedright
Définition
\end{minipage} & \begin{minipage}[b]{\linewidth}\raggedright
Sensibilité aux valeurs extrêmes
\end{minipage} \\
\midrule\noalign{}
\endhead
\bottomrule\noalign{}
\endlastfoot
\textbf{Moyenne} & Somme des valeurs divisée par leur nombre & Élevée \\
\textbf{Médiane} & Valeur qui sépare la distribution en deux moitiés
égales & Faible \\
\textbf{Mode} & Valeur la plus fréquente & Faible (mais peu informatif
pour une variable continue) \\
\end{longtable}

\begin{Shaded}
\begin{Highlighting}[]
\NormalTok{df[}\StringTok{"prix"}\NormalTok{].mean()}
\NormalTok{df[}\StringTok{"prix"}\NormalTok{].median()}
\NormalTok{df[}\StringTok{"quartier"}\NormalTok{].mode()[}\DecValTok{0}\NormalTok{]}
\end{Highlighting}
\end{Shaded}

\begin{tcolorbox}[enhanced jigsaw, arc=.35mm, bottomrule=.15mm, bottomtitle=1mm, breakable, colback=white, colbacktitle=quarto-callout-warning-color!10!white, colframe=quarto-callout-warning-color-frame, coltitle=black, left=2mm, leftrule=.75mm, opacityback=0, opacitybacktitle=0.6, rightrule=.15mm, title=\textcolor{quarto-callout-warning-color}{\faExclamationTriangle}\hspace{0.5em}{Moyenne ou médiane ?}, titlerule=0mm, toprule=.15mm, toptitle=1mm]

Sur une distribution fortement asymétrique --- typiquement les salaires
ou les prix de l'immobilier --- la moyenne peut être fortement tirée
vers le haut par quelques valeurs extrêmes. La \textbf{médiane} donne
alors une image plus représentative de la « valeur typique ».

\end{tcolorbox}

\subsection{1.2. Dispersion}\label{dispersion}

\begin{itemize}
\tightlist
\item
  l'\textbf{écart-type} (\(\sigma\)) et la \textbf{variance}
  (\(\sigma^2\)) mesurent l'écart moyen à la moyenne ;
\item
  l'\textbf{écart interquartile} (\(IQR = Q_3 - Q_1\)), déjà rencontré
  au chapitre précédent, mesure la dispersion de la moitié centrale des
  données, de façon robuste aux valeurs extrêmes ;
\item
  l'\textbf{étendue} (maximum − minimum) donne une idée rapide, mais
  très sensible aux valeurs aberrantes.
\end{itemize}

\begin{Shaded}
\begin{Highlighting}[]
\NormalTok{df[}\StringTok{"prix"}\NormalTok{].std()}
\NormalTok{df[}\StringTok{"prix"}\NormalTok{].var()}
\NormalTok{df[}\StringTok{"prix"}\NormalTok{].quantile([}\FloatTok{0.25}\NormalTok{, }\FloatTok{0.5}\NormalTok{, }\FloatTok{0.75}\NormalTok{])}
\end{Highlighting}
\end{Shaded}

Pour un petit tableau, on peut calculer les indicateurs à la main. Pour
les valeurs \(2, 4, 4, 6, 9\), la moyenne est \(25/5=5\) et la médiane
est \(4\). L'étendue vaut \(9-2=7\). Les quartiles dépendent de la
convention de calcul utilisée par l'outil ; dans un cours débutant,
l'important est de préciser la convention et de comparer les résultats
avec la même méthode.

L'écart-type est exprimé dans la même unité que la variable,
contrairement à la variance. Il est utile pour décrire la dispersion,
mais il est sensible aux valeurs extrêmes. L'IQR est souvent plus facile
à interpréter pour une distribution asymétrique. Le \textbf{coefficient
de variation}, défini par \(\sigma/\bar{x}\) lorsque la moyenne est
positive et non nulle, permet parfois de comparer la dispersion relative
de variables ayant des échelles différentes ; il est présenté ici comme
un repère et non comme un indicateur obligatoire.

\begin{tcolorbox}[enhanced jigsaw, arc=.35mm, bottomrule=.15mm, bottomtitle=1mm, breakable, colback=white, colbacktitle=quarto-callout-note-color!10!white, colframe=quarto-callout-note-color-frame, coltitle=black, left=2mm, leftrule=.75mm, opacityback=0, opacitybacktitle=0.6, rightrule=.15mm, title=\textcolor{quarto-callout-note-color}{\faInfo}\hspace{0.5em}{Priorité pour débuter}, titlerule=0mm, toprule=.15mm, toptitle=1mm]

Les indicateurs indispensables sont l'effectif, la moyenne, la médiane,
le minimum, le maximum, les quartiles et l'écart-type ou l'IQR.
L'asymétrie, l'aplatissement et le coefficient de variation sont des
compléments : il suffit d'en comprendre l'idée et de savoir quand
demander une explication supplémentaire.

\end{tcolorbox}

\subsection{1.3. Forme de la
distribution}\label{forme-de-la-distribution}

Deux indicateurs complètent la description d'une distribution :

\begin{itemize}
\tightlist
\item
  l'\textbf{asymétrie} (\emph{skewness}) : une valeur positive indique
  une distribution étalée vers la droite (queue de valeurs élevées), une
  valeur négative vers la gauche ;
\item
  l'\textbf{aplatissement} (\emph{kurtosis}) : indique si la
  distribution comporte davantage de valeurs extrêmes qu'une
  distribution normale.
\end{itemize}

\begin{Shaded}
\begin{Highlighting}[]
\NormalTok{df[}\StringTok{"prix"}\NormalTok{].skew()}
\NormalTok{df[}\StringTok{"prix"}\NormalTok{].kurt()}
\end{Highlighting}
\end{Shaded}

La méthode \texttt{describe()} de \texttt{pandas} regroupe l'essentiel
de ces indicateurs en une seule commande :

\begin{Shaded}
\begin{Highlighting}[]
\NormalTok{df.describe()}
\end{Highlighting}
\end{Shaded}

\section{2. Analyse univariée, bivariée et
multivariée}\label{analyse-univariuxe9e-bivariuxe9e-et-multivariuxe9e}

\begin{longtable}[]{@{}
  >{\raggedright\arraybackslash}p{(\linewidth - 4\tabcolsep) * \real{0.3333}}
  >{\raggedright\arraybackslash}p{(\linewidth - 4\tabcolsep) * \real{0.3333}}
  >{\raggedright\arraybackslash}p{(\linewidth - 4\tabcolsep) * \real{0.3333}}@{}}
\caption{Les trois niveaux de l'analyse
exploratoire}\label{tbl-niveaux-analyse}\tabularnewline
\toprule\noalign{}
\begin{minipage}[b]{\linewidth}\raggedright
Niveau d'analyse
\end{minipage} & \begin{minipage}[b]{\linewidth}\raggedright
Nombre de variables étudiées simultanément
\end{minipage} & \begin{minipage}[b]{\linewidth}\raggedright
Exemple
\end{minipage} \\
\midrule\noalign{}
\endfirsthead
\toprule\noalign{}
\begin{minipage}[b]{\linewidth}\raggedright
Niveau d'analyse
\end{minipage} & \begin{minipage}[b]{\linewidth}\raggedright
Nombre de variables étudiées simultanément
\end{minipage} & \begin{minipage}[b]{\linewidth}\raggedright
Exemple
\end{minipage} \\
\midrule\noalign{}
\endhead
\bottomrule\noalign{}
\endlastfoot
\textbf{Univariée} & Une seule & Distribution des prix de vente \\
\textbf{Bivariée} & Deux & Relation entre surface et prix \\
\textbf{Multivariée} & Trois ou plus & Relation entre surface, quartier
et prix \\
\end{longtable}

L'\textbf{analyse univariée} décrit une variable isolément (statistiques
descriptives, distribution). L'\textbf{analyse bivariée} cherche une
relation entre deux variables --- par exemple à l'aide d'un tableau
croisé pour deux variables catégorielles, ou d'un coefficient de
corrélation pour deux variables numériques. L'\textbf{analyse
multivariée} étend cette logique à plusieurs variables simultanément,
souvent à l'aide de tableaux d'agrégation ou de visualisations combinant
plusieurs dimensions (couleur, taille, facettes), présentées au chapitre
suivant.

\section{3. Corrélations et
associations}\label{corruxe9lations-et-associations}

\subsection{3.1. Variables numériques : Pearson et
Spearman}\label{variables-numuxe9riques-pearson-et-spearman}

\begin{itemize}
\tightlist
\item
  le \textbf{coefficient de corrélation de Pearson} (\(r\), entre −1 et
  1) mesure l'intensité d'une relation \textbf{linéaire} entre deux
  variables numériques ;
\item
  le \textbf{coefficient de corrélation de Spearman} mesure une relation
  \textbf{monotone} (pas nécessairement linéaire), en travaillant sur
  les rangs des valeurs plutôt que sur les valeurs elles-mêmes --- plus
  robuste aux valeurs aberrantes.
\end{itemize}

\begin{Shaded}
\begin{Highlighting}[]
\NormalTok{df[[}\StringTok{"surface"}\NormalTok{, }\StringTok{"prix"}\NormalTok{]].corr(method}\OperatorTok{=}\StringTok{"pearson"}\NormalTok{)}
\NormalTok{df[[}\StringTok{"surface"}\NormalTok{, }\StringTok{"prix"}\NormalTok{]].corr(method}\OperatorTok{=}\StringTok{"spearman"}\NormalTok{)}
\end{Highlighting}
\end{Shaded}

\begin{tcolorbox}[enhanced jigsaw, arc=.35mm, bottomrule=.15mm, bottomtitle=1mm, breakable, colback=white, colbacktitle=quarto-callout-warning-color!10!white, colframe=quarto-callout-warning-color-frame, coltitle=black, left=2mm, leftrule=.75mm, opacityback=0, opacitybacktitle=0.6, rightrule=.15mm, title=\textcolor{quarto-callout-warning-color}{\faExclamationTriangle}\hspace{0.5em}{Corrélation n'est pas causalité}, titlerule=0mm, toprule=.15mm, toptitle=1mm]

Une corrélation élevée entre deux variables ne signifie pas que l'une
provoque l'autre : elle peut résulter d'une \textbf{variable cachée}
commune, ou d'une pure coïncidence statistique. Un exemple classique est
la corrélation observée entre les ventes de glaces et le nombre de
noyades : les deux augmentent en été, sans lien de cause à effet direct
entre elles.

\end{tcolorbox}

\subsection{3.2. Variables catégorielles : tableaux croisés et V de
Cramér}\label{variables-catuxe9gorielles-tableaux-croisuxe9s-et-v-de-cramuxe9r}

Pour deux variables catégorielles, on utilise un \textbf{tableau croisé}
(\emph{crosstab}), éventuellement complété par le \textbf{V de Cramér},
qui résume en un seul indicateur (entre 0 et 1) l'intensité de
l'association :

\begin{Shaded}
\begin{Highlighting}[]
\NormalTok{tableau\_croise }\OperatorTok{=}\NormalTok{ pd.crosstab(df[}\StringTok{"quartier"}\NormalTok{], df[}\StringTok{"type\_bien"}\NormalTok{])}
\end{Highlighting}
\end{Shaded}

\section{\texorpdfstring{4. Agrégation : \texttt{groupby} et tableaux
croisés
dynamiques}{4. Agrégation : groupby et tableaux croisés dynamiques}}\label{agruxe9gation-groupby-et-tableaux-croisuxe9s-dynamiques}

L'outil central de l'exploration multivariée avec \texttt{pandas} est la
méthode \texttt{groupby()}, qui permet de calculer un indicateur pour
chaque sous-groupe d'une variable catégorielle :

\begin{Shaded}
\begin{Highlighting}[]
\NormalTok{df.groupby(}\StringTok{"quartier"}\NormalTok{)[}\StringTok{"prix"}\NormalTok{].mean().sort\_values(ascending}\OperatorTok{=}\VariableTok{False}\NormalTok{)}

\NormalTok{df.groupby([}\StringTok{"quartier"}\NormalTok{, }\StringTok{"type\_bien"}\NormalTok{]).agg(}
\NormalTok{    prix\_moyen}\OperatorTok{=}\NormalTok{(}\StringTok{"prix"}\NormalTok{, }\StringTok{"mean"}\NormalTok{),}
\NormalTok{    nombre\_biens}\OperatorTok{=}\NormalTok{(}\StringTok{"prix"}\NormalTok{, }\StringTok{"count"}\NormalTok{)}
\NormalTok{)}
\end{Highlighting}
\end{Shaded}

Le \textbf{tableau croisé dynamique} (\emph{pivot table}) offre une
présentation plus lisible lorsque l'on croise deux dimensions
catégorielles avec une mesure numérique :

\begin{Shaded}
\begin{Highlighting}[]
\NormalTok{pd.pivot\_table(}
\NormalTok{    df,}
\NormalTok{    values}\OperatorTok{=}\StringTok{"prix"}\NormalTok{,}
\NormalTok{    index}\OperatorTok{=}\StringTok{"quartier"}\NormalTok{,}
\NormalTok{    columns}\OperatorTok{=}\StringTok{"type\_bien"}\NormalTok{,}
\NormalTok{    aggfunc}\OperatorTok{=}\StringTok{"mean"}
\NormalTok{)}
\end{Highlighting}
\end{Shaded}

Ces deux outils constituent la base de la quasi-totalité des tableaux de
bord et des rapports produits par un Data Analyst.

\section{5. Choisir une visualisation
simple}\label{choisir-une-visualisation-simple}

Les statistiques répondent à des questions numériques ; les graphiques
rendent les formes et les comparaisons plus faciles à voir. Pour
débuter, on peut retenir les associations suivantes :

{\def\LTcaptype{none} % do not increment counter
\begin{longtable}[]{@{}
  >{\raggedright\arraybackslash}p{(\linewidth - 4\tabcolsep) * \real{0.3333}}
  >{\raggedright\arraybackslash}p{(\linewidth - 4\tabcolsep) * \real{0.3333}}
  >{\raggedright\arraybackslash}p{(\linewidth - 4\tabcolsep) * \real{0.3333}}@{}}
\toprule\noalign{}
\begin{minipage}[b]{\linewidth}\raggedright
Question
\end{minipage} & \begin{minipage}[b]{\linewidth}\raggedright
Graphique conseillé
\end{minipage} & \begin{minipage}[b]{\linewidth}\raggedright
Ce que l'on observe
\end{minipage} \\
\midrule\noalign{}
\endhead
\bottomrule\noalign{}
\endlastfoot
Quelle est la répartition d'une variable numérique ? & Histogramme &
Forme, concentration et asymétrie \\
Y a-t-il des valeurs atypiques ? & Boîte à moustaches & Médiane,
quartiles et valeurs éloignées \\
Comment une valeur évolue-t-elle dans le temps ? & Courbe & Tendance,
saisonnalité et ruptures \\
Comment comparer des catégories ? & Barres & Effectifs, sommes ou
moyennes par catégorie \\
Deux variables numériques évoluent-elles ensemble ? & Nuage de points &
Relation, groupes et valeurs atypiques \\
\end{longtable}
}

Un graphique ne doit pas seulement être « joli ». Son titre doit
préciser la mesure, l'unité et la période ; ses axes doivent être
lisibles ; la source et les filtres importants doivent être connus. Pour
comparer des catégories, l'effectif et la moyenne ne racontent pas la
même histoire : une catégorie peut avoir une moyenne élevée avec très
peu d'observations.

\section{6. Profilage automatique des
données}\label{profilage-automatique-des-donnuxe9es}

Sur un nouveau jeu de données, il est souvent utile d'obtenir rapidement
une vue d'ensemble avant de se lancer dans une exploration manuelle
détaillée.

\begin{tcolorbox}[enhanced jigsaw, arc=.35mm, bottomrule=.15mm, bottomtitle=1mm, breakable, colback=white, colbacktitle=quarto-callout-important-color!10!white, colframe=quarto-callout-important-color-frame, coltitle=black, left=2mm, leftrule=.75mm, opacityback=0, opacitybacktitle=0.6, rightrule=.15mm, title=\textcolor{quarto-callout-important-color}{\faExclamation}\hspace{0.5em}{Zoom sur les tendances : rapports d'exploration automatisés}, titlerule=0mm, toprule=.15mm, toptitle=1mm]

Des bibliothèques comme
\textbf{\href{https://github.com/ydataai/ydata-profiling}{\texttt{ydata-profiling}}}
(anciennement \texttt{pandas-profiling}) ou
\textbf{\href{https://github.com/fbdesignpro/sweetviz}{\texttt{Sweetviz}}}
génèrent, en une seule instruction, un rapport HTML complet :
statistiques descriptives par variable, distributions, valeurs
manquantes, corrélations entre variables, alertes de qualité (variables
constantes, fortement corrélées, etc.).

\begin{Shaded}
\begin{Highlighting}[]
\ImportTok{from}\NormalTok{ ydata\_profiling }\ImportTok{import}\NormalTok{ ProfileReport}

\NormalTok{profil }\OperatorTok{=}\NormalTok{ ProfileReport(df, title}\OperatorTok{=}\StringTok{"Exploration du jeu de données"}\NormalTok{)}
\NormalTok{profil.to\_file(}\StringTok{"rapport\_exploration.html"}\NormalTok{)}
\end{Highlighting}
\end{Shaded}

Ces outils ne remplacent pas le raisonnement de l'analyste --- ils
accélèrent la première prise de contact avec un jeu de données et aident
à repérer rapidement les points nécessitant une attention particulière,
avant une exploration ciblée avec les techniques présentées dans ce
chapitre.

\end{tcolorbox}

\begin{tcolorbox}[enhanced jigsaw, arc=.35mm, bottomrule=.15mm, bottomtitle=1mm, breakable, colback=white, colbacktitle=quarto-callout-note-color!10!white, colframe=quarto-callout-note-color-frame, coltitle=black, left=2mm, leftrule=.75mm, opacityback=0, opacitybacktitle=0.6, rightrule=.15mm, title=\textcolor{quarto-callout-note-color}{\faInfo}\hspace{0.5em}{Auto-évaluation}, titlerule=0mm, toprule=.15mm, toptitle=1mm]

Avant de passer au chapitre suivant, assurez-vous de pouvoir répondre «
oui » aux questions suivantes :

\begin{itemize}
\tightlist
\item[$\square$]
  Je sais expliquer pourquoi la médiane est parfois préférable à la
  moyenne.
\item[$\square$]
  Je sais distinguer une analyse univariée, bivariée et multivariée avec
  un exemple pour chacune.
\item[$\square$]
  Je sais expliquer pourquoi une forte corrélation n'implique pas une
  relation de cause à effet.
\item[$\square$]
  Je sais construire un tableau croisé dynamique avec \texttt{pandas}.
\end{itemize}

\end{tcolorbox}

\section{Exercices}\label{exercices-2}

\begin{exercise}[Interpréter des statistiques
descriptives]\protect\hypertarget{exr-eda-1}{}\label{exr-eda-1}

Pour une variable \texttt{revenu\_mensuel} (en MAD), on obtient :
moyenne = 8 500, médiane = 6 200, écart-type = 5 400, asymétrie
(\emph{skewness}) = 2,3.

\begin{enumerate}
\def\labelenumi{\alph{enumi}.}
\tightlist
\item
  Que suggère l'écart important entre la moyenne et la médiane ?
\item
  Que confirme la valeur positive et élevée de l'asymétrie ?
\end{enumerate}

\end{exercise}

\begin{tcolorbox}[enhanced jigsaw, arc=.35mm, bottomrule=.15mm, bottomtitle=1mm, breakable, colback=white, colbacktitle=quarto-callout-tip-color!10!white, colframe=quarto-callout-tip-color-frame, coltitle=black, left=2mm, leftrule=.75mm, opacityback=0, opacitybacktitle=0.6, rightrule=.15mm, title=\textcolor{quarto-callout-tip-color}{\faLightbulb}\hspace{0.5em}{Solution de l'exercice Exercice~\ref{exr-eda-1}}, titlerule=0mm, toprule=.15mm, toptitle=1mm]

\begin{enumerate}
\def\labelenumi{\alph{enumi}.}
\tightlist
\item
  Un écart important entre une moyenne nettement supérieure à la médiane
  suggère la présence de quelques revenus très élevés qui tirent la
  moyenne vers le haut, sans être représentatifs de la majorité des
  observations.
\item
  Une asymétrie positive et élevée confirme une distribution étalée vers
  la droite (une « longue traîne » de hauts revenus), cohérente avec
  l'écart observé entre moyenne et médiane.
\end{enumerate}

\end{tcolorbox}

\begin{exercise}[Choisir le bon coefficient de
corrélation]\protect\hypertarget{exr-eda-2}{}\label{exr-eda-2}

Pour chacune des situations suivantes, indiquez s'il est préférable
d'utiliser Pearson ou Spearman :

\begin{enumerate}
\def\labelenumi{\alph{enumi}.}
\tightlist
\item
  La relation entre l'ancienneté d'un employé (en années) et son
  salaire, globalement croissante mais pas forcément linéaire.
\item
  La relation entre la taille et le poids d'un ensemble d'individus
  adultes, approximativement linéaire.
\end{enumerate}

\end{exercise}

\begin{tcolorbox}[enhanced jigsaw, arc=.35mm, bottomrule=.15mm, bottomtitle=1mm, breakable, colback=white, colbacktitle=quarto-callout-tip-color!10!white, colframe=quarto-callout-tip-color-frame, coltitle=black, left=2mm, leftrule=.75mm, opacityback=0, opacitybacktitle=0.6, rightrule=.15mm, title=\textcolor{quarto-callout-tip-color}{\faLightbulb}\hspace{0.5em}{Solution de l'exercice Exercice~\ref{exr-eda-2}}, titlerule=0mm, toprule=.15mm, toptitle=1mm]

\begin{enumerate}
\def\labelenumi{\alph{enumi}.}
\tightlist
\item
  \textbf{Spearman}, car la relation est monotone (croissante) sans être
  nécessairement linéaire ; Spearman capte ce type de relation même s'il
  n'est pas de forme parfaitement linéaire.
\item
  \textbf{Pearson}, adapté à une relation linéaire entre deux variables
  numériques.
\end{enumerate}

\end{tcolorbox}

\begin{exercise}[Construire un plan
d'exploration]\protect\hypertarget{exr-eda-3}{}\label{exr-eda-3}

On dispose d'un jeu de données de transactions d'un site e-commerce
marocain, avec les colonnes suivantes : \texttt{date\_commande},
\texttt{categorie\_produit}, \texttt{prix}, \texttt{region\_client},
\texttt{moyen\_paiement}.

Proposez, pour ce jeu de données, un exemple concret d'analyse
univariée, un exemple d'analyse bivariée et un exemple d'analyse
multivariée, en vous appuyant sur ces colonnes.

\end{exercise}

\begin{tcolorbox}[enhanced jigsaw, arc=.35mm, bottomrule=.15mm, bottomtitle=1mm, breakable, colback=white, colbacktitle=quarto-callout-tip-color!10!white, colframe=quarto-callout-tip-color-frame, coltitle=black, left=2mm, leftrule=.75mm, opacityback=0, opacitybacktitle=0.6, rightrule=.15mm, title=\textcolor{quarto-callout-tip-color}{\faLightbulb}\hspace{0.5em}{Solution de l'exercice Exercice~\ref{exr-eda-3}}, titlerule=0mm, toprule=.15mm, toptitle=1mm]

\begin{itemize}
\tightlist
\item
  \textbf{Univariée} : distribution du \texttt{prix} des commandes
  (moyenne, médiane, dispersion).
\item
  \textbf{Bivariée} : prix moyen selon la \texttt{categorie\_produit}
  (via \texttt{groupby}).
\item
  \textbf{Multivariée} : prix moyen croisé par \texttt{region\_client}
  \textbf{et} \texttt{moyen\_paiement} (via un tableau croisé
  dynamique), permettant de voir si certaines régions privilégient
  certains moyens de paiement pour des achats plus ou moins coûteux.
\end{itemize}

\end{tcolorbox}

\begin{exercise}[Calculer des indicateurs
simples]\protect\hypertarget{exr-eda-4}{}\label{exr-eda-4}

Les durées de livraison, en jours, sont : 2, 3, 3, 4, 8.

\begin{enumerate}
\def\labelenumi{\alph{enumi}.}
\tightlist
\item
  Calculez la moyenne, la médiane et l'étendue.
\item
  Quel indicateur est le moins influencé par la durée de 8 jours ?
\item
  Quelle question supplémentaire poseriez-vous avant de conclure que le
  service est lent ?
\end{enumerate}

\end{exercise}

\begin{tcolorbox}[enhanced jigsaw, arc=.35mm, bottomrule=.15mm, bottomtitle=1mm, breakable, colback=white, colbacktitle=quarto-callout-tip-color!10!white, colframe=quarto-callout-tip-color-frame, coltitle=black, left=2mm, leftrule=.75mm, opacityback=0, opacitybacktitle=0.6, rightrule=.15mm, title=\textcolor{quarto-callout-tip-color}{\faLightbulb}\hspace{0.5em}{Solution de l'exercice Exercice~\ref{exr-eda-4}}, titlerule=0mm, toprule=.15mm, toptitle=1mm]

\begin{enumerate}
\def\labelenumi{\alph{enumi}.}
\tightlist
\item
  Moyenne \(=(2+3+3+4+8)/5=\textbf{4}\) jours ; médiane = \textbf{3
  jours} ; étendue \(=8-2=\textbf{6}\) jours.
\item
  La \textbf{médiane}.
\item
  Il faut connaître le nombre total de livraisons, la période, la région
  et la proportion de livraisons de 8 jours : une seule valeur extrême
  ne suffit pas.
\end{enumerate}

\end{tcolorbox}

\begin{exercise}[Lire une
dispersion]\protect\hypertarget{exr-eda-5}{}\label{exr-eda-5}

Deux classes ont la même moyenne à un examen : 12/20. La classe A a un
écart-type de 1,5 et la classe B un écart-type de 4.

\begin{enumerate}
\def\labelenumi{\alph{enumi}.}
\tightlist
\item
  Quelle classe est la plus homogène ?
\item
  Que signifie l'écart-type plus élevé de la classe B ?
\item
  Peut-on dire que les étudiants de B ont de moins bonnes notes ?
\end{enumerate}

\end{exercise}

\begin{tcolorbox}[enhanced jigsaw, arc=.35mm, bottomrule=.15mm, bottomtitle=1mm, breakable, colback=white, colbacktitle=quarto-callout-tip-color!10!white, colframe=quarto-callout-tip-color-frame, coltitle=black, left=2mm, leftrule=.75mm, opacityback=0, opacitybacktitle=0.6, rightrule=.15mm, title=\textcolor{quarto-callout-tip-color}{\faLightbulb}\hspace{0.5em}{Solution de l'exercice Exercice~\ref{exr-eda-5}}, titlerule=0mm, toprule=.15mm, toptitle=1mm]

\begin{enumerate}
\def\labelenumi{\alph{enumi}.}
\tightlist
\item
  La classe \textbf{A} est plus homogène.
\item
  Les notes de B sont plus éloignées de leur moyenne : les niveaux sont
  davantage dispersés.
\item
  Non. Les deux moyennes sont identiques ; B est seulement plus
  hétérogène.
\end{enumerate}

\end{tcolorbox}

\begin{exercise}[Choisir le bon
graphique]\protect\hypertarget{exr-eda-6}{}\label{exr-eda-6}

Associez chaque besoin à un graphique : histogramme, courbe, barres,
boîte à moustaches ou nuage de points.

\begin{enumerate}
\def\labelenumi{\alph{enumi}.}
\tightlist
\item
  Comparer le chiffre d'affaires de quatre régions.
\item
  Observer le nombre de commandes chaque mois.
\item
  Repérer des valeurs extrêmes dans les salaires.
\item
  Étudier le lien entre surface et prix d'un logement.
\end{enumerate}

\end{exercise}

\begin{tcolorbox}[enhanced jigsaw, arc=.35mm, bottomrule=.15mm, bottomtitle=1mm, breakable, colback=white, colbacktitle=quarto-callout-tip-color!10!white, colframe=quarto-callout-tip-color-frame, coltitle=black, left=2mm, leftrule=.75mm, opacityback=0, opacitybacktitle=0.6, rightrule=.15mm, title=\textcolor{quarto-callout-tip-color}{\faLightbulb}\hspace{0.5em}{Solution de l'exercice Exercice~\ref{exr-eda-6}}, titlerule=0mm, toprule=.15mm, toptitle=1mm]

\begin{enumerate}
\def\labelenumi{\alph{enumi}.}
\tightlist
\item
  \textbf{Barres} ; b. \textbf{Courbe} ; c.~\textbf{Boîte à moustaches}
  ; d.~\textbf{Nuage de points}. Un histogramme serait adapté pour la
  répartition d'une seule variable numérique.
\end{enumerate}

\end{tcolorbox}

\begin{exercise}[Corrélation et
causalité]\protect\hypertarget{exr-eda-7}{}\label{exr-eda-7}

Dans une ville, le nombre de glaces vendues et le nombre de baignades
augmentent tous les deux en été.

\begin{enumerate}
\def\labelenumi{\alph{enumi}.}
\tightlist
\item
  Quelle relation statistique peut-on observer ?
\item
  Pourquoi ne peut-on pas conclure que les glaces provoquent les
  baignades ?
\item
  Quelle variable cachée peut expliquer les deux évolutions ?
\end{enumerate}

\end{exercise}

\begin{tcolorbox}[enhanced jigsaw, arc=.35mm, bottomrule=.15mm, bottomtitle=1mm, breakable, colback=white, colbacktitle=quarto-callout-tip-color!10!white, colframe=quarto-callout-tip-color-frame, coltitle=black, left=2mm, leftrule=.75mm, opacityback=0, opacitybacktitle=0.6, rightrule=.15mm, title=\textcolor{quarto-callout-tip-color}{\faLightbulb}\hspace{0.5em}{Solution de l'exercice Exercice~\ref{exr-eda-7}}, titlerule=0mm, toprule=.15mm, toptitle=1mm]

\begin{enumerate}
\def\labelenumi{\alph{enumi}.}
\tightlist
\item
  Une corrélation positive.
\item
  Une corrélation décrit une association, pas une cause.
\item
  La \textbf{température} ou la saison : le beau temps favorise à la
  fois l'achat de glaces et les baignades.
\end{enumerate}

\end{tcolorbox}

\begin{exercise}[Interpréter un tableau
agrégé]\protect\hypertarget{exr-eda-8}{}\label{exr-eda-8}

Un tableau par région présente les résultats suivants :

{\def\LTcaptype{none} % do not increment counter
\begin{longtable}[]{@{}lrr@{}}
\toprule\noalign{}
Région & Nombre de commandes & Montant moyen (MAD) \\
\midrule\noalign{}
\endhead
\bottomrule\noalign{}
\endlastfoot
Nord & 100 & 180 \\
Centre & 20 & 350 \\
Sud & 80 & 200 \\
\end{longtable}
}

\begin{enumerate}
\def\labelenumi{\alph{enumi}.}
\tightlist
\item
  Quelle région a le montant moyen le plus élevé ?
\item
  Pourquoi faut-il également regarder le nombre de commandes ?
\item
  Calculez le montant total estimé pour chaque région et identifiez la
  région au total le plus élevé.
\end{enumerate}

\end{exercise}

\begin{tcolorbox}[enhanced jigsaw, arc=.35mm, bottomrule=.15mm, bottomtitle=1mm, breakable, colback=white, colbacktitle=quarto-callout-tip-color!10!white, colframe=quarto-callout-tip-color-frame, coltitle=black, left=2mm, leftrule=.75mm, opacityback=0, opacitybacktitle=0.6, rightrule=.15mm, title=\textcolor{quarto-callout-tip-color}{\faLightbulb}\hspace{0.5em}{Solution de l'exercice Exercice~\ref{exr-eda-8}}, titlerule=0mm, toprule=.15mm, toptitle=1mm]

\begin{enumerate}
\def\labelenumi{\alph{enumi}.}
\tightlist
\item
  Le \textbf{Centre} (350 MAD).
\item
  Sa moyenne repose sur seulement 20 commandes ; elle est donc moins
  représentative d'un volume important que celle d'une région avec
  beaucoup de commandes.
\item
  Nord : \(100\times180=\textbf{18\,000}\) MAD ; Centre :
  \(20\times350=\textbf{7\,000}\) MAD ; Sud :
  \(80\times200=\textbf{16\,000}\) MAD. Le total le plus élevé est celui
  du \textbf{Nord}.
\end{enumerate}

\end{tcolorbox}

\bookmarksetup{startatroot}

\chapter{Introduction à l'entrepôt de
données}\label{introduction-uxe0-lentrepuxf4t-de-donnuxe9es}

\begin{tcolorbox}[enhanced jigsaw, arc=.35mm, bottomrule=.15mm, bottomtitle=1mm, breakable, colback=white, colbacktitle=quarto-callout-tip-color!10!white, colframe=quarto-callout-tip-color-frame, coltitle=black, left=2mm, leftrule=.75mm, opacityback=0, opacitybacktitle=0.6, rightrule=.15mm, title=\textcolor{quarto-callout-tip-color}{\faLightbulb}\hspace{0.5em}{Objectifs du chapitre}, titlerule=0mm, toprule=.15mm, toptitle=1mm]

À la fin de ce chapitre, l'étudiant sera capable de :

\begin{itemize}
\tightlist
\item
  \textbf{Expliquer} pourquoi une organisation sépare souvent les
  systèmes opérationnels des données destinées à l'analyse ;
\item
  \textbf{Distinguer} une base transactionnelle, un entrepôt de données,
  un lac de données et un datamart ;
\item
  \textbf{Lire} un modèle en étoile composé d'une table de faits et de
  tables de dimensions ;
\item
  \textbf{Définir} le grain d'une table de faits et choisir des mesures
  adaptées ;
\item
  \textbf{Décrire} les grandes étapes d'un flux ETL/ELT, de la source au
  tableau de bord.
\end{itemize}

\end{tcolorbox}

Un tableau de bord fiable ne commence pas par un graphique. Il commence
par des données organisées, documentées et mises à jour de façon
régulière. L'\textbf{entrepôt de données} (\emph{data warehouse})
fournit cette base pour l'analyse.

\section{1. Pourquoi un entrepôt de données
?}\label{pourquoi-un-entrepuxf4t-de-donnuxe9es}

\subsection{1.1. Les systèmes
opérationnels}\label{les-systuxe8mes-opuxe9rationnels}

Une application de vente, une caisse ou un site web utilise une base
\textbf{opérationnelle} pour enregistrer les événements du quotidien :
créer une commande, modifier une adresse, enregistrer un paiement. Ces
systèmes doivent répondre vite à de petites opérations et éviter les
incohérences.

Une base opérationnelle est souvent organisée pour limiter la répétition
des informations. Une commande peut être répartie entre
\texttt{clients}, \texttt{commandes} et \texttt{lignes\_commande},
reliées par des identifiants.

Cette organisation est adaptée à l'écriture, mais une question
analytique comme « Quel est le chiffre d'affaires mensuel par région et
par catégorie depuis trois ans ? » peut nécessiter de nombreuses
jointures et plusieurs sources. Exécuter ce type de requête directement
sur la base de production peut ralentir l'application.

\subsection{1.2. La base analytique}\label{la-base-analytique}

Un \textbf{entrepôt de données} est un espace centralisé, organisé pour
analyser l'historique de l'organisation. Il rassemble des données
provenant de plusieurs sources, conserve leur contexte et les transforme
dans un format plus simple à interroger.

{\def\LTcaptype{none} % do not increment counter
\begin{longtable}[]{@{}
  >{\raggedright\arraybackslash}p{(\linewidth - 4\tabcolsep) * \real{0.3333}}
  >{\raggedright\arraybackslash}p{(\linewidth - 4\tabcolsep) * \real{0.3333}}
  >{\raggedright\arraybackslash}p{(\linewidth - 4\tabcolsep) * \real{0.3333}}@{}}
\toprule\noalign{}
\begin{minipage}[b]{\linewidth}\raggedright
Système
\end{minipage} & \begin{minipage}[b]{\linewidth}\raggedright
But principal
\end{minipage} & \begin{minipage}[b]{\linewidth}\raggedright
Exemple de question
\end{minipage} \\
\midrule\noalign{}
\endhead
\bottomrule\noalign{}
\endlastfoot
Base opérationnelle (OLTP) & Enregistrer les événements courants & «
Quel est le statut de la commande 1524 ? » \\
Entrepôt analytique (OLAP) & Comparer, agréger et suivre l'historique &
« Comment les ventes ont-elles évolué par région ? » \\
\end{longtable}
}

Les deux systèmes sont complémentaires. L'entrepôt n'a pas pour objectif
de remplacer l'application métier.

\subsection{1.3. Entrepôt, lac et
datamart}\label{entrepuxf4t-lac-et-datamart}

\begin{itemize}
\tightlist
\item
  Un \textbf{data warehouse} contient des données structurées, préparées
  et gouvernées pour l'analyse.
\item
  Un \textbf{data lake} peut stocker des données brutes de formats
  variés (fichiers, textes, événements) avant leur usage ; il demande
  des règles claires pour rester compréhensible.
\item
  Un \textbf{datamart} est une partie orientée vers un besoin ou un
  service, par exemple un datamart des ventes ou des ressources
  humaines.
\end{itemize}

\begin{tcolorbox}[enhanced jigsaw, arc=.35mm, bottomrule=.15mm, bottomtitle=1mm, breakable, colback=white, colbacktitle=quarto-callout-warning-color!10!white, colframe=quarto-callout-warning-color-frame, coltitle=black, left=2mm, leftrule=.75mm, opacityback=0, opacitybacktitle=0.6, rightrule=.15mm, title=\textcolor{quarto-callout-warning-color}{\faExclamationTriangle}\hspace{0.5em}{Un entrepôt ne corrige pas une mauvaise question}, titlerule=0mm, toprule=.15mm, toptitle=1mm]

Centraliser beaucoup de données ne suffit pas. Il faut définir le
vocabulaire métier, les unités, la fréquence de mise à jour et les
indicateurs attendus. Sinon, on obtient plusieurs chiffres différents
pour un même « chiffre d'affaires ».

\end{tcolorbox}

\section{2. Le modèle en étoile}\label{le-moduxe8le-en-uxe9toile}

Le modèle en étoile est une organisation pédagogique et professionnelle
très fréquente pour les données analytiques. Il place une \textbf{table
de faits} au centre et des \textbf{tables de dimensions} autour.

\subsection{2.1. La table de faits}\label{la-table-de-faits}

Une table de faits contient les événements mesurables : ventes,
livraisons, appels ou visites. Elle comprend généralement :

\begin{itemize}
\tightlist
\item
  des clés vers les dimensions (\texttt{client\_id},
  \texttt{produit\_id}, \texttt{date\_id}, \texttt{magasin\_id}) ;
\item
  des mesures numériques (\texttt{quantite}, \texttt{montant},
  \texttt{remise}) ;
\item
  une ligne par événement selon un niveau de détail défini.
\end{itemize}

Le \textbf{grain} est la phrase qui précise ce que représente une ligne.
Par exemple : « une ligne représente un produit d'une commande, dans un
magasin, à une date donnée ». Le grain doit être défini avant de choisir
les mesures.

\subsection{2.2. Les tables de
dimensions}\label{les-tables-de-dimensions}

Une dimension décrit le contexte d'un fait :

\begin{itemize}
\tightlist
\item
  \texttt{dim\_date} : jour, mois, trimestre, année ;
\item
  \texttt{dim\_produit} : produit, catégorie, marque ;
\item
  \texttt{dim\_client} : ville, région, segment ;
\item
  \texttt{dim\_magasin} : magasin, ville, zone.
\end{itemize}

Les dimensions contiennent les attributs utilisés pour filtrer,
regrouper et expliquer les mesures. Elles sont souvent dénormalisées
pour rendre l'exploration plus simple.

\begin{Shaded}
\begin{Highlighting}[]
\NormalTok{                         dim\_date}
\NormalTok{                            |}
\NormalTok{dim\_client {-}{-}{-}{-} fact\_ventes {-}{-}{-}{-} dim\_produit}
\NormalTok{                            |}
\NormalTok{                       dim\_magasin}
\end{Highlighting}
\end{Shaded}

\subsection{2.3. Exemple de modèle de
ventes}\label{exemple-de-moduxe8le-de-ventes}

\begin{Shaded}
\begin{Highlighting}[]
\NormalTok{fact\_ventes}
\NormalTok{{-}{-}{-}{-}{-}{-}{-}{-}{-}{-}{-}}
\NormalTok{date\_id}
\NormalTok{produit\_id}
\NormalTok{client\_id}
\NormalTok{magasin\_id}
\NormalTok{quantite}
\NormalTok{prix\_unitaire}
\NormalTok{remise}
\NormalTok{montant\_net}

\NormalTok{dim\_produit}
\NormalTok{{-}{-}{-}{-}{-}{-}{-}{-}{-}{-}{-}}
\NormalTok{produit\_id}
\NormalTok{libelle}
\NormalTok{categorie}
\NormalTok{marque}

\NormalTok{dim\_date}
\NormalTok{{-}{-}{-}{-}{-}{-}{-}{-}}
\NormalTok{date\_id}
\NormalTok{date\_complete}
\NormalTok{mois}
\NormalTok{trimestre}
\NormalTok{annee}
\end{Highlighting}
\end{Shaded}

Le montant net peut être calculé par :

\[
\text{montant\_net}
=\text{quantité}\times\text{prix unitaire}-\text{remise}.
\]

Un modèle en étoile ne dicte pas un seul outil : il peut être utilisé
dans une base SQL, un fichier analytique ou un modèle Power BI.

\begin{tcolorbox}[enhanced jigsaw, arc=.35mm, bottomrule=.15mm, bottomtitle=1mm, breakable, colback=white, colbacktitle=quarto-callout-note-color!10!white, colframe=quarto-callout-note-color-frame, coltitle=black, left=2mm, leftrule=.75mm, opacityback=0, opacitybacktitle=0.6, rightrule=.15mm, title=\textcolor{quarto-callout-note-color}{\faInfo}\hspace{0.5em}{Le rôle de la dimension Date}, titlerule=0mm, toprule=.15mm, toptitle=1mm]

Même si la source contient déjà une date, une dimension Date facilite
les regroupements par mois, trimestre, semaine ou jour ouvré. Elle
permet aussi de définir clairement les périodes comparées.

\end{tcolorbox}

\section{3. Les mesures et les
indicateurs}\label{les-mesures-et-les-indicateurs}

Une \textbf{mesure additive} peut être additionnée sur toutes les
dimensions, comme le nombre d'articles ou le montant des ventes. Une
mesure \textbf{semi-additive}, comme un stock de fin de journée, peut
être additionnée par produit mais pas dans le temps. Une mesure
\textbf{non additive}, comme un taux ou une moyenne, doit être
recalculée à partir de ses éléments plutôt que sommée.

Exemple : si une région a 10 commandes et une autre 2 commandes, la
moyenne globale ne doit pas être calculée comme la moyenne des deux
moyennes si les effectifs sont différents. Il faut utiliser :

\[
\text{moyenne globale} =
\frac{\text{total des montants}}{\text{nombre total de commandes}}.
\]

Les noms des indicateurs doivent préciser leur définition. « Ventes »
peut signifier chiffre d'affaires brut, chiffre d'affaires net,
commandes payées ou unités vendues. Dans un dictionnaire métier, on
documente la formule, l'unité, la période et les filtres.

\section{4. Le flux ETL/ELT}\label{le-flux-etlelt}

Les données passent généralement par un flux en plusieurs étapes :

\begin{enumerate}
\def\labelenumi{\arabic{enumi}.}
\tightlist
\item
  \textbf{Extract (extraire)} : lire les sources opérationnelles,
  fichiers ou API ;
\item
  \textbf{Transform (transformer)} : convertir les dates, harmoniser les
  unités, dédupliquer, contrôler les valeurs et calculer les champs
  dérivés ;
\item
  \textbf{Load (charger)} : déposer les données préparées dans les
  tables de l'entrepôt.
\end{enumerate}

Dans une approche \textbf{ELT}, on charge d'abord les données dans une
plateforme capable de les traiter, puis on réalise les transformations à
l'intérieur de cette plateforme. Les deux approches poursuivent le même
but : rendre les données fiables et exploitables.

Un chargement quotidien doit pouvoir répondre à quatre questions simples
: quelles sources ont été lues ? combien de lignes ont été reçues ?
combien ont été rejetées ? quand le chargement s'est-il terminé ?

\begin{tcolorbox}[enhanced jigsaw, arc=.35mm, bottomrule=.15mm, bottomtitle=1mm, breakable, colback=white, colbacktitle=quarto-callout-important-color!10!white, colframe=quarto-callout-important-color-frame, coltitle=black, left=2mm, leftrule=.75mm, opacityback=0, opacitybacktitle=0.6, rightrule=.15mm, title=\textcolor{quarto-callout-important-color}{\faExclamation}\hspace{0.5em}{Qualité et traçabilité}, titlerule=0mm, toprule=.15mm, toptitle=1mm]

Conservez les données sources ou un journal de chargement, le nombre de
lignes avant et après transformation, les erreurs détectées et la date
de mise à jour. Ces informations permettent de comprendre pourquoi un
indicateur a changé.

\end{tcolorbox}

\section{5. Gouvernance et sécurité, sans complexité
inutile}\label{gouvernance-et-suxe9curituxe9-sans-complexituxe9-inutile}

La \textbf{gouvernance} désigne les règles qui rendent les données
compréhensibles, accessibles aux bonnes personnes et utilisées de
manière responsable. Pour commencer, un projet peut mettre en place :

\begin{itemize}
\tightlist
\item
  un dictionnaire des tables et des indicateurs ;
\item
  un responsable pour chaque donnée importante ;
\item
  des règles de qualité (unicité d'un identifiant, prix positif, date
  obligatoire) ;
\item
  des droits d'accès adaptés, notamment pour les données personnelles ;
\item
  une fréquence de mise à jour annoncée.
\end{itemize}

Un tableau de bord ne devrait pas exposer inutilement des noms, numéros
de téléphone ou adresses. On peut agréger les résultats, masquer les
colonnes sensibles et limiter les accès selon le rôle de l'utilisateur.

\begin{tcolorbox}[enhanced jigsaw, arc=.35mm, bottomrule=.15mm, bottomtitle=1mm, breakable, colback=white, colbacktitle=quarto-callout-note-color!10!white, colframe=quarto-callout-note-color-frame, coltitle=black, left=2mm, leftrule=.75mm, opacityback=0, opacitybacktitle=0.6, rightrule=.15mm, title=\textcolor{quarto-callout-note-color}{\faInfo}\hspace{0.5em}{Auto-évaluation}, titlerule=0mm, toprule=.15mm, toptitle=1mm]

Avant de passer au chapitre suivant, assurez-vous de pouvoir répondre «
oui » aux questions suivantes :

\begin{itemize}
\tightlist
\item[$\square$]
  Je sais expliquer la différence entre OLTP et OLAP.
\item[$\square$]
  Je peux définir le grain d'une table de faits.
\item[$\square$]
  Je peux reconnaître une dimension et une mesure dans un exemple.
\item[$\square$]
  Je sais distinguer ETL et ELT.
\item[$\square$]
  Je comprends pourquoi les indicateurs doivent être documentés.
\end{itemize}

\end{tcolorbox}

\section{Exercices}\label{exercices-3}

\begin{tcolorbox}[enhanced jigsaw, arc=.35mm, bottomrule=.15mm, bottomtitle=1mm, breakable, colback=white, colbacktitle=quarto-callout-note-color!10!white, colframe=quarto-callout-note-color-frame, coltitle=black, left=2mm, leftrule=.75mm, opacityback=0, opacitybacktitle=0.6, rightrule=.15mm, title=\textcolor{quarto-callout-note-color}{\faInfo}\hspace{0.5em}{Consigne}, titlerule=0mm, toprule=.15mm, toptitle=1mm]

Les exercices vérifient la compréhension et peuvent être résolus sans
ordinateur. Les calculs sont volontairement courts.

\end{tcolorbox}

\begin{exercise}[OLTP ou OLAP
?]\protect\hypertarget{exr-dwh-1}{}\label{exr-dwh-1}

Indiquez si la question relève plutôt d'un système opérationnel ou
analytique :

\begin{enumerate}
\def\labelenumi{\alph{enumi}.}
\tightlist
\item
  Modifier l'adresse d'un client.
\item
  Comparer les ventes des quatre derniers trimestres par région.
\item
  Vérifier si une commande est payée.
\item
  Calculer la médiane des délais de livraison de l'année.
\end{enumerate}

\end{exercise}

\begin{tcolorbox}[enhanced jigsaw, arc=.35mm, bottomrule=.15mm, bottomtitle=1mm, breakable, colback=white, colbacktitle=quarto-callout-tip-color!10!white, colframe=quarto-callout-tip-color-frame, coltitle=black, left=2mm, leftrule=.75mm, opacityback=0, opacitybacktitle=0.6, rightrule=.15mm, title=\textcolor{quarto-callout-tip-color}{\faLightbulb}\hspace{0.5em}{Solution de l'exercice Exercice~\ref{exr-dwh-1}}, titlerule=0mm, toprule=.15mm, toptitle=1mm]

\begin{enumerate}
\def\labelenumi{\alph{enumi}.}
\tightlist
\item
  \textbf{OLTP} ; b. \textbf{OLAP} ; c.~\textbf{OLTP} ;
  d.~\textbf{OLAP}.
\end{enumerate}

\end{tcolorbox}

\begin{exercise}[Identifier le
grain]\protect\hypertarget{exr-dwh-2}{}\label{exr-dwh-2}

Une table \texttt{fact\_ventes} contient une ligne par produit présent
dans une commande. Formulez le grain avec une phrase et expliquez
pourquoi une commande de trois produits peut produire trois lignes.

\end{exercise}

\begin{tcolorbox}[enhanced jigsaw, arc=.35mm, bottomrule=.15mm, bottomtitle=1mm, breakable, colback=white, colbacktitle=quarto-callout-tip-color!10!white, colframe=quarto-callout-tip-color-frame, coltitle=black, left=2mm, leftrule=.75mm, opacityback=0, opacitybacktitle=0.6, rightrule=.15mm, title=\textcolor{quarto-callout-tip-color}{\faLightbulb}\hspace{0.5em}{Solution de l'exercice Exercice~\ref{exr-dwh-2}}, titlerule=0mm, toprule=.15mm, toptitle=1mm]

Le grain est : \textbf{une ligne représente un produit d'une commande}.
Une commande contenant trois produits possède donc trois lignes de
détail, même si elle n'a qu'un seul numéro de commande.

\end{tcolorbox}

\begin{exercise}[Faits ou
dimensions]\protect\hypertarget{exr-dwh-3}{}\label{exr-dwh-3}

Classez les éléments suivants : \texttt{montant\_net},
\texttt{categorie}, \texttt{annee}, \texttt{quantite},
\texttt{ville\_client}, \texttt{prix\_unitaire}.

\end{exercise}

\begin{tcolorbox}[enhanced jigsaw, arc=.35mm, bottomrule=.15mm, bottomtitle=1mm, breakable, colback=white, colbacktitle=quarto-callout-tip-color!10!white, colframe=quarto-callout-tip-color-frame, coltitle=black, left=2mm, leftrule=.75mm, opacityback=0, opacitybacktitle=0.6, rightrule=.15mm, title=\textcolor{quarto-callout-tip-color}{\faLightbulb}\hspace{0.5em}{Solution de l'exercice Exercice~\ref{exr-dwh-3}}, titlerule=0mm, toprule=.15mm, toptitle=1mm]

Faits ou mesures : \texttt{montant\_net}, \texttt{quantite},
\texttt{prix\_unitaire}. Dimensions ou attributs de dimensions :
\texttt{categorie}, \texttt{annee}, \texttt{ville\_client}.

\end{tcolorbox}

\begin{exercise}[Calculer une
mesure]\protect\hypertarget{exr-dwh-4}{}\label{exr-dwh-4}

Une ligne indique 3 articles à 120 MAD chacun et une remise de 40 MAD.
Calculez le montant net. Une deuxième ligne indique 2 articles à 80 MAD
sans remise. Quel est le montant total des deux lignes ?

\end{exercise}

\begin{tcolorbox}[enhanced jigsaw, arc=.35mm, bottomrule=.15mm, bottomtitle=1mm, breakable, colback=white, colbacktitle=quarto-callout-tip-color!10!white, colframe=quarto-callout-tip-color-frame, coltitle=black, left=2mm, leftrule=.75mm, opacityback=0, opacitybacktitle=0.6, rightrule=.15mm, title=\textcolor{quarto-callout-tip-color}{\faLightbulb}\hspace{0.5em}{Solution de l'exercice Exercice~\ref{exr-dwh-4}}, titlerule=0mm, toprule=.15mm, toptitle=1mm]

Première ligne : \(3\times120-40=\textbf{320}\) MAD. Deuxième ligne :
\(2\times80=\textbf{160}\) MAD. Total : \(320+160=\textbf{480}\) MAD.

\end{tcolorbox}

\begin{exercise}[Mesure additive ou non
?]\protect\hypertarget{exr-dwh-5}{}\label{exr-dwh-5}

Indiquez si l'on peut additionner directement la mesure sur les périodes
suivantes :

\begin{enumerate}
\def\labelenumi{\alph{enumi}.}
\tightlist
\item
  Le montant de ventes.
\item
  Le stock disponible à la fin de chaque jour.
\item
  Un taux de conversion en pourcentage.
\end{enumerate}

\end{exercise}

\begin{tcolorbox}[enhanced jigsaw, arc=.35mm, bottomrule=.15mm, bottomtitle=1mm, breakable, colback=white, colbacktitle=quarto-callout-tip-color!10!white, colframe=quarto-callout-tip-color-frame, coltitle=black, left=2mm, leftrule=.75mm, opacityback=0, opacitybacktitle=0.6, rightrule=.15mm, title=\textcolor{quarto-callout-tip-color}{\faLightbulb}\hspace{0.5em}{Solution de l'exercice Exercice~\ref{exr-dwh-5}}, titlerule=0mm, toprule=.15mm, toptitle=1mm]

\begin{enumerate}
\def\labelenumi{\alph{enumi}.}
\tightlist
\item
  Le montant de ventes est \textbf{additif}.
\item
  Le stock de fin de journée est \textbf{semi-additif} : on peut
  comparer ou prendre le dernier stock, mais le sommer sur les jours n'a
  généralement pas de sens.
\item
  Un taux est \textbf{non additif} : il faut recalculer le taux à partir
  des numérateurs et dénominateurs.
\end{enumerate}

\end{tcolorbox}

\begin{exercise}[Comprendre
ETL]\protect\hypertarget{exr-dwh-6}{}\label{exr-dwh-6}

Remettez dans l'ordre : charger les données préparées, extraire un
fichier CSV, convertir les dates et contrôler les prix négatifs.

\end{exercise}

\begin{tcolorbox}[enhanced jigsaw, arc=.35mm, bottomrule=.15mm, bottomtitle=1mm, breakable, colback=white, colbacktitle=quarto-callout-tip-color!10!white, colframe=quarto-callout-tip-color-frame, coltitle=black, left=2mm, leftrule=.75mm, opacityback=0, opacitybacktitle=0.6, rightrule=.15mm, title=\textcolor{quarto-callout-tip-color}{\faLightbulb}\hspace{0.5em}{Solution de l'exercice Exercice~\ref{exr-dwh-6}}, titlerule=0mm, toprule=.15mm, toptitle=1mm]

\begin{enumerate}
\def\labelenumi{\arabic{enumi}.}
\tightlist
\item
  Extraire le fichier CSV.
\item
  Convertir les dates et contrôler les prix négatifs.
\item
  Charger les données préparées.
\end{enumerate}

\end{tcolorbox}

\begin{exercise}[Moyenne
globale]\protect\hypertarget{exr-dwh-7}{}\label{exr-dwh-7}

La région A compte 10 commandes de 100 MAD en moyenne. La région B
compte 2 commandes de 300 MAD en moyenne.

\begin{enumerate}
\def\labelenumi{\alph{enumi}.}
\tightlist
\item
  Quel est le total estimé de chaque région ?
\item
  Calculez la moyenne globale.
\item
  Pourquoi la moyenne des deux moyennes, 200 MAD, est-elle incorrecte ?
\end{enumerate}

\end{exercise}

\begin{tcolorbox}[enhanced jigsaw, arc=.35mm, bottomrule=.15mm, bottomtitle=1mm, breakable, colback=white, colbacktitle=quarto-callout-tip-color!10!white, colframe=quarto-callout-tip-color-frame, coltitle=black, left=2mm, leftrule=.75mm, opacityback=0, opacitybacktitle=0.6, rightrule=.15mm, title=\textcolor{quarto-callout-tip-color}{\faLightbulb}\hspace{0.5em}{Solution de l'exercice Exercice~\ref{exr-dwh-7}}, titlerule=0mm, toprule=.15mm, toptitle=1mm]

\begin{enumerate}
\def\labelenumi{\alph{enumi}.}
\tightlist
\item
  A : \(10\times100=\textbf{1\,000}\) MAD ; B :
  \(2\times300=\textbf{600}\) MAD.
\item
  Moyenne globale : \((1000+600)/(10+2)=1600/12=\textbf{133,33}\) MAD
  environ.
\item
  Les régions n'ont pas le même nombre de commandes ; la moyenne de 200
  MAD donne le même poids aux deux groupes.
\end{enumerate}

\end{tcolorbox}

\begin{exercise}[Qualité et
sécurité]\protect\hypertarget{exr-dwh-8}{}\label{exr-dwh-8}

Proposez une règle de qualité et une mesure de protection pour chacune
des situations :

\begin{enumerate}
\def\labelenumi{\alph{enumi}.}
\tightlist
\item
  La colonne \texttt{prix} contient parfois des valeurs négatives.
\item
  Un tableau de bord doit être partagé avec des responsables régionaux.
\end{enumerate}

\end{exercise}

\begin{tcolorbox}[enhanced jigsaw, arc=.35mm, bottomrule=.15mm, bottomtitle=1mm, breakable, colback=white, colbacktitle=quarto-callout-tip-color!10!white, colframe=quarto-callout-tip-color-frame, coltitle=black, left=2mm, leftrule=.75mm, opacityback=0, opacitybacktitle=0.6, rightrule=.15mm, title=\textcolor{quarto-callout-tip-color}{\faLightbulb}\hspace{0.5em}{Solution de l'exercice Exercice~\ref{exr-dwh-8}}, titlerule=0mm, toprule=.15mm, toptitle=1mm]

\begin{enumerate}
\def\labelenumi{\alph{enumi}.}
\tightlist
\item
  Règle : \texttt{prix\ \textgreater{}=\ 0}, avec rejet ou examen des
  exceptions.
\item
  Protection : filtrer les données selon la région de l'utilisateur et
  ne partager que les colonnes nécessaires ; les noms ou coordonnées
  personnelles peuvent être masqués.
\end{enumerate}

\end{tcolorbox}

\bookmarksetup{startatroot}

\chapter{Business Intelligence avec Power BI et
Taipy}\label{business-intelligence-avec-power-bi-et-taipy}

\begin{tcolorbox}[enhanced jigsaw, arc=.35mm, bottomrule=.15mm, bottomtitle=1mm, breakable, colback=white, colbacktitle=quarto-callout-tip-color!10!white, colframe=quarto-callout-tip-color-frame, coltitle=black, left=2mm, leftrule=.75mm, opacityback=0, opacitybacktitle=0.6, rightrule=.15mm, title=\textcolor{quarto-callout-tip-color}{\faLightbulb}\hspace{0.5em}{Objectifs du chapitre}, titlerule=0mm, toprule=.15mm, toptitle=1mm]

À la fin de ce chapitre, l'étudiant sera capable de :

\begin{itemize}
\tightlist
\item
  \textbf{Définir} la Business Intelligence et distinguer un indicateur,
  un rapport et un tableau de bord ;
\item
  \textbf{Construire} mentalement le parcours d'une donnée depuis la
  source jusqu'à une décision ;
\item
  \textbf{Importer et préparer} des données dans Power BI avec Power
  Query ;
\item
  \textbf{Comprendre} un modèle en étoile et créer des mesures simples
  en DAX ;
\item
  \textbf{Choisir} entre Power BI et une application Python avec Taipy
  selon le besoin ;
\item
  \textbf{Expliquer} les principes d'un tableau de bord lisible :
  objectif, public, filtres, contexte et action.
\end{itemize}

\end{tcolorbox}

La \textbf{Business Intelligence (BI)} transforme des données préparées
en informations régulièrement consultées pour piloter une activité. Elle
ne consiste pas à accumuler des graphiques : un bon tableau de bord aide
une personne précise à répondre à des questions précises et à décider.

\section{1. Comprendre la BI}\label{comprendre-la-bi}

\subsection{1.1. Du besoin à
l'indicateur}\label{du-besoin-uxe0-lindicateur}

Un \textbf{indicateur} est une mesure définie par une formule, une
période et une population. Un \textbf{KPI} (\emph{Key Performance
Indicator}) est un indicateur considéré comme essentiel pour suivre un
objectif.

Exemple pour un site de livraison :

{\def\LTcaptype{none} % do not increment counter
\begin{longtable}[]{@{}ll@{}}
\toprule\noalign{}
Élément & Exemple \\
\midrule\noalign{}
\endhead
\bottomrule\noalign{}
\endlastfoot
Objectif & Livrer les commandes à temps \\
KPI & Taux de livraison à temps \\
Formule & livraisons à temps / livraisons terminées \\
Période & semaine ou mois \\
Cible & au moins 95 \% \\
Action possible & examiner les villes sous la cible \\
\end{longtable}
}

Un rapport peut contenir plusieurs pages et analyses détaillées. Un
tableau de bord est généralement plus synthétique : il met en évidence
les quelques indicateurs et alertes nécessaires au suivi.

\subsection{1.2. Le parcours d'une solution
BI}\label{le-parcours-dune-solution-bi}

\begin{Shaded}
\begin{Highlighting}[]
\NormalTok{Sources → préparation → modèle sémantique → mesures → visuels}
\NormalTok{                                      ↓}
\NormalTok{                             rapport / tableau de bord}
\NormalTok{                                      ↓}
\NormalTok{                              décision et action}
\end{Highlighting}
\end{Shaded}

Le modèle \textbf{sémantique} donne un sens commun aux tables, aux
relations et aux mesures. Deux personnes doivent obtenir le même chiffre
lorsqu'elles utilisent le même KPI.

\subsection{1.3. Règles d'un tableau de bord
débutant}\label{ruxe8gles-dun-tableau-de-bord-duxe9butant}

Un tableau de bord lisible :

\begin{enumerate}
\def\labelenumi{\arabic{enumi}.}
\tightlist
\item
  commence par le public et la décision attendue ;
\item
  affiche une période et une unité ;
\item
  présente quelques KPI importants avant les détails ;
\item
  utilise une couleur pour signaler une information, pas pour décorer
  chaque objet ;
\item
  permet d'explorer sans cacher la définition des indicateurs ;
\item
  indique la date de dernière actualisation.
\end{enumerate}

\section{2. Power BI : une démarche
structurée}\label{power-bi-une-duxe9marche-structuruxe9e}

Power BI est une plateforme de Business Intelligence qui permet de
connecter des sources, de les transformer, de créer un modèle de
données, de calculer des mesures et de construire des rapports
interactifs. Le parcours présenté ici commence dans Power BI Desktop ;
le partage dépend ensuite de l'environnement et des droits de
l'organisation.

\subsection{2.1. Étape 1 --- Importer les
sources}\label{uxe9tape-1-importer-les-sources}

On peut importer un fichier CSV, un classeur Excel, une base SQL ou
d'autres sources disponibles dans l'environnement. À l'import, il faut
vérifier :

\begin{itemize}
\tightlist
\item
  les noms de colonnes et leurs types ;
\item
  le séparateur et l'encodage d'un CSV ;
\item
  le type date, nombre ou texte ;
\item
  le nombre de lignes et la période couverte ;
\item
  la présence de doublons et de valeurs manquantes.
\end{itemize}

Importer n'est pas encore analyser : il faut comprendre ce qui a été
chargé.

\subsection{2.2. Étape 2 --- Transformer avec Power
Query}\label{uxe9tape-2-transformer-avec-power-query}

\textbf{Power Query} sert à préparer les données avant leur utilisation
dans le modèle. On peut supprimer des colonnes inutiles, renommer un
champ, convertir un type, filtrer des lignes, remplacer une valeur,
fusionner des tables ou ajouter une colonne calculée.

Une transformation doit rester explicable et reproductible. Power Query
conserve les étapes appliquées ; il vaut mieux créer plusieurs étapes
nommées et simples qu'une modification manuelle impossible à refaire.

Exemples de contrôles :

\begin{itemize}
\tightlist
\item
  \texttt{Montant} est numérique et positif ;
\item
  \texttt{DateCommande} est bien de type date ;
\item
  \texttt{ProduitID} et \texttt{DateID} sont présents ;
\item
  la catégorie vide est traitée explicitement ;
\item
  le nombre de lignes final est comparé au nombre attendu.
\end{itemize}

\subsection{2.3. Étape 3 --- Construire le
modèle}\label{uxe9tape-3-construire-le-moduxe8le}

Pour un modèle de ventes, on privilégie un schéma en étoile :

\begin{Shaded}
\begin{Highlighting}[]
\NormalTok{dim\_Date       dim\_Produit       dim\_Client}
\NormalTok{      \textbackslash{}             |             /}
\NormalTok{             fact\_Ventes}
\NormalTok{                   |}
\NormalTok{              dim\_Magasin}
\end{Highlighting}
\end{Shaded}

La table de faits contient les lignes de vente et les dimensions
permettent de filtrer par date, produit, client ou magasin. Les
relations doivent être vérifiées : une dimension possède généralement
une ligne par identifiant, et plusieurs lignes de faits peuvent lui être
liées.

\begin{tcolorbox}[enhanced jigsaw, arc=.35mm, bottomrule=.15mm, bottomtitle=1mm, breakable, colback=white, colbacktitle=quarto-callout-warning-color!10!white, colframe=quarto-callout-warning-color-frame, coltitle=black, left=2mm, leftrule=.75mm, opacityback=0, opacitybacktitle=0.6, rightrule=.15mm, title=\textcolor{quarto-callout-warning-color}{\faExclamationTriangle}\hspace{0.5em}{Le piège des relations ambiguës}, titlerule=0mm, toprule=.15mm, toptitle=1mm]

Si une dimension contient plusieurs lignes pour le même identifiant, ou
si deux tables sont reliées de manière inattendue, les totaux peuvent
être multipliés. Vérifiez toujours le grain, les clés et un total connu
avant de publier un rapport.

\end{tcolorbox}

\subsection{2.4. Étape 4 --- Créer des mesures en
DAX}\label{uxe9tape-4-cruxe9er-des-mesures-en-dax}

DAX est le langage de formules utilisé dans les modèles Power BI. Une
\textbf{mesure} est recalculée selon le contexte de filtre du visuel, de
la page ou de l'utilisateur.

\begin{Shaded}
\begin{Highlighting}[]
\NormalTok{ChiffreAffaires = SUM(fact\_Ventes[MontantNet])}

\NormalTok{NombreCommandes = DISTINCTCOUNT(fact\_Ventes[CommandeID])}

\NormalTok{PanierMoyen = DIVIDE([ChiffreAffaires], [NombreCommandes])}
\end{Highlighting}
\end{Shaded}

\texttt{DIVIDE} évite une erreur lorsque le dénominateur est nul. Pour
débuter, il faut distinguer une mesure d'une colonne calculée : la
mesure s'adapte aux filtres, tandis qu'une colonne calculée est créée
ligne par ligne dans la table.

Une mesure peut être comparée à une période précédente, mais la logique
de dates doit être préparée soigneusement. Dans un premier tableau de
bord, trois ou quatre mesures bien définies valent mieux que de
nombreuses formules difficiles à expliquer.

\subsection{2.5. Étape 5 --- Construire les
visuels}\label{uxe9tape-5-construire-les-visuels}

Pour une page de suivi des ventes, une composition simple peut être :

\begin{itemize}
\tightlist
\item
  une ligne de cartes KPI : chiffre d'affaires, commandes, panier moyen,
  taux de livraison à temps ;
\item
  une courbe du chiffre d'affaires par mois ;
\item
  un graphique en barres par région ou catégorie ;
\item
  un tableau détaillé filtrable ;
\item
  des segments (\emph{slicers}) pour l'année, la région et la catégorie.
\end{itemize}

Chaque visuel doit répondre à une question. Un graphique avec trop de
séries ou trop de couleurs devient difficile à lire.

\subsection{2.6. Étape 6 --- Vérifier et
partager}\label{uxe9tape-6-vuxe9rifier-et-partager}

Avant tout partage, vérifiez les totaux sur un petit exemple connu, la
date d'actualisation, les filtres, les titres et les droits d'accès. Un
rapport partagé avec des données personnelles doit respecter les règles
de l'organisation et le principe du moindre accès.

\section{3. Taipy : créer une application BI en
Python}\label{taipy-cruxe9er-une-application-bi-en-python}

Taipy permet de construire des applications Python interactives avec des
contrôles, des graphiques, des tableaux et des filtres. Il est
particulièrement intéressant quand l'équipe veut intégrer l'analyse dans
une application métier ou conserver une logique Python personnalisée.

\subsection{3.1. Quand choisir Taipy ?}\label{quand-choisir-taipy}

Taipy peut être pertinent lorsque :

\begin{itemize}
\tightlist
\item
  les traitements sont déjà écrits en Python ;
\item
  l'application doit combiner un tableau de bord avec une logique métier
  ;
\item
  l'on veut contrôler l'interface et le déploiement comme un projet
  logiciel ;
\item
  l'utilisateur doit modifier un paramètre et observer immédiatement un
  résultat.
\end{itemize}

Power BI peut être plus adapté lorsqu'une organisation veut une solution
de BI prête à l'emploi, un catalogue de rapports, des connexions gérées
et des usages majoritairement orientés vers l'exploration métier.

Ce ne sont pas des outils absolument opposés. Une équipe peut préparer
les données avec Python, les exposer dans Power BI, ou intégrer un
indicateur Power BI dans un parcours applicatif selon son architecture.

\subsection{3.2. Un exemple Taipy
minimal}\label{un-exemple-taipy-minimal}

Le code suivant illustre le principe avec des données déjà disponibles.
Il ne remplace pas la préparation et la validation des données.

\begin{Shaded}
\begin{Highlighting}[]
\ImportTok{import}\NormalTok{ pandas }\ImportTok{as}\NormalTok{ pd}
\ImportTok{from}\NormalTok{ taipy.gui }\ImportTok{import}\NormalTok{ Gui}

\NormalTok{df }\OperatorTok{=}\NormalTok{ pd.DataFrame(\{}
    \StringTok{"Mois"}\NormalTok{: [}\StringTok{"Janvier"}\NormalTok{, }\StringTok{"Février"}\NormalTok{, }\StringTok{"Mars"}\NormalTok{],}
    \StringTok{"Ventes"}\NormalTok{: [}\DecValTok{12000}\NormalTok{, }\DecValTok{14500}\NormalTok{, }\DecValTok{13200}\NormalTok{],}
\NormalTok{\})}

\NormalTok{mois }\OperatorTok{=} \StringTok{"Tous"}
\NormalTok{texte }\OperatorTok{=} \StringTok{"\#\#\# Suivi des ventes"}

\NormalTok{page }\OperatorTok{=} \StringTok{"""}
\StringTok{\# Suivi des ventes}

\StringTok{Sélection : \textless{}|}\SpecialCharTok{\{mois\}}\StringTok{|selector|lov=Tous;Janvier;Février;Mars|\textgreater{}}

\StringTok{\textless{}|}\SpecialCharTok{\{df\}}\StringTok{|table|\textgreater{}}
\StringTok{"""}

\NormalTok{Gui(page}\OperatorTok{=}\NormalTok{page).run()}
\end{Highlighting}
\end{Shaded}

Dans cet exemple, \texttt{df} contient les données affichées,
\texttt{page} décrit l'interface et le sélecteur permet de choisir une
valeur. Dans un vrai projet, on ajoute une fonction qui filtre les
données, des indicateurs calculés et des graphiques adaptés. Les noms
exacts de certains contrôles peuvent évoluer selon la version de Taipy :
il faut vérifier la documentation de la version utilisée.

\subsection{3.3. Organiser une petite
application}\label{organiser-une-petite-application}

Même une application pédagogique gagne à séparer :

\begin{Shaded}
\begin{Highlighting}[]
\NormalTok{projet/}
\NormalTok{├── data/              \# données sources ou exemple}
\NormalTok{├── preparation.py     \# lecture et nettoyage}
\NormalTok{├── indicateurs.py     \# calculs métier}
\NormalTok{├── app.py             \# interface Taipy}
\NormalTok{└── README.md          \# installation et utilisation}
\end{Highlighting}
\end{Shaded}

Cette séparation permet de tester les indicateurs indépendamment de
l'interface. L'application ne doit pas recalculer silencieusement un
indicateur avec une définition différente de celle du rapport.

\section{4. Comparaison synthétique}\label{comparaison-synthuxe9tique}

{\def\LTcaptype{none} % do not increment counter
\begin{longtable}[]{@{}
  >{\raggedright\arraybackslash}p{(\linewidth - 4\tabcolsep) * \real{0.3333}}
  >{\raggedright\arraybackslash}p{(\linewidth - 4\tabcolsep) * \real{0.3333}}
  >{\raggedright\arraybackslash}p{(\linewidth - 4\tabcolsep) * \real{0.3333}}@{}}
\toprule\noalign{}
\begin{minipage}[b]{\linewidth}\raggedright
Critère
\end{minipage} & \begin{minipage}[b]{\linewidth}\raggedright
Power BI
\end{minipage} & \begin{minipage}[b]{\linewidth}\raggedright
Taipy
\end{minipage} \\
\midrule\noalign{}
\endhead
\bottomrule\noalign{}
\endlastfoot
Profil & Plateforme BI avec interface graphique &
Bibliothèque/applications Python \\
Préparation & Power Query et connexions & Python, pandas ou autres
outils \\
Mesures & DAX et modèle sémantique & Fonctions et code Python \\
Usage principal & Rapports, exploration et partage BI & Application
interactive intégrée à un projet \\
Point d'attention & Relations, contexte des mesures, droits &
Organisation du code, déploiement, versions \\
\end{longtable}
}

Le choix doit partir du besoin, pas de la préférence pour un outil. Dans
les deux cas, les mêmes bases restent indispensables : question métier,
qualité des données, modèle clair, indicateurs documentés et
visualisation lisible.

\begin{tcolorbox}[enhanced jigsaw, arc=.35mm, bottomrule=.15mm, bottomtitle=1mm, breakable, colback=white, colbacktitle=quarto-callout-note-color!10!white, colframe=quarto-callout-note-color-frame, coltitle=black, left=2mm, leftrule=.75mm, opacityback=0, opacitybacktitle=0.6, rightrule=.15mm, title=\textcolor{quarto-callout-note-color}{\faInfo}\hspace{0.5em}{Auto-évaluation}, titlerule=0mm, toprule=.15mm, toptitle=1mm]

Avant le cas pratique, assurez-vous de pouvoir répondre « oui » aux
questions suivantes :

\begin{itemize}
\tightlist
\item[$\square$]
  Je peux distinguer un KPI, une mesure et un visuel.
\item[$\square$]
  Je sais expliquer le rôle de Power Query.
\item[$\square$]
  Je peux décrire la place de la table de faits dans un modèle Power BI.
\item[$\square$]
  Je sais pourquoi une mesure DAX n'est pas toujours une colonne
  calculée.
\item[$\square$]
  Je peux citer une situation où Taipy est pertinent.
\item[$\square$]
  Je sais vérifier un tableau de bord avant de le partager.
\end{itemize}

\end{tcolorbox}

\section{Exercices}\label{exercices-4}

\begin{tcolorbox}[enhanced jigsaw, arc=.35mm, bottomrule=.15mm, bottomtitle=1mm, breakable, colback=white, colbacktitle=quarto-callout-note-color!10!white, colframe=quarto-callout-note-color-frame, coltitle=black, left=2mm, leftrule=.75mm, opacityback=0, opacitybacktitle=0.6, rightrule=.15mm, title=\textcolor{quarto-callout-note-color}{\faInfo}\hspace{0.5em}{Consigne}, titlerule=0mm, toprule=.15mm, toptitle=1mm]

Ces exercices sont conceptuels ou demandent seulement un calcul court.
Aucun logiciel n'est nécessaire.

\end{tcolorbox}

\begin{exercise}[Indicateur ou action
?]\protect\hypertarget{exr-bi-1}{}\label{exr-bi-1}

Pour l'objectif « réduire les retards de livraison », proposez un KPI,
sa formule et une action possible si sa valeur diminue.

\end{exercise}

\begin{tcolorbox}[enhanced jigsaw, arc=.35mm, bottomrule=.15mm, bottomtitle=1mm, breakable, colback=white, colbacktitle=quarto-callout-tip-color!10!white, colframe=quarto-callout-tip-color-frame, coltitle=black, left=2mm, leftrule=.75mm, opacityback=0, opacitybacktitle=0.6, rightrule=.15mm, title=\textcolor{quarto-callout-tip-color}{\faLightbulb}\hspace{0.5em}{Solution de l'exercice Exercice~\ref{exr-bi-1}}, titlerule=0mm, toprule=.15mm, toptitle=1mm]

Exemple : KPI = \textbf{taux de livraison à temps} ; formule =
livraisons à temps / livraisons terminées × 100 ; action = analyser les
villes ou transporteurs dont le taux est le plus faible.

\end{tcolorbox}

\begin{exercise}[Lire un parcours
BI]\protect\hypertarget{exr-bi-2}{}\label{exr-bi-2}

Associez chaque étape à son rôle : préparer, modéliser, visualiser,
décider.

\begin{enumerate}
\def\labelenumi{\alph{enumi}.}
\tightlist
\item
  Relier les ventes à la dimension Date.
\item
  Convertir une colonne texte en date.
\item
  Afficher les ventes par mois.
\item
  Renforcer un transporteur après avoir constaté des retards.
\end{enumerate}

\end{exercise}

\begin{tcolorbox}[enhanced jigsaw, arc=.35mm, bottomrule=.15mm, bottomtitle=1mm, breakable, colback=white, colbacktitle=quarto-callout-tip-color!10!white, colframe=quarto-callout-tip-color-frame, coltitle=black, left=2mm, leftrule=.75mm, opacityback=0, opacitybacktitle=0.6, rightrule=.15mm, title=\textcolor{quarto-callout-tip-color}{\faLightbulb}\hspace{0.5em}{Solution de l'exercice Exercice~\ref{exr-bi-2}}, titlerule=0mm, toprule=.15mm, toptitle=1mm]

\begin{enumerate}
\def\labelenumi{\alph{enumi}.}
\tightlist
\item
  \textbf{Modéliser} ; b. \textbf{Préparer} ; c.~\textbf{Visualiser} ;
  d.~\textbf{Décider}.
\end{enumerate}

\end{tcolorbox}

\begin{exercise}[Calculer un
KPI]\protect\hypertarget{exr-bi-3}{}\label{exr-bi-3}

Sur 250 livraisons terminées, 225 ont été livrées à temps. Calculez le
taux de livraison à temps et indiquez s'il atteint une cible de 95 \%.

\end{exercise}

\begin{tcolorbox}[enhanced jigsaw, arc=.35mm, bottomrule=.15mm, bottomtitle=1mm, breakable, colback=white, colbacktitle=quarto-callout-tip-color!10!white, colframe=quarto-callout-tip-color-frame, coltitle=black, left=2mm, leftrule=.75mm, opacityback=0, opacitybacktitle=0.6, rightrule=.15mm, title=\textcolor{quarto-callout-tip-color}{\faLightbulb}\hspace{0.5em}{Solution de l'exercice Exercice~\ref{exr-bi-3}}, titlerule=0mm, toprule=.15mm, toptitle=1mm]

\(225/250\times100=90\%\). Le taux n'atteint pas la cible de 95 \%.

\end{tcolorbox}

\begin{exercise}[Power Query ou DAX
?]\protect\hypertarget{exr-bi-4}{}\label{exr-bi-4}

Indiquez l'outil le plus adapté :

\begin{enumerate}
\def\labelenumi{\alph{enumi}.}
\tightlist
\item
  Corriger le type d'une colonne \texttt{Date}.
\item
  Calculer le chiffre d'affaires selon les filtres du rapport.
\item
  Supprimer une colonne inutilisée avant le modèle.
\end{enumerate}

\end{exercise}

\begin{tcolorbox}[enhanced jigsaw, arc=.35mm, bottomrule=.15mm, bottomtitle=1mm, breakable, colback=white, colbacktitle=quarto-callout-tip-color!10!white, colframe=quarto-callout-tip-color-frame, coltitle=black, left=2mm, leftrule=.75mm, opacityback=0, opacitybacktitle=0.6, rightrule=.15mm, title=\textcolor{quarto-callout-tip-color}{\faLightbulb}\hspace{0.5em}{Solution de l'exercice Exercice~\ref{exr-bi-4}}, titlerule=0mm, toprule=.15mm, toptitle=1mm]

\begin{enumerate}
\def\labelenumi{\alph{enumi}.}
\tightlist
\item
  \textbf{Power Query} ; b. \textbf{mesure DAX} ; c.~\textbf{Power
  Query}.
\end{enumerate}

\end{tcolorbox}

\begin{exercise}[Vérifier une
relation]\protect\hypertarget{exr-bi-5}{}\label{exr-bi-5}

Une dimension \texttt{Produit} contient deux lignes pour le même
\texttt{ProduitID}. Quel risque cela crée-t-il lorsqu'on relie cette
dimension à une table de ventes ? Quelle vérification simple
proposeriez-vous ?

\end{exercise}

\begin{tcolorbox}[enhanced jigsaw, arc=.35mm, bottomrule=.15mm, bottomtitle=1mm, breakable, colback=white, colbacktitle=quarto-callout-tip-color!10!white, colframe=quarto-callout-tip-color-frame, coltitle=black, left=2mm, leftrule=.75mm, opacityback=0, opacitybacktitle=0.6, rightrule=.15mm, title=\textcolor{quarto-callout-tip-color}{\faLightbulb}\hspace{0.5em}{Solution de l'exercice Exercice~\ref{exr-bi-5}}, titlerule=0mm, toprule=.15mm, toptitle=1mm]

Les ventes peuvent être comptées plusieurs fois lors d'un regroupement
ou d'un filtre. Il faut vérifier l'unicité de \texttt{ProduitID} dans la
dimension et retrouver la cause du doublon avant de publier.

\end{tcolorbox}

\begin{exercise}[Choisir un
outil]\protect\hypertarget{exr-bi-6}{}\label{exr-bi-6}

Une équipe possède déjà plusieurs fonctions Python et veut intégrer un
indicateur interactif dans une application interne. Une autre équipe
veut surtout créer rapidement des rapports métiers partagés. Quel outil
recommander à chacune, et pourquoi ?

\end{exercise}

\begin{tcolorbox}[enhanced jigsaw, arc=.35mm, bottomrule=.15mm, bottomtitle=1mm, breakable, colback=white, colbacktitle=quarto-callout-tip-color!10!white, colframe=quarto-callout-tip-color-frame, coltitle=black, left=2mm, leftrule=.75mm, opacityback=0, opacitybacktitle=0.6, rightrule=.15mm, title=\textcolor{quarto-callout-tip-color}{\faLightbulb}\hspace{0.5em}{Solution de l'exercice Exercice~\ref{exr-bi-6}}, titlerule=0mm, toprule=.15mm, toptitle=1mm]

La première équipe peut choisir \textbf{Taipy}, pour réutiliser sa
logique Python dans une application. La seconde peut choisir
\textbf{Power BI}, orienté vers les modèles, rapports et usages BI
partagés. Le contexte de sécurité et de déploiement doit aussi être
vérifié.

\end{tcolorbox}

\begin{exercise}[Lire une
mesure]\protect\hypertarget{exr-bi-7}{}\label{exr-bi-7}

Le chiffre d'affaires est de 48 000 MAD pour 300 commandes.

\begin{enumerate}
\def\labelenumi{\alph{enumi}.}
\tightlist
\item
  Calculez le panier moyen.
\item
  Pourquoi faut-il utiliser le nombre de commandes et non le nombre de
  lignes de produits si le KPI est défini par commande ?
\end{enumerate}

\end{exercise}

\begin{tcolorbox}[enhanced jigsaw, arc=.35mm, bottomrule=.15mm, bottomtitle=1mm, breakable, colback=white, colbacktitle=quarto-callout-tip-color!10!white, colframe=quarto-callout-tip-color-frame, coltitle=black, left=2mm, leftrule=.75mm, opacityback=0, opacitybacktitle=0.6, rightrule=.15mm, title=\textcolor{quarto-callout-tip-color}{\faLightbulb}\hspace{0.5em}{Solution de l'exercice Exercice~\ref{exr-bi-7}}, titlerule=0mm, toprule=.15mm, toptitle=1mm]

\begin{enumerate}
\def\labelenumi{\alph{enumi}.}
\tightlist
\item
  \(48\,000/300=\textbf{160}\) MAD par commande.
\item
  Une commande peut contenir plusieurs lignes de produits ; compter les
  lignes gonflerait le dénominateur et changerait la définition du KPI.
\end{enumerate}

\end{tcolorbox}

\begin{exercise}[Critiquer une
page]\protect\hypertarget{exr-bi-8}{}\label{exr-bi-8}

Une page affiche 25 graphiques sans unité, sans date de mise à jour et
avec une couleur différente pour chaque barre. Citez trois améliorations
prioritaires et justifiez-les.

\end{exercise}

\begin{tcolorbox}[enhanced jigsaw, arc=.35mm, bottomrule=.15mm, bottomtitle=1mm, breakable, colback=white, colbacktitle=quarto-callout-tip-color!10!white, colframe=quarto-callout-tip-color-frame, coltitle=black, left=2mm, leftrule=.75mm, opacityback=0, opacitybacktitle=0.6, rightrule=.15mm, title=\textcolor{quarto-callout-tip-color}{\faLightbulb}\hspace{0.5em}{Solution de l'exercice Exercice~\ref{exr-bi-8}}, titlerule=0mm, toprule=.15mm, toptitle=1mm]

Exemples : réduire le nombre de visuels pour mettre en avant les KPI ;
ajouter titres, unités et période ; afficher la date d'actualisation ;
réserver les couleurs aux catégories ou alertes utiles ; ajouter des
filtres compréhensibles. Ces changements facilitent la lecture et
évitent des interprétations différentes.

\end{tcolorbox}

\bookmarksetup{startatroot}

\chapter{Cas pratique complet : analyser les ventes d'une petite chaîne
de
magasins}\label{cas-pratique-complet-analyser-les-ventes-dune-petite-chauxeene-de-magasins}

\begin{tcolorbox}[enhanced jigsaw, arc=.35mm, bottomrule=.15mm, bottomtitle=1mm, breakable, colback=white, colbacktitle=quarto-callout-tip-color!10!white, colframe=quarto-callout-tip-color-frame, coltitle=black, left=2mm, leftrule=.75mm, opacityback=0, opacitybacktitle=0.6, rightrule=.15mm, title=\textcolor{quarto-callout-tip-color}{\faLightbulb}\hspace{0.5em}{Objectifs du chapitre}, titlerule=0mm, toprule=.15mm, toptitle=1mm]

À la fin de ce cas pratique, l'étudiant sera capable de :

\begin{itemize}
\tightlist
\item
  \textbf{Cadrer} une question d'analyse à partir d'un besoin métier ;
\item
  \textbf{Repérer} les problèmes de qualité dans un petit jeu de données
  ;
\item
  \textbf{Calculer et interpréter} des indicateurs simples : chiffre
  d'affaires, quantité, panier moyen, médiane et taux ;
\item
  \textbf{Choisir} un modèle en étoile adapté à une analyse de ventes ;
\item
  \textbf{Proposer} une page Power BI et une vue Taipy cohérentes avec
  les questions ;
\item
  \textbf{Rédiger} une conclusion prudente, fondée sur les données et
  accompagnée de limites.
\end{itemize}

\end{tcolorbox}

Ce chapitre suit un projet de bout en bout. L'entreprise fictive
\textbf{Atlas Marchés} possède trois magasins à Casablanca, Rabat et
Tanger. La direction souhaite comprendre les ventes du premier trimestre
2026 et décider où concentrer ses efforts commerciaux. Les données sont
volontairement petites afin que les calculs puissent être vérifiés à la
main.

\section{1. Le besoin métier}\label{le-besoin-muxe9tier}

La direction pose quatre questions :

\begin{enumerate}
\def\labelenumi{\arabic{enumi}.}
\tightlist
\item
  Quel est le chiffre d'affaires net par mois et par magasin ?
\item
  Quels produits ou catégories contribuent le plus aux ventes ?
\item
  Le panier moyen est-il différent selon les magasins ?
\item
  Quelles informations faudrait-il surveiller dans un tableau de bord
  mensuel ?
\end{enumerate}

La question n'est pas « faire un graphique ». Le livrable attendu est
une analyse descriptive, quelques explications prudentes et une
proposition de suivi.

\subsection{1.1. Périmètre et
définitions}\label{puxe9rimuxe8tre-et-duxe9finitions}

Le périmètre est limité aux ventes \textbf{payées} enregistrées entre le
1er janvier et le 31 mars 2026. Une ligne du fichier correspond à une
ligne de produit dans une commande. Le montant net est défini ainsi :

\[
\text{montant net}
=\text{quantité}\times\text{prix unitaire}-\text{remise}.
\]

Le \textbf{panier moyen} est :

\[
\text{panier moyen}
=\frac{\text{chiffre d'affaires net}}{\text{nombre de commandes distinctes}}.
\]

Ces définitions sont annoncées avant les calculs ; elles évitent de
changer de formule au milieu de l'analyse.

\section{2. Le jeu de données}\label{le-jeu-de-donnuxe9es}

Voici un extrait pédagogique. Dans un vrai projet, le fichier complet
pourrait contenir beaucoup plus de lignes.

{\def\LTcaptype{none} % do not increment counter
\begin{longtable}[]{@{}
  >{\raggedleft\arraybackslash}p{(\linewidth - 18\tabcolsep) * \real{0.1000}}
  >{\raggedright\arraybackslash}p{(\linewidth - 18\tabcolsep) * \real{0.1000}}
  >{\raggedright\arraybackslash}p{(\linewidth - 18\tabcolsep) * \real{0.1000}}
  >{\raggedright\arraybackslash}p{(\linewidth - 18\tabcolsep) * \real{0.1000}}
  >{\raggedright\arraybackslash}p{(\linewidth - 18\tabcolsep) * \real{0.1000}}
  >{\raggedright\arraybackslash}p{(\linewidth - 18\tabcolsep) * \real{0.1000}}
  >{\raggedleft\arraybackslash}p{(\linewidth - 18\tabcolsep) * \real{0.1000}}
  >{\raggedleft\arraybackslash}p{(\linewidth - 18\tabcolsep) * \real{0.1000}}
  >{\raggedleft\arraybackslash}p{(\linewidth - 18\tabcolsep) * \real{0.1000}}
  >{\raggedright\arraybackslash}p{(\linewidth - 18\tabcolsep) * \real{0.1000}}@{}}
\toprule\noalign{}
\begin{minipage}[b]{\linewidth}\raggedleft
Ligne
\end{minipage} & \begin{minipage}[b]{\linewidth}\raggedright
Commande
\end{minipage} & \begin{minipage}[b]{\linewidth}\raggedright
Date
\end{minipage} & \begin{minipage}[b]{\linewidth}\raggedright
Magasin
\end{minipage} & \begin{minipage}[b]{\linewidth}\raggedright
Catégorie
\end{minipage} & \begin{minipage}[b]{\linewidth}\raggedright
Produit
\end{minipage} & \begin{minipage}[b]{\linewidth}\raggedleft
Qté
\end{minipage} & \begin{minipage}[b]{\linewidth}\raggedleft
Prix unitaire
\end{minipage} & \begin{minipage}[b]{\linewidth}\raggedleft
Remise
\end{minipage} & \begin{minipage}[b]{\linewidth}\raggedright
Statut
\end{minipage} \\
\midrule\noalign{}
\endhead
\bottomrule\noalign{}
\endlastfoot
1 & C001 & 05/01/2026 & Casablanca & Épicerie & Huile & 2 & 75 & 0 &
Payée \\
2 & C001 & 05/01/2026 & Casablanca & Épicerie & Riz & 3 & 20 & 0 &
Payée \\
3 & C002 & 08/01/2026 & Rabat & Boissons & Jus & 4 & 18 & 5 & Payée \\
4 & C003 & 15/01/2026 & Tanger & Hygiène & Savon & 5 & 12 & 0 & Payée \\
5 & C004 & 03/02/2026 & Casablanca & Épicerie & Huile & 1 & 75 & 0 &
Payée \\
6 & C005 & 10/02/2026 & Rabat & Épicerie & Riz & 5 & 20 & 10 & Payée \\
7 & C006 & 18/02/2026 & Tanger & Boissons & Jus & 2 & 18 & 0 &
Annulée \\
8 & C007 & 22/02/2026 & Casablanca & Hygiène & Savon & 6 & 12 & 0 &
Payée \\
9 & C008 & 06/03/2026 & Rabat & Boissons & Jus & 6 & 18 & 0 & Payée \\
10 & C009 & 12/03/2026 & Tanger & Épicerie & Riz & 2 & 20 & 0 & Payée \\
11 & C010 & 20/03/2026 & Casablanca & Boissons & Jus & 3 & 18 & 5 &
Payée \\
12 & C010 & 20/03/2026 & Casablanca & Hygiène & Savon & 2 & 12 & 0 &
Payée \\
\end{longtable}
}

\subsection{2.1. Première lecture}\label{premiuxe8re-lecture}

Le grain est \textbf{une ligne de produit dans une commande}. La
commande C001 apparaît deux fois car elle contient deux produits. Pour
compter les commandes, il faut compter les identifiants distincts, et
non les lignes.

La ligne C006 est annulée. Si le chiffre d'affaires concerne les ventes
payées, elle doit être exclue, mais cette décision doit être documentée.
Elle ne doit pas être supprimée du fichier original : on conserve le
statut pour pouvoir analyser les annulations.

\section{3. Préparer les données}\label{pruxe9parer-les-donnuxe9es}

\subsection{3.1. Dictionnaire minimal}\label{dictionnaire-minimal}

{\def\LTcaptype{none} % do not increment counter
\begin{longtable}[]{@{}lll@{}}
\toprule\noalign{}
Colonne & Type attendu & Règle \\
\midrule\noalign{}
\endhead
\bottomrule\noalign{}
\endlastfoot
\texttt{Commande} & Texte & Identifiant non vide \\
\texttt{Date} & Date & Date comprise dans la période \\
\texttt{Magasin} & Catégorie & Casablanca, Rabat ou Tanger \\
\texttt{Qté} & Entier & Strictement positive pour une vente \\
\texttt{Prix\ unitaire} & Décimal & Positif, en MAD \\
\texttt{Remise} & Décimal & Supérieure ou égale à 0, en MAD \\
\texttt{Statut} & Catégorie & Payée ou Annulée \\
\end{longtable}
}

Il faut également vérifier que la remise ne dépasse pas le montant brut
de la ligne, que les dates ne sont pas inversées et que les identifiants
ne comportent pas d'espaces cachés.

\subsection{3.2. Transformation}\label{transformation}

Pour chaque ligne, on calcule :

\begin{Shaded}
\begin{Highlighting}[]
\NormalTok{montant\_brut = quantité × prix\_unitaire}
\NormalTok{montant\_net = montant\_brut − remise}
\NormalTok{mois = mois(date)}
\end{Highlighting}
\end{Shaded}

Puis on filtre \texttt{Statut\ =\ Payée} pour les indicateurs de chiffre
d'affaires. On conserve une copie de l'ensemble initial et on compte les
lignes exclues.

Pour l'extrait, 11 lignes sur 12 sont payées. La ligne annulée
représente 1 ligne sur 12, soit environ 8,3 \% des lignes ; ce
pourcentage ne doit pas être interprété comme un taux de commandes
annulées, car une commande peut comporter plusieurs lignes.

\section{4. Calculer les indicateurs}\label{calculer-les-indicateurs}

\subsection{4.1. Montants ligne par
ligne}\label{montants-ligne-par-ligne}

Sur les lignes payées :

{\def\LTcaptype{none} % do not increment counter
\begin{longtable}[]{@{}lllr@{}}
\toprule\noalign{}
Commande & Produit & Calcul & Montant net \\
\midrule\noalign{}
\endhead
\bottomrule\noalign{}
\endlastfoot
C001 & Huile & \(2\times75-0\) & 150 \\
C001 & Riz & \(3\times20-0\) & 60 \\
C002 & Jus & \(4\times18-5\) & 67 \\
C003 & Savon & \(5\times12-0\) & 60 \\
C004 & Huile & \(1\times75-0\) & 75 \\
C005 & Riz & \(5\times20-10\) & 90 \\
C007 & Savon & \(6\times12-0\) & 72 \\
C008 & Jus & \(6\times18-0\) & 108 \\
C009 & Riz & \(2\times20-0\) & 40 \\
C010 & Jus & \(3\times18-5\) & 49 \\
C010 & Savon & \(2\times12-0\) & 24 \\
\end{longtable}
}

Le chiffre d'affaires net de l'extrait est donc :

\[
150+60+67+60+75+90+72+108+40+49+24
=\textbf{795 MAD}.
\]

\subsection{4.2. Par mois}\label{par-mois}

{\def\LTcaptype{none} % do not increment counter
\begin{longtable}[]{@{}lrrr@{}}
\toprule\noalign{}
Mois & Chiffre d'affaires net & Commandes distinctes & Panier moyen \\
\midrule\noalign{}
\endhead
\bottomrule\noalign{}
\endlastfoot
Janvier & 337 & 3 & 112,33 \\
Février & 237 & 3 & 79,00 \\
Mars & 221 & 3 & 73,67 \\
\textbf{Total} & \textbf{795} & \textbf{9} & \textbf{88,33} \\
\end{longtable}
}

Le panier moyen total est \(795/9=88,33\) MAD. Il ne faut pas faire la
moyenne simple des trois paniers mensuels sans vérifier les effectifs ;
ici ils sont égaux, mais ce n'est pas toujours le cas.

\subsection{4.3. Par magasin et par
catégorie}\label{par-magasin-et-par-catuxe9gorie}

En regroupant les lignes payées :

{\def\LTcaptype{none} % do not increment counter
\begin{longtable}[]{@{}
  >{\raggedright\arraybackslash}p{(\linewidth - 6\tabcolsep) * \real{0.2500}}
  >{\raggedleft\arraybackslash}p{(\linewidth - 6\tabcolsep) * \real{0.2500}}
  >{\raggedleft\arraybackslash}p{(\linewidth - 6\tabcolsep) * \real{0.2500}}
  >{\raggedleft\arraybackslash}p{(\linewidth - 6\tabcolsep) * \real{0.2500}}@{}}
\toprule\noalign{}
\begin{minipage}[b]{\linewidth}\raggedright
Magasin
\end{minipage} & \begin{minipage}[b]{\linewidth}\raggedleft
Chiffre d'affaires net
\end{minipage} & \begin{minipage}[b]{\linewidth}\raggedleft
Commandes distinctes
\end{minipage} & \begin{minipage}[b]{\linewidth}\raggedleft
Panier moyen
\end{minipage} \\
\midrule\noalign{}
\endhead
\bottomrule\noalign{}
\endlastfoot
Casablanca & 435 & 4 & 108,75 \\
Rabat & 157 & 3 & 52,33 \\
Tanger & 100 & 2 & 50,00 \\
\end{longtable}
}

{\def\LTcaptype{none} % do not increment counter
\begin{longtable}[]{@{}lr@{}}
\toprule\noalign{}
Catégorie & Chiffre d'affaires net \\
\midrule\noalign{}
\endhead
\bottomrule\noalign{}
\endlastfoot
Épicerie & 415 \\
Boissons & 224 \\
Hygiène & 156 \\
\end{longtable}
}

Casablanca contribue le plus au chiffre d'affaires et possède le panier
moyen le plus élevé dans cet extrait. Cela ne suffit pas à conclure que
le magasin est le plus performant : il faudrait connaître le nombre de
visiteurs, les coûts, la taille des magasins et une période plus longue.

\subsection{4.4. Moyenne et médiane des montants de
lignes}\label{moyenne-et-muxe9diane-des-montants-de-lignes}

Les montants des lignes payées sont :

\[60,\ 60,\ 67,\ 72,\ 75,\ 90,\ 108,\ 24,\ 40,\ 49,\ 150\]

Triés :

\[24,\ 40,\ 49,\ 60,\ 60,\ 67,\ 72,\ 75,\ 90,\ 108,\ 150.\]

La médiane est la sixième valeur, soit \textbf{67 MAD}. La moyenne est
\(795/11 \approx \textbf{72,27}\) MAD. La moyenne un peu supérieure à la
médiane suggère que quelques lignes élevées, notamment 150 MAD, tirent
la moyenne vers le haut.

\section{5. Construire le modèle
analytique}\label{construire-le-moduxe8le-analytique}

À partir de ce fichier, un modèle en étoile simple pourrait être :

\begin{Shaded}
\begin{Highlighting}[]
\NormalTok{dim\_date {-}{-}{-}{-}{-}{-}{-}{-}{-}{-}{-}\textbackslash{}}
\NormalTok{dim\_produit {-}{-}{-}{-}{-}{-}{-}{-}{-} fact\_ventes}
\NormalTok{dim\_magasin {-}{-}{-}{-}{-}{-}{-}{-}/}
\end{Highlighting}
\end{Shaded}

\texttt{fact\_ventes} contient \texttt{Commande}, \texttt{DateID},
\texttt{ProduitID}, \texttt{MagasinID}, \texttt{Quantite},
\texttt{PrixUnitaire}, \texttt{Remise}, \texttt{MontantNet} et
\texttt{Statut}. Les dimensions décrivent les dates, produits et
magasins.

Les mesures attendues sont :

\begin{Shaded}
\begin{Highlighting}[]
\NormalTok{ChiffreAffaires = somme de MontantNet pour les lignes payées}
\NormalTok{Commandes = nombre d\textquotesingle{}identifiants Commande distincts pour les lignes payées}
\NormalTok{PanierMoyen = ChiffreAffaires / Commandes}
\NormalTok{Quantités = somme de Quantite pour les lignes payées}
\end{Highlighting}
\end{Shaded}

Le tableau de bord ne doit pas calculer le panier moyen comme la moyenne
des montants de lignes : une commande peut contenir plusieurs lignes. Le
grain explique cette différence.

\section{6. Concevoir la page Power
BI}\label{concevoir-la-page-power-bi}

Une page « Vue générale » peut contenir :

\begin{enumerate}
\def\labelenumi{\arabic{enumi}.}
\tightlist
\item
  quatre cartes : chiffre d'affaires net, commandes, panier moyen et
  quantités vendues ;
\item
  une courbe du chiffre d'affaires par mois ;
\item
  des barres du chiffre d'affaires par magasin ;
\item
  des barres du chiffre d'affaires par catégorie ;
\item
  des segments pour le mois, le magasin et la catégorie ;
\item
  une note indiquant « données payées, période janvier--mars 2026,
  dernière actualisation ».
\end{enumerate}

Une deuxième page « Diagnostic » peut détailler les produits, les
annulations et les différences entre magasins. La première page répond à
« où en sommes-nous ? », la seconde aide à explorer « où regarder
ensuite ? ».

Avant le partage, on compare les cartes aux totaux calculés manuellement
: 795 MAD, 9 commandes distinctes et 88,33 MAD de panier moyen.

\section{7. Construire une vue Taipy}\label{construire-une-vue-taipy}

La même analyse peut être intégrée à une application Python. La logique
doit être séparée de l'affichage :

\begin{Shaded}
\begin{Highlighting}[]
\KeywordTok{def}\NormalTok{ indicateurs(df):}
\NormalTok{    payees }\OperatorTok{=}\NormalTok{ df[df[}\StringTok{"Statut"}\NormalTok{] }\OperatorTok{==} \StringTok{"Payée"}\NormalTok{]}
\NormalTok{    chiffre\_affaires }\OperatorTok{=}\NormalTok{ payees[}\StringTok{"MontantNet"}\NormalTok{].}\BuiltInTok{sum}\NormalTok{()}
\NormalTok{    commandes }\OperatorTok{=}\NormalTok{ payees[}\StringTok{"Commande"}\NormalTok{].nunique()}
\NormalTok{    panier\_moyen }\OperatorTok{=}\NormalTok{ chiffre\_affaires }\OperatorTok{/}\NormalTok{ commandes }\ControlFlowTok{if}\NormalTok{ commandes }\ControlFlowTok{else} \DecValTok{0}
    \ControlFlowTok{return}\NormalTok{ \{}
        \StringTok{"chiffre\_affaires"}\NormalTok{: chiffre\_affaires,}
        \StringTok{"commandes"}\NormalTok{: commandes,}
        \StringTok{"panier\_moyen"}\NormalTok{: panier\_moyen,}
\NormalTok{    \}}
\end{Highlighting}
\end{Shaded}

Une interface Taipy peut ensuite proposer un filtre \texttt{Magasin},
afficher les trois indicateurs et mettre à jour une table mensuelle. Le
résultat doit rester cohérent avec le dictionnaire métier et les calculs
Power BI.

\section{8. Interpréter et
communiquer}\label{interpruxe9ter-et-communiquer}

Une conclusion raisonnable pour l'extrait serait :

\begin{quote}
Sur les ventes payées de l'extrait, le chiffre d'affaires net atteint
795 MAD pour 9 commandes, avec un panier moyen de 88,33 MAD. Casablanca
représente 435 MAD et le panier moyen y est supérieur à celui de Rabat
et Tanger. Les lignes de vente ont une médiane de 67 MAD et une moyenne
de 72,27 MAD, ce qui signale une légère influence de quelques montants
élevés. Ces résultats sont exploratoires : l'échantillon est très petit,
ne couvre que trois mois et ne permet pas d'expliquer les causes ni de
mesurer la rentabilité.
\end{quote}

La conclusion sépare les \textbf{faits calculés}, les
\textbf{interprétations prudentes} et les \textbf{limites}. Elle ne
transforme pas une différence observée en causalité.

\begin{tcolorbox}[enhanced jigsaw, arc=.35mm, bottomrule=.15mm, bottomtitle=1mm, breakable, colback=white, colbacktitle=quarto-callout-note-color!10!white, colframe=quarto-callout-note-color-frame, coltitle=black, left=2mm, leftrule=.75mm, opacityback=0, opacitybacktitle=0.6, rightrule=.15mm, title=\textcolor{quarto-callout-note-color}{\faInfo}\hspace{0.5em}{Auto-évaluation}, titlerule=0mm, toprule=.15mm, toptitle=1mm]

Avant de considérer le cas comme terminé, assurez-vous de pouvoir
répondre « oui » aux questions suivantes :

\begin{itemize}
\tightlist
\item[$\square$]
  Je peux formuler le grain du fichier.
\item[$\square$]
  Je sais pourquoi la ligne annulée est filtrée pour le chiffre
  d'affaires.
\item[$\square$]
  Je peux recalculer le total de 795 MAD.
\item[$\square$]
  Je distingue lignes de produits et commandes distinctes.
\item[$\square$]
  Je peux expliquer la différence entre moyenne et médiane dans ce cas.
\item[$\square$]
  Je peux proposer une page de tableau de bord sans surcharger
  l'utilisateur.
\end{itemize}

\end{tcolorbox}

\section{Exercices}\label{exercices-5}

\begin{tcolorbox}[enhanced jigsaw, arc=.35mm, bottomrule=.15mm, bottomtitle=1mm, breakable, colback=white, colbacktitle=quarto-callout-note-color!10!white, colframe=quarto-callout-note-color-frame, coltitle=black, left=2mm, leftrule=.75mm, opacityback=0, opacitybacktitle=0.6, rightrule=.15mm, title=\textcolor{quarto-callout-note-color}{\faInfo}\hspace{0.5em}{Consigne}, titlerule=0mm, toprule=.15mm, toptitle=1mm]

Répondez sans ordinateur. Les données de l'énoncé suffisent ; les
calculs sont courts et peuvent être réalisés avec une calculatrice.

\end{tcolorbox}

\begin{exercise}[Définir le
grain]\protect\hypertarget{exr-cas-1}{}\label{exr-cas-1}

Quelle phrase décrit le grain du fichier ? Pourquoi la commande C001
occupe-t-elle deux lignes ?

\end{exercise}

\begin{tcolorbox}[enhanced jigsaw, arc=.35mm, bottomrule=.15mm, bottomtitle=1mm, breakable, colback=white, colbacktitle=quarto-callout-tip-color!10!white, colframe=quarto-callout-tip-color-frame, coltitle=black, left=2mm, leftrule=.75mm, opacityback=0, opacitybacktitle=0.6, rightrule=.15mm, title=\textcolor{quarto-callout-tip-color}{\faLightbulb}\hspace{0.5em}{Solution de l'exercice Exercice~\ref{exr-cas-1}}, titlerule=0mm, toprule=.15mm, toptitle=1mm]

Une ligne représente \textbf{un produit dans une commande}. C001
contient de l'huile et du riz : elle possède donc deux lignes de
produits.

\end{tcolorbox}

\begin{exercise}[Exclure une
annulation]\protect\hypertarget{exr-cas-2}{}\label{exr-cas-2}

Pourquoi la ligne C006 n'est-elle pas incluse dans le chiffre d'affaires
payé ? Peut-on la supprimer du fichier original ?

\end{exercise}

\begin{tcolorbox}[enhanced jigsaw, arc=.35mm, bottomrule=.15mm, bottomtitle=1mm, breakable, colback=white, colbacktitle=quarto-callout-tip-color!10!white, colframe=quarto-callout-tip-color-frame, coltitle=black, left=2mm, leftrule=.75mm, opacityback=0, opacitybacktitle=0.6, rightrule=.15mm, title=\textcolor{quarto-callout-tip-color}{\faLightbulb}\hspace{0.5em}{Solution de l'exercice Exercice~\ref{exr-cas-2}}, titlerule=0mm, toprule=.15mm, toptitle=1mm]

Elle est marquée \textbf{Annulée}, donc elle ne représente pas une vente
payée selon la définition du périmètre. Il faut la conserver dans la
source afin de pouvoir suivre les annulations et documenter le filtre.

\end{tcolorbox}

\begin{exercise}[Calculer une
ligne]\protect\hypertarget{exr-cas-3}{}\label{exr-cas-3}

Calculez le montant net de 4 jus à 18 MAD avec une remise de 5 MAD, puis
celui de 5 savons à 12 MAD sans remise.

\end{exercise}

\begin{tcolorbox}[enhanced jigsaw, arc=.35mm, bottomrule=.15mm, bottomtitle=1mm, breakable, colback=white, colbacktitle=quarto-callout-tip-color!10!white, colframe=quarto-callout-tip-color-frame, coltitle=black, left=2mm, leftrule=.75mm, opacityback=0, opacitybacktitle=0.6, rightrule=.15mm, title=\textcolor{quarto-callout-tip-color}{\faLightbulb}\hspace{0.5em}{Solution de l'exercice Exercice~\ref{exr-cas-3}}, titlerule=0mm, toprule=.15mm, toptitle=1mm]

Jus : \(4\times18-5=\textbf{67}\) MAD. Savons :
\(5\times12=\textbf{60}\) MAD.

\end{tcolorbox}

\begin{exercise}[Vérifier le
total]\protect\hypertarget{exr-cas-4}{}\label{exr-cas-4}

Le chiffre d'affaires net est de 795 MAD et le magasin de Casablanca
réalise 435 MAD.

\begin{enumerate}
\def\labelenumi{\alph{enumi}.}
\tightlist
\item
  Quelle part du chiffre d'affaires total représente Casablanca ?
\item
  Quel montant reste pour Rabat et Tanger ensemble ?
\end{enumerate}

\end{exercise}

\begin{tcolorbox}[enhanced jigsaw, arc=.35mm, bottomrule=.15mm, bottomtitle=1mm, breakable, colback=white, colbacktitle=quarto-callout-tip-color!10!white, colframe=quarto-callout-tip-color-frame, coltitle=black, left=2mm, leftrule=.75mm, opacityback=0, opacitybacktitle=0.6, rightrule=.15mm, title=\textcolor{quarto-callout-tip-color}{\faLightbulb}\hspace{0.5em}{Solution de l'exercice Exercice~\ref{exr-cas-4}}, titlerule=0mm, toprule=.15mm, toptitle=1mm]

\begin{enumerate}
\def\labelenumi{\alph{enumi}.}
\tightlist
\item
  \(435/795\times100 \approx 54,7\%\).
\item
  \(795-435=\textbf{360}\) MAD.
\end{enumerate}

\end{tcolorbox}

\begin{exercise}[Commandes et
lignes]\protect\hypertarget{exr-cas-5}{}\label{exr-cas-5}

Il y a 11 lignes payées mais seulement 9 commandes distinctes.

\begin{enumerate}
\def\labelenumi{\alph{enumi}.}
\tightlist
\item
  Pourquoi ces deux nombres sont-ils différents ?
\item
  Quel nombre utiliser pour calculer le panier moyen ?
\end{enumerate}

\end{exercise}

\begin{tcolorbox}[enhanced jigsaw, arc=.35mm, bottomrule=.15mm, bottomtitle=1mm, breakable, colback=white, colbacktitle=quarto-callout-tip-color!10!white, colframe=quarto-callout-tip-color-frame, coltitle=black, left=2mm, leftrule=.75mm, opacityback=0, opacitybacktitle=0.6, rightrule=.15mm, title=\textcolor{quarto-callout-tip-color}{\faLightbulb}\hspace{0.5em}{Solution de l'exercice Exercice~\ref{exr-cas-5}}, titlerule=0mm, toprule=.15mm, toptitle=1mm]

\begin{enumerate}
\def\labelenumi{\alph{enumi}.}
\tightlist
\item
  Certaines commandes contiennent plusieurs produits et donc plusieurs
  lignes ; C001 et C010 en sont des exemples.
\item
  Il faut utiliser \textbf{9 commandes distinctes}, car le panier moyen
  est défini par commande.
\end{enumerate}

\end{tcolorbox}

\begin{exercise}[Comparer moyenne et
médiane]\protect\hypertarget{exr-cas-6}{}\label{exr-cas-6}

La moyenne des montants de lignes est 72,27 MAD et la médiane est 67
MAD. Que suggère cette différence ? Quelle valeur de la liste contribue
particulièrement à tirer la moyenne vers le haut ?

\end{exercise}

\begin{tcolorbox}[enhanced jigsaw, arc=.35mm, bottomrule=.15mm, bottomtitle=1mm, breakable, colback=white, colbacktitle=quarto-callout-tip-color!10!white, colframe=quarto-callout-tip-color-frame, coltitle=black, left=2mm, leftrule=.75mm, opacityback=0, opacitybacktitle=0.6, rightrule=.15mm, title=\textcolor{quarto-callout-tip-color}{\faLightbulb}\hspace{0.5em}{Solution de l'exercice Exercice~\ref{exr-cas-6}}, titlerule=0mm, toprule=.15mm, toptitle=1mm]

La moyenne supérieure à la médiane suggère une légère asymétrie à droite
: quelques montants élevés influencent la moyenne. La ligne à 150 MAD y
contribue particulièrement.

\end{tcolorbox}

\begin{exercise}[Choisir les indicateurs du tableau de
bord]\protect\hypertarget{exr-cas-7}{}\label{exr-cas-7}

Parmi les éléments suivants, lesquels doivent apparaître en priorité sur
la page générale : chiffre d'affaires net, identifiant de chaque client,
commandes distinctes, panier moyen, une liste de 25 couleurs, date
d'actualisation ? Justifiez les choix.

\end{exercise}

\begin{tcolorbox}[enhanced jigsaw, arc=.35mm, bottomrule=.15mm, bottomtitle=1mm, breakable, colback=white, colbacktitle=quarto-callout-tip-color!10!white, colframe=quarto-callout-tip-color-frame, coltitle=black, left=2mm, leftrule=.75mm, opacityback=0, opacitybacktitle=0.6, rightrule=.15mm, title=\textcolor{quarto-callout-tip-color}{\faLightbulb}\hspace{0.5em}{Solution de l'exercice Exercice~\ref{exr-cas-7}}, titlerule=0mm, toprule=.15mm, toptitle=1mm]

Les priorités sont le \textbf{chiffre d'affaires net}, les
\textbf{commandes distinctes}, le \textbf{panier moyen} et la
\textbf{date d'actualisation}. Un identifiant individuel n'est pas
nécessaire pour la vue générale et peut poser un problème de
confidentialité ; 25 couleurs nuisent à la lecture.

\end{tcolorbox}

\begin{exercise}[Limites et
décisions]\protect\hypertarget{exr-cas-8}{}\label{exr-cas-8}

La direction demande : « Casablanca est-elle le meilleur magasin ? »

\begin{enumerate}
\def\labelenumi{\alph{enumi}.}
\tightlist
\item
  Donnez une conclusion que les données permettent de soutenir.
\item
  Citez deux informations manquantes avant de conclure sur la
  performance globale.
\end{enumerate}

\end{exercise}

\begin{tcolorbox}[enhanced jigsaw, arc=.35mm, bottomrule=.15mm, bottomtitle=1mm, breakable, colback=white, colbacktitle=quarto-callout-tip-color!10!white, colframe=quarto-callout-tip-color-frame, coltitle=black, left=2mm, leftrule=.75mm, opacityback=0, opacitybacktitle=0.6, rightrule=.15mm, title=\textcolor{quarto-callout-tip-color}{\faLightbulb}\hspace{0.5em}{Solution de l'exercice Exercice~\ref{exr-cas-8}}, titlerule=0mm, toprule=.15mm, toptitle=1mm]

\begin{enumerate}
\def\labelenumi{\alph{enumi}.}
\tightlist
\item
  On peut dire que Casablanca a le chiffre d'affaires net et le panier
  moyen les plus élevés dans \textbf{cet extrait de ventes payées}.
\item
  Il manque par exemple les coûts et la marge, le nombre de visiteurs,
  la taille des magasins, le stock disponible, une période plus longue
  et les objectifs de chaque magasin. Le chiffre d'affaires seul ne
  suffit pas à définir la performance globale.
\end{enumerate}

\end{tcolorbox}

\bookmarksetup{startatroot}

\chapter*{References}\label{references}
\addcontentsline{toc}{chapter}{References}

\markboth{References}{References}

\protect\phantomsection\label{refs}
\begin{CSLReferences}{1}{1}
\bibitem[\citeproctext]{ref-knaflic2015}
Knaflic, Cole. 2015. \emph{Storytelling with Data}. Wiley.

\bibitem[\citeproctext]{ref-mckinney2022}
McKinney, Wes. 2022. \emph{Python for Data Analysis}. O'Reilly Media.

\end{CSLReferences}


\backmatter


\end{document}
